% Stéréochimie : isomérie et représentation des molécules — Chimie Tle TI — Cours signé OnBuch+
% Programme officiel MINESEC, Terminale TI. Compiler avec : tectonic L05-stereochimie-isomerie.tex
\newcommand{\DOCMATIERE}{Chimie}
\newcommand{\DOCNIVEAU}{Tle TI}
\newcommand{\DOCMODULE}{Module 1 · Chimie organique}
\newcommand{\DOCLECON}{05}
\newcommand{\DOCTITRE}{Stéréochimie : isomérie et représentation des molécules}
% OnBuch+ — préambule « cours signés » ENRICHI (compilé avec tectonic / XeLaTeX)
% Système visuel « Pop » d'OnBuch+ : fond crème (jamais blanc), bordures encre
% épaisses, ombres « dures » décalées, titres Archivo Black en capitales.
% Version enrichie pour les cours de chimie : chimie (mhchem, chemfig),
% graphes (pgfplots), boîtes pédagogiques supplémentaires (définition,
% propriété, méthode, attention, le savais-tu, exercice résolu, exercice,
% TP, bilan), page de garde refaite et signature OnBuch+ ancrée en bas.
%
% Macros attendues dans main.tex AVANT \input{preamble} :
%   \DOCMATIERE  \DOCNIVEAU  \DOCMODULE  \DOCLECON (numéro)  \DOCTITRE
\documentclass[11pt]{article}

\usepackage{fontspec}
\usepackage[french]{babel}
\usepackage[a4paper,margin=2cm,top=3.4cm,bottom=2.8cm,headheight=1.4cm,footskip=1.3cm]{geometry}
\usepackage{amsmath,amssymb,amsfonts}
\usepackage[table]{xcolor} % [table] : fournit aussi \rowcolors (alternance de lignes)
\usepackage{pagecolor}
\usepackage{tikz}
\usetikzlibrary{arrows.meta,positioning,calc,shapes.geometric,shapes.misc,shapes.arrows,decorations.pathmorphing,decorations.markings,patterns,shadows,mindmap,trees,fit,backgrounds,matrix,intersections}
\usepackage{pgfplots}
\pgfplotsset{compat=1.18}
\usepackage[version=4]{mhchem}
\usepackage{chemfig}
\usepackage{enumitem}
\usepackage{fancyhdr}
% [most] charge toutes les bibliothèques tcolorbox (listings, minted, raster,
% documentation...) et gonfle inutilement la mémoire TeX sur un document long
% (~300 boîtes) : on ne charge que ce qui est réellement utilisé.
\usepackage[skins,breakable]{tcolorbox}
\usepackage{graphicx}
\usepackage{colortbl}
\usepackage{booktabs}
\usepackage{tabularx}
\usepackage{array}
\usepackage{multicol}
\usepackage{siunitx}
% parse-numbers=false : accepte aussi \SI{2,5 \times 10^{-3}}{...}
\sisetup{locale=FR,output-decimal-marker={,},group-minimum-digits=4,parse-numbers=false}
\DeclareSIUnit\ohm{\ensuremath{\symup{\Omega}}} % U+2126 absent de la police
\usepackage{qrcode}
% \degree : commande fréquemment utilisée par les modèles pour un angle ou une
% température en dehors de \SI{}{} (non fournie par les paquets chargés).
\providecommand{\degree}{^{\circ}}
\usepackage{lastpage}
\usepackage{hyperref}
\hypersetup{hidelinks}

% ============================================================
% Palette « Pop » OnBuch+ — mêmes hex que la classe P (plus_ui.dart)
% ============================================================
\definecolor{poporange}{HTML}{F59321}
\definecolor{popdark}{HTML}{E07A0C}
\definecolor{popcream}{HTML}{FFF6E8}
\definecolor{popink}{HTML}{1C1714}
\definecolor{poppurple}{HTML}{7A5AE0}
\definecolor{popgreen}{HTML}{2E9E6A}
\definecolor{poppink}{HTML}{E0457B}
\definecolor{popblue}{HTML}{2F7DE1}
\definecolor{popgold}{HTML}{C98A12}
\definecolor{popmuted}{HTML}{8A7F76}
\definecolor{popline}{HTML}{EADFD2}
\definecolor{popgray}{HTML}{8A7F76}
\colorlet{popgrey}{popgray}
% Alias tolérants (le modèle nomme parfois les couleurs différemment)
\colorlet{popwhite}{white}
\colorlet{popblack}{popink}
\colorlet{popred}{poppink}
\colorlet{popyellow}{popgold}
% Teintes claires pour fonds de tableaux / zones
\colorlet{popblueL}{popblue!10!white}
\colorlet{poporangeL}{poporange!14!white}
\colorlet{popgreenL}{popgreen!12!white}
\colorlet{poppurpleL}{poppurple!12!white}
\colorlet{poppinkL}{poppink!12!white}
\colorlet{popgoldL}{popgold!14!white}

\pagecolor{popcream}

% ============================================================
% Typographie
% ============================================================
\defaultfontfeatures{Ligatures=TeX}
\setmainfont{PlusJakartaSans-Medium.ttf}[Path=../../../fonts/,BoldFont=PlusJakartaSans-Bold.ttf]
% Maths en sans-serif (Fira Math) avec chiffres et lettres droites de Plus
% Jakarta, pour que les formules se fondent dans le texte.
\usepackage[math-style=ISO,bold-style=ISO]{unicode-math}
\setmathfont{FiraMath-Regular.otf}[Path=../../../fonts/]
\setmathfont{PlusJakartaSans-Medium.ttf}[Path=../../../fonts/,range={up/num,up/{Latin,latin},\mathup}]
\usepackage{icomma} % virgule décimale sans espace en mode math
% Caractères absents de la police de texte (sinon carré vide) : rendus par
% leur équivalent mathématique, en texte comme en formule.
\usepackage{newunicodechar}
\newunicodechar{α}{\ensuremath{\alpha}}
\newunicodechar{Α}{\ensuremath{\alpha}}
\newunicodechar{β}{\ensuremath{\beta}}
\newunicodechar{δ}{\ensuremath{\delta}}
\newunicodechar{✓}{\popcheck{popgreen}}
\newfontfamily\popdisplay{ArchivoBlack-Regular.ttf}[Path=../../../fonts/]
\newfontfamily\poplogo{SpaceGrotesk-Bold.ttf}[Path=../../../fonts/]
\newfontfamily\popsemi{PlusJakartaSans-SemiBold.ttf}[Path=../../../fonts/]
\newfontfamily\popxbold{PlusJakartaSans-ExtraBold.ttf}[Path=../../../fonts/]
\newfontfamily\popxboldls{PlusJakartaSans-ExtraBold.ttf}[Path=../../../fonts/,LetterSpace=3.0]

\setlist[enumerate]{leftmargin=1.7em,itemsep=5pt,topsep=4pt}
\setlist[itemize]{leftmargin=1.7em,itemsep=4pt,topsep=4pt,label={\color{poporange}\textbullet}}
\setlength{\parindent}{0pt}
\setlength{\parskip}{5pt}
\renewcommand{\arraystretch}{1.25}

% chemfig : liaisons un peu plus épaisses, encre Pop
\setchemfig{atom sep=2em,bond style={line width=0.9pt,popink},cram width=3pt}

% pgfplots : style Pop par défaut
\pgfplotsset{
  every axis/.append style={axis line style={popink,line width=1pt},tick style={popink},
    grid style={popline,line width=0.5pt},label style={font=\small},tick label style={font=\footnotesize}},
  popcurve/.style={poporange,line width=1.8pt,no markers,smooth},
}

% ============================================================
% Pill / squiggle
% ============================================================
\newcommand{\poppill}[3]{%
  \tikz[baseline=-0.55ex]\node[draw=popink,line width=1.2pt,rounded corners=2.6pt,%
    fill=#2,inner xsep=6.5pt,inner ysep=3pt]{%
    \popxboldls\fontsize{7}{7}\selectfont\color{#3}\MakeUppercase{#1}};%
}
\newcommand{\popsquiggle}[1]{%
  \tikz[baseline=-0.4ex]{
    \draw[#1,line width=2.6pt,line cap=round]
      plot [smooth] coordinates {(0,0.09) (0.34,0.01) (0.68,0.21) (1.02,0.01) (1.36,0.10)};
  }%
}
% Mot-clé surligné (marqueur orange sous le texte)
\newcommand{\cle}[1]{{\popxbold\color{popdark}#1}}

% ============================================================
% En-tête / pied
% ============================================================
\pagestyle{fancy}
\fancyhf{}
\renewcommand{\headrulewidth}{0pt}
\renewcommand{\footrulewidth}{0pt}
\lhead{\raisebox{-0.15ex}{\poplogo\fontsize{13}{13}\selectfont\color{popink}OnBuch\color{poporange}+}}
\rhead{\poppill{\DOCMATIERE\ \textbullet\ \DOCNIVEAU\ \textbullet\ Leçon \DOCLECON}{white}{popink}}
\lfoot{\popxbold\fontsize{7.6}{7.6}\selectfont\color{popmuted}\MakeUppercase{Cours signé OnBuch+}\ \textbullet\ {\popsemi\DOCTITRE}}
\rfoot{\tikz[baseline=-0.6ex]\node[rounded corners=8.5pt,fill=popink,minimum height=17pt,minimum width=17pt,inner xsep=5pt,inner sep=0]{\popxbold\fontsize{8}{8}\selectfont\color{popcream}\thepage\,/\,\pageref*{LastPage}};}

% ============================================================
% Titres
% ============================================================
\newcommand{\chaptitre}[2]{%
  \begin{tcolorbox}[enhanced,colback=poporange,colframe=popink,boxrule=2.6pt,arc=11pt,
      drop shadow=popink,left=15pt,right=15pt,top=12pt,bottom=11pt,before skip=3pt,after skip=18pt]
    \poppill{#2}{popink}{popcream}\par\vspace{9pt}
    {\raggedright\hyphenpenalty=10000\exhyphenpenalty=10000\popdisplay\fontsize{22}{24}\selectfont\color{popcream}\MakeUppercase{#1}\par}
  \end{tcolorbox}
}
% Section numérotée : pastille orange avec le numéro + titre Archivo
\newcounter{popsec}
\newcommand{\coursec}[1]{%
  \stepcounter{popsec}\setcounter{popsub}{0}%
  \par\vspace{16pt}\noindent
  \begin{minipage}{\linewidth}
  \tikz[baseline=-0.7ex]\node[draw=popink,line width=1.6pt,fill=poporange,rounded corners=4pt,minimum size=22pt,inner sep=0]{\popdisplay\fontsize{12}{12}\selectfont\color{popcream}\thepopsec};%
  \hspace{8pt}{\popdisplay\fontsize{15}{17}\selectfont\color{popink}\MakeUppercase{#1}}\par
  \vspace{4pt}\noindent{\color{poporange}\rule{36pt}{3.6pt}}
  \end{minipage}\par\nopagebreak\vspace{8pt}
}
\newcounter{popsub}
\newcommand{\courssub}[1]{%
  \stepcounter{popsub}%
  \par\vspace{9pt}\noindent{\popxbold\fontsize{11.5}{13}\selectfont\color{popink}{\color{poporange}\thepopsec.\thepopsub}\hspace{6pt}#1}\par\nopagebreak\vspace{4pt}
}

% ============================================================
% Boîtes pédagogiques — même recette : fond blanc, bordure épaisse colorée,
% ombre dure encre, badge plein. Titre optionnel [#1].
% ============================================================
\tcbset{popbase/.style={breakable,enhanced,colback=white,boxrule=2.3pt,arc=9pt,
  shadow={3pt}{-3pt}{0pt}{popink},left=14pt,right=14pt,top=11pt,bottom=11pt,before skip=11pt,after skip=15pt,
  fontupper=\popsemi\fontsize{10.6}{14.5}\selectfont\color{popink},
  fontlower=\popsemi\fontsize{10.6}{14.5}\selectfont\color{popink}}}
% Coche dessinée (le glyphe ✓ n'existe pas dans les polices du document)
\newcommand{\popcheck}[1]{\tikz[baseline=-0.5ex]\draw[#1,line width=1.6pt,line cap=round,line join=round](-0.13,0.0)--(-0.04,-0.09)--(0.13,0.1);}
\newcommand{\popboxhead}[3]{\poppill{#1}{#2}{white}\ifx\relax#3\relax\else\hspace{6pt}{\popxbold\fontsize{10.6}{12}\selectfont\color{popink}#3}\fi\par\vspace{7pt}}

\newtcolorbox{aretenir}[1][]{popbase,colframe=popgreen,before upper={\popboxhead{À retenir}{popgreen}{#1}}}
\newtcolorbox{exemplebox}[1][]{popbase,colframe=poporange,before upper={\popboxhead{Exemple}{poporange}{#1}}}
\newtcolorbox{definition}[1][]{popbase,colframe=popblue,before upper={\popboxhead{Définition}{popblue}{#1}}}
\newtcolorbox{propriete}[1][]{popbase,colframe=poppurple,before upper={\popboxhead{Propriété}{poppurple}{#1}}}
\newtcolorbox{methode}[1][]{popbase,colframe=popgold,before upper={\popboxhead{Méthode}{popgold}{#1}}}
\newtcolorbox{attention}[1][]{popbase,colframe=poppink,before upper={\popboxhead{Attention}{poppink}{#1}}}
\newtcolorbox{experience}[1][]{popbase,colframe=popblue,colback=popblueL,before upper={\popboxhead{Expérience}{popblue}{#1}}}
% « Le savais-tu ? » : carte violette pleine (contexte, culture, Cameroun)
\newtcolorbox{savaistu}[1][]{popbase,colback=poppurple,colframe=popink,
  fontupper=\popsemi\fontsize{10.4}{14.2}\selectfont\color{white},
  before upper={\poppill{Le savais-tu\,?}{white}{poppurple}\ifx\relax#1\relax\else\hspace{6pt}{\popxbold\color{white}#1}\fi\par\vspace{7pt}}}
% Exercice résolu : énoncé en haut, \tcblower puis la solution sur fond crème
\newcounter{popexo}\newcounter{popexr}
\newtcolorbox{exoresolu}[1][]{popbase,colframe=poporange,colbacklower=popcream,
  segmentation style={popink,line width=1.2pt,dashed},
  before upper={\stepcounter{popexr}\popboxhead{Exercice résolu \thepopexr}{poporange}{#1}},
  before lower={\poppill{Solution}{popink}{popcream}\par\vspace{6pt}}}
% Exercice (non résolu en place) — la correction est dans la partie Corrigés
\newtcolorbox{exercice}[1][]{popbase,colframe=popink,
  before upper={\stepcounter{popexo}\popboxhead{Exercice \thepopexo}{popink}{#1}}}
% Corrigé
\newtcolorbox{corrige}[1][]{popbase,colframe=popgreen,colback=popgreenL,
  before upper={\popboxhead{Corrigé}{popgreen}{#1}}}
% Objectifs / prérequis / bilan
\newtcolorbox{objectifs}{popbase,colframe=popink,colback=poporangeL,
  before upper={\popboxhead{Objectifs de la leçon}{popink}{}}}
\newtcolorbox{prerequis}{popbase,colframe=popink,colback=popblueL,
  before upper={\popboxhead{Prérequis}{popblue}{}}}
\newtcolorbox{bilan}{popbase,colframe=popink,colback=white,boxrule=2.8pt,
  before upper={\popboxhead{Fiche bilan}{poporange}{L'essentiel à connaître pour l'examen}}}
% Figure encadrée avec légende
\newtcolorbox{popfigure}[1][]{enhanced,breakable=false,colback=white,colframe=popline,boxrule=1.4pt,arc=8pt,
  left=10pt,right=10pt,top=10pt,bottom=8pt,before skip=10pt,after skip=12pt,halign=center,
  lower separated=false,halign lower=center,
  fontlower=\popsemi\fontsize{9}{11.5}\selectfont\color{popmuted},
  #1}
\newcommand{\legende}[1]{\par\vspace{6pt}{\popsemi\fontsize{9}{11.5}\selectfont\color{popmuted}#1}}

% ============================================================
% Signature OnBuch+
% ============================================================
\newcommand{\poplogoinline}[1]{{\poplogo\fontsize{#1}{#1}\selectfont\color{popink}OnBuch\color{poporange}+}}
\newcommand{\signatureonbuch}{%
  \begin{tcolorbox}[enhanced,colback=popink,colframe=popink,boxrule=2.2pt,arc=10pt,
      drop shadow=poporange,left=14pt,right=14pt,top=10pt,bottom=10pt,before skip=0pt,after skip=0pt]
    \tikz[baseline=-0.7ex]\node[circle,fill=popgreen,draw=popcream,line width=1.3pt,minimum size=22pt,inner sep=0]{\popcheck{white}};%
    \hspace{9pt}\begin{minipage}[c]{0.62\linewidth}
      {\popxbold\fontsize{12}{13}\selectfont\color{popcream}Signé\ \textbullet\ {\poplogo OnBuch\color{poporange}+}}\par\vspace{2pt}
      {\raggedright\popsemi\fontsize{8}{10.5}\selectfont\color{popline}\MakeUppercase{Équipe pédagogique OnBuch+}\\\MakeUppercase{Conforme au programme officiel MINESEC}\par}
    \end{minipage}\hfill
    {\poplogo\fontsize{22}{22}\selectfont\color{popcream}OnBuch\color{poporange}+}
  \end{tcolorbox}%
}

% ============================================================
% Page de garde
% #1 titre · #2 matière · #3 niveau · #4 module · #5 n° leçon · #6 durée ·
% #7 description / objectifs (texte ou itemize)
% ============================================================
\newcommand{\pagegarde}[7]{%
  \thispagestyle{empty}
  \newgeometry{a4paper,margin=1.7cm,top=1.6cm,bottom=1.4cm}%
  \begin{tikzpicture}[remember picture,overlay]
    \fill[poporange!18] ([xshift=-3.2cm,yshift=-3.6cm]current page.north east) circle (4.6cm);
    \fill[poppurple!14] ([xshift=2.4cm,yshift=7.6cm]current page.south west) circle (3.8cm);
  \end{tikzpicture}%
  {\poplogo\fontsize{30}{30}\selectfont\color{popink}OnBuch\color{poporange}+}\hfill
  \raisebox{0.6ex}{\poppill{Cours signé \textbullet\ Premium}{popink}{popcream}}\par
  \vspace{1pt}\popsquiggle{poppurple}\par
  \vspace{8pt}
  \begin{tcolorbox}[enhanced,colback=poporange,colframe=popink,boxrule=2.8pt,arc=13pt,
      drop shadow=popink,left=18pt,right=18pt,top=12pt,bottom=13pt,after skip=10pt]
    \poppill{#2 \textbullet\ #3}{popink}{popcream}\hspace{5pt}\poppill{#4}{white}{popink}\par\vspace{8pt}
    {\popdisplay\fontsize{14}{14}\selectfont\color{popink}LEÇON #5}\par\vspace{4pt}
    {\raggedright\hyphenpenalty=10000\exhyphenpenalty=10000\popdisplay\fontsize{25}{27}\selectfont\color{popcream}\MakeUppercase{#1}\par}
    \vspace{8pt}
    {\popxbold\fontsize{9.5}{9.5}\selectfont\color{popink}\MakeUppercase{Durée conseillée : #6}}
  \end{tcolorbox}
  \begin{tcolorbox}[enhanced,colback=white,colframe=popink,boxrule=2.2pt,arc=10pt,
      drop shadow=popink,left=16pt,right=16pt,top=10pt,bottom=10pt,after skip=8pt]
    \poppill{Dans cette leçon}{poppurple}{white}\par\vspace{5pt}
    \popsemi\fontsize{9.3}{12.6}\selectfont\color{popink}#7
  \end{tcolorbox}
  \noindent\begin{minipage}[t]{0.487\linewidth}
    \begin{tcolorbox}[enhanced,colback=popgreen,colframe=popink,boxrule=2.2pt,arc=10pt,
        drop shadow=popink,left=11pt,right=11pt,top=10pt,bottom=10pt,height=3.1cm,valign=center]
      \tcbox[colback=white,colframe=popink,boxrule=1.3pt,arc=5pt,boxsep=0pt,nobeforeafter,
        left=3pt,right=3pt,top=3pt,bottom=3pt,valign=center]{\qrcode[height=1.9cm]{https://play.google.com/store/apps/details?id=com.luvvix.onbuch&hl=fr}}%
      \hspace{8pt}\begin{minipage}[c]{0.5\linewidth}
        {\popdisplay\fontsize{11}{12.5}\selectfont\color{white}\MakeUppercase{Télécharge OnBuch}}\par\vspace{4pt}
        {\raggedright\popsemi\fontsize{8.4}{11}\selectfont\color{white}Gratuit \textbullet\ Google Play\\Tuteur IA, annales, concours, résultats\par}
      \end{minipage}
    \end{tcolorbox}
  \end{minipage}\hfill
  \begin{minipage}[t]{0.487\linewidth}
    \begin{tcolorbox}[enhanced,colback=poppurple,colframe=popink,boxrule=2.2pt,arc=10pt,
        drop shadow=popink,left=11pt,right=11pt,top=10pt,bottom=10pt,height=3.1cm,valign=center]
      \tikz[baseline=-0.7ex]\node[circle,fill=white,draw=popink,line width=1.3pt,minimum size=1.6cm,inner sep=0]{\popdisplay\fontsize{24}{24}\selectfont\color{poppurple}+};%
      \hspace{8pt}\begin{minipage}[c]{0.6\linewidth}
        {\popdisplay\fontsize{11}{12.5}\selectfont\color{white}\MakeUppercase{Passe à OnBuch+}}\par\vspace{4pt}
        {\raggedright\popsemi\fontsize{8.4}{11}\selectfont\color{white}Tous les cours signés, Tuteur IA illimité, fiches PDF — onglet OnBuch+ de l'app\par}
      \end{minipage}
    \end{tcolorbox}
  \end{minipage}
  \vspace{8pt}
  \popcontenu
  \vfill
  \signatureonbuch
  \restoregeometry
  \newpage
}

% Rangée « Ce cours contient » (page de garde)
\newcommand{\poptile}[3]{% #1 couleur #2 gros texte #3 légende
  \begin{tcolorbox}[enhanced,colback=#1,colframe=popink,boxrule=2pt,arc=9pt,shadow={2.5pt}{-2.5pt}{0pt}{popink},
      left=6pt,right=6pt,top=9pt,bottom=9pt,halign=center,valign=center,height=1.75cm,width=\linewidth,nobeforeafter]
    {\popdisplay\fontsize{12.5}{13}\selectfont\color{white}\MakeUppercase{#2}}\par\vspace{3pt}
    {\centering\popsemi\fontsize{7.8}{9.5}\selectfont\color{white}#3\par}
  \end{tcolorbox}}
\newcommand{\popcontenu}{%
  \par\noindent{\popxboldls\fontsize{7.5}{7.5}\selectfont\color{popmuted}CE COURS SIGNÉ CONTIENT}\par\vspace{7pt}
  \noindent\begin{minipage}[t]{0.235\linewidth}\poptile{popblue}{Cours}{complet, illustré, pas à pas}\end{minipage}\hfill
  \begin{minipage}[t]{0.235\linewidth}\poptile{poporange}{Méthodes}{et exercices résolus}\end{minipage}\hfill
  \begin{minipage}[t]{0.235\linewidth}\poptile{poppink}{Exercices}{type examen + corrigés}\end{minipage}\hfill
  \begin{minipage}[t]{0.235\linewidth}\poptile{popgreen}{Fiche bilan}{l'essentiel à retenir}\end{minipage}\par}

% Sommaire
\newtcolorbox{sommaire}{enhanced,colback=white,colframe=popink,boxrule=2.2pt,arc=10pt,
  drop shadow=popink,left=18pt,right=18pt,top=13pt,bottom=13pt,before skip=6pt,after skip=20pt,
  before upper={\poppill{Sommaire}{poppurple}{white}\par\vspace{9pt}\popsemi\fontsize{10.6}{14.5}\selectfont\color{popink}}}

% Signature de fin de cours — poussée tout en bas de la dernière page
\newcommand{\findecours}{%
  \par\vfill\nopagebreak
  \begin{center}\popsquiggle{poporange}\end{center}\vspace{4pt}
  \signatureonbuch
}
% Glyphes absents de Fira Math : redéfinis à partir de glyphes disponibles
\AtBeginDocument{%
  \def\setminus{\mathbin{\backslash}}\let\smallsetminus\setminus
  \def\vdots{\mathord{\vbox{\baselineskip3.2pt\lineskiplimit0pt\kern4pt\hbox{.}\hbox{.}\hbox{.}}}}%
  \def\ddots{\mathinner{\mkern1mu\raise6.5pt\hbox{.}\mkern2mu\raise3.6pt\hbox{.}\mkern2mu\raise0.7pt\hbox{.}\mkern1mu}}%
  \def\triangle{\mathord{\scalebox{0.9}{$\symup{\Delta}$}}}%
  \def\nabla{\mathord{\raisebox{\depth}{\scalebox{1}[-1]{$\symup{\Delta}$}}}}%
  \def\checkmark{\popcheck{popdark}}%
  \def\star{\mathbin{\ast}}\def\bigstar{\mathord{\ast}}%
  \def\diamond{\mathbin{\scalebox{0.6}{$\Diamond$}}}%
  \def\bigcup{\mathop{\mathchoice{\raisebox{-0.3em}{\scalebox{1.7}{$\cup$}}}{\scalebox{1.25}{$\cup$}}{\cup}{\cup}}\displaylimits}%
  \def\bigcap{\mathop{\mathchoice{\raisebox{-0.3em}{\scalebox{1.7}{$\cap$}}}{\scalebox{1.25}{$\cap$}}{\cap}{\cap}}\displaylimits}%
  \def\lesseqqgtr{\mathrel{\lessgtr}}\def\bigoplus{\mathop{\oplus}\displaylimits}%
}
\providecommand{\ssi}{si et seulement si} % abréviation fréquente

%% FIGFIX-BEGIN v5 — correctifs globaux des figures (docs/onbuch-generation/tools/figfix)
\usepackage{adjustbox}\usepackage{etoolbox}\tcbuselibrary{raster}
% toute figure TikZ/pgfplots plus large que la ligne est réduite à la largeur disponible
\BeforeBeginEnvironment{tikzpicture}{\begin{adjustbox}{max width=\linewidth}}
\AfterEndEnvironment{tikzpicture}{\end{adjustbox}}
% légendes pgfplots lisibles par-dessus les courbes
\pgfplotsset{every axis/.append style={legend style={fill=white,fill opacity=0.92,text opacity=1,draw=black!15,inner sep=3pt}}}
% étiquettes d'arêtes (syntaxe edge["..."]) avec fond blanc
\tikzset{every edge quotes/.append style={fill=white,inner sep=1.2pt}}
%% FIGFIX-END
\begin{document}

\pagegarde{Stéréochimie : isomérie et représentation des molécules}{Chimie}{Tle TI}{Module 1 · Chimie organique}{05}{4 h}{Cette leçon explore les différentes façons dont des molécules de même formule brute peuvent différer : par leur enchaînement d'atomes (isomérie de constitution) ou par leur arrangement dans l'espace (stéréo-isomérie). Elle donne les outils de représentation (perspective, Fischer, Newman) indispensables pour décrire et comparer ces isomères au Baccalauréat TI.
\vspace{4pt}
\begin{multicols}{2}\small
\begin{itemize}
\item À la fin de cette leçon, je sais définir et distinguer les trois types d'isomérie de constitution (chaîne, position, fonction)
\item À la fin de cette leçon, je sais identifier un carbone asymétrique dans une molécule et justifier la chiralité
\item À la fin de cette leçon, je sais représenter une molécule chirale en perspective (Cram) et en projection de Fischer
\item À la fin de cette leçon, je sais distinguer deux énantiomères de deux diastéréoisomères
\item À la fin de cette leçon, je sais expliquer l'activité optique d'une solution et le rôle du polarimètre
\item À la fin de cette leçon, je sais attribuer la configuration Z ou E à un alcène et le justifier
\end{itemize}\end{multicols}}

\begin{sommaire}
\begin{multicols}{2}
\begin{enumerate}
\item Isomérie de constitution
\item Carbone asymétrique, chiralité et énantiomères
\item Activité optique et diastéréoisomères
\item Isomérie Z/E des alcènes
\item Stéréo-isomérie de conformation
\item Importance de la stéréo-isomérie et synthèse
\item Travaux pratiques
\item Méthodes et exercices résolus
\item Exercices
\item Corrigés des exercices
\item Fiche bilan
\end{enumerate}
\end{multicols}
\end{sommaire}

\begin{objectifs}
\begin{itemize}
\item À la fin de cette leçon, je sais définir et distinguer les trois types d'isomérie de constitution (chaîne, position, fonction)
\item À la fin de cette leçon, je sais identifier un carbone asymétrique dans une molécule et justifier la chiralité
\item À la fin de cette leçon, je sais représenter une molécule chirale en perspective (Cram) et en projection de Fischer
\item À la fin de cette leçon, je sais distinguer deux énantiomères de deux diastéréoisomères
\item À la fin de cette leçon, je sais expliquer l'activité optique d'une solution et le rôle du polarimètre
\item À la fin de cette leçon, je sais attribuer la configuration Z ou E à un alcène et le justifier
\item À la fin de cette leçon, je sais construire et interpréter une représentation de Newman (formes éclipsée et décalée)
\item À la fin de cette leçon, je sais représenter les conformations chaise et bateau du cyclohexane
\item À la fin de cette leçon, je sais citer des exemples concrets montrant l'importance de la stéréo-isomérie en pharmacie et dans la vie quotidienne
\end{itemize}
\end{objectifs}

\begin{prerequis}
\begin{itemize}
\item Formules brutes, semi-développées et développées des molécules organiques (Première)
\item Géométrie tétraédrique du carbone à 4 liaisons simples et géométrie plane du carbone à double liaison (Première)
\item Nomenclature des alcanes, alcènes et familles fonctionnelles (leçons 01 à 03)
\item Notion de polarisation de la lumière (Physique, Première)
\item Écriture et comptage des isomères à partir d'une formule brute (Première)
\end{itemize}
\end{prerequis}

\newpage

\coursec{Isomérie de constitution}

\textbf{Situation d'accroche.} En 1960, un médicament contre les nausées des femmes enceintes, la \cle{thalidomide}, a provoqué de graves malformations chez des milliers de bébés : une forme de la molécule soignait, son image miroir était toxique. Pourtant, ces deux molécules ont exactement la même formule brute ! De même, le \cle{limonène} sent l'orange sous une forme et le citron sous son image miroir. Comment deux molécules identiques sur le papier peuvent-elles avoir des effets si différents ? C'est toute la question de la \cle{stéréochimie}, c'est-à-dire de l'arrangement des atomes dans l'espace, que nous allons explorer dans cette leçon.

\textbf{Problématique.} Une formule brute comme $C_4H_{10}O$ ne dit rien de la manière dont les atomes sont reliés entre eux. Or, c'est précisément cet arrangement qui détermine l'odeur, le point d'ébullition, la toxicité ou l'activité médicamenteuse d'un composé. Comment classer et dénombrer toutes les molécules possibles associées à une même formule brute ? C'est l'objet de cette première section.

\courssub{La notion d'isomérie : intuition et définition}

Imagine que tu disposes de quatre briques rouges et dix briques blanches (les atomes de carbone et d'hydrogène du butane $C_4H_{10}$). Tu peux les assembler en une longue file droite, ou bien en une file de trois avec une brique greffée au milieu. Tu as utilisé exactement les mêmes briques, pourtant les deux constructions ne sont pas superposables : ce sont deux objets différents. En chimie, c'est exactement ce phénomène que l'on appelle l'\cle{isomérie}.

\begin{definition}[Isomères]
On appelle \cle{isomères} des composés qui possèdent la \textbf{même formule brute} mais des \textbf{structures différentes}, c'est-à-dire des arrangements d'atomes différents. Les isomères sont des espèces chimiques distinctes : elles ont des propriétés physiques (températures de fusion et d'ébullition, densité, solubilité) et des propriétés chimiques différentes.
\end{definition}

On distingue deux grandes familles d'isomérie :
\begin{itemize}
\item l'\cle{isomérie de constitution} : les isomères diffèrent par l'\textbf{ordre de liaison des atomes} (la chaîne carbonée, la position d'un groupe, la nature de la fonction) ;
\item la \cle{stéréo-isomérie} : les isomères ont le même enchaînement d'atomes mais diffèrent par leur \textbf{arrangement dans l'espace} (énantiomères, diastéréoisomères, isomères $Z$/$E$, conformations). Elle fera l'objet des sections suivantes.
\end{itemize}

Dans cette section, nous nous concentrons sur l'isomérie de constitution, qui se décline en trois types : l'isomérie de \textbf{chaîne}, l'isomérie de \textbf{position} et l'isomérie de \textbf{fonction}.

\courssub{Isomérie de chaîne}

\begin{definition}[Isomérie de chaîne]
Deux composés sont des \cle{isomères de chaîne} lorsqu'ils possèdent la même formule brute, la \textbf{même fonction chimique}, mais des \textbf{squelettes carbonés différents} (chaîne linéaire ou ramifiée).
\end{definition}

\textbf{Exemple fondateur : le butane $C_4H_{10}$.} La formule brute $C_4H_{10}$ correspond à deux alcanes différents :

\begin{center}
\chemfig{CH_3-CH_2-CH_2-CH_3} \hspace{2cm} \chemfig{CH_3-CH(-[2]CH_3)-CH_3}
\end{center}

\begin{itemize}
\item le \textbf{butane} (chaîne linéaire de 4 atomes de carbone) : $T_{\text{éb}} = \SI{-0,5}{\celsius}$ ;
\item le \textbf{2-méthylpropane} ou isobutane (chaîne ramifiée) : $T_{\text{éb}} = \SI{-11,7}{\celsius}$.
\end{itemize}

Ces deux gaz sont vendus mélangés dans les bouteilles de gaz domestique utilisées dans les cuisines camerounaises. Remarque déjà une conséquence capitale : la ramification \textbf{diminue} la température d'ébullition, car les molécules ramifiées, plus compactes, s'imbriquent moins bien les unes dans les autres, ce qui affaiblit les interactions intermoléculaires.

\begin{exemplebox}[Les trois isomères du pentane $C_5H_{12}$]
La formule brute $C_5H_{12}$ admet exactement \textbf{trois} isomères de chaîne :
\begin{center}
\begin{tabularx}{\linewidth}{|l|X|X|}
\hline
\rowcolor{popblueL}
\textbf{Nom} & \textbf{Formule semi-développée} & $T_{\text{éb}}$ (\si{\celsius}) \\
\hline
Pentane & \ce{CH3-CH2-CH2-CH2-CH3} & 36,1 \\
\hline
2-Méthylbutane & \ce{CH3-CH(CH3)-CH2-CH3} & 27,9 \\
\hline
2,2-Diméthylpropane & \ce{C(CH3)4} & 9,5 \\
\hline
\end{tabularx}
\end{center}
Plus la chaîne est ramifiée, plus la température d'ébullition est basse : c'est une règle générale et très utile pour vérifier la cohérence de tes résultats.
\end{exemplebox}

\courssub{Isomérie de position}

\begin{definition}[Isomérie de position]
Deux composés sont des \cle{isomères de position} lorsqu'ils possèdent la même formule brute, la \textbf{même chaîne carbonée} et la \textbf{même fonction chimique}, mais que le \textbf{groupe caractéristique} (ou la double liaison) n'occupe pas la même \textbf{position} sur la chaîne.
\end{definition}

\textbf{Exemple 1 : les alcools en $C_4H_{10}O$.} Le groupe hydroxyle \ce{-OH} peut occuper deux positions non équivalentes sur la chaîne butane :

\begin{center}
\chemfig{CH_3-CH_2-CH_2-CH_2-OH} \hspace{1.5cm} \chemfig{CH_3-CH(-[2]OH)-CH_2-CH_3}
\end{center}

\begin{itemize}
\item le \textbf{butan-1-ol} (alcool primaire, \ce{-OH} en bout de chaîne) : $T_{\text{éb}} = \SI{117}{\celsius}$ ;
\item le \textbf{butan-2-ol} (alcool secondaire, \ce{-OH} sur le carbone n° 2) : $T_{\text{éb}} = \SI{99,5}{\celsius}$.
\end{itemize}

\textbf{Exemple 2 : les alcènes en $C_4H_8$.} La double liaison peut aussi changer de place :

\begin{center}
\chemfig{CH_2=CH-CH_2-CH_3} \hspace{2cm} \chemfig{CH_3-CH=CH-CH_3}
\end{center}

On obtient le \textbf{but-1-ène} ($T_{\text{éb}} = \SI{-6,3}{\celsius}$) et le \textbf{but-2-ène} ($T_{\text{éb}} \approx \SI{1}{\celsius}$). Même squelette, même fonction, position différente : ce sont des isomères de position.

\begin{attention}[Ne pas confondre position et chaîne]
Le butan-2-ol \chemfig{CH_3-CH(-[2]OH)-CH_2-CH_3} et le 2-méthylpropan-1-ol \chemfig{CH_3-CH(-[2]CH_3)-CH_2-OH} ont tous deux la formule $C_4H_{10}O$, mais leurs squelettes carbonés diffèrent : ce sont des isomères de \textbf{chaîne}, pas de position. Avant de conclure, identifie toujours d'abord la chaîne la plus longue, puis la position du groupe caractéristique.
\end{attention}

\courssub{Isomérie de fonction}

\begin{definition}[Isomérie de fonction]
Deux composés sont des \cle{isomères de fonction} lorsqu'ils possèdent la même formule brute mais appartiennent à des \textbf{familles chimiques différentes} : leurs groupes caractéristiques sont différents.
\end{definition}

C'est le type d'isomérie le plus spectaculaire, car les propriétés chimiques changent du tout au tout.

\textbf{Exemple fondateur : $C_2H_6O$.} Cette formule brute correspond à deux composés que tout oppose :

\begin{center}
\chemfig{CH_3-CH_2-OH} \hspace{2cm} \chemfig{CH_3-O-CH_3}
\end{center}

\begin{itemize}
\item l'\textbf{éthanol} \ce{CH3CH2OH}, un alcool : liquide à température ambiante, $T_{\text{éb}} = \SI{78,4}{\celsius}$, soluble dans l'eau, réagit avec le sodium ;
\item le \textbf{méthoxyméthane} \ce{CH3-O-CH3}, un éther : gaz à température ambiante, $T_{\text{éb}} = \SI{-24}{\celsius}$, pratiquement insoluble dans l'eau, inerte vis-à-vis du sodium.
\end{itemize}

L'écart de plus de \SI{100}{\celsius} entre les températures d'ébullition s'explique par la présence, dans l'alcool seul, de liaisons hydrogène entre molécules.

\textbf{Autres couples classiques d'isomères de fonction :}

\begin{center}
\begin{tabularx}{\linewidth}{|l|X|X|}
\hline
\rowcolor{popblueL}
\textbf{Formule brute} & \textbf{Famille 1} & \textbf{Famille 2} \\
\hline
$C_3H_6O$ & Propanal (aldéhyde) \ce{CH3-CH2-CHO} & Propanone (cétone) \ce{CH3-CO-CH3} \\
\hline
$C_3H_6O_2$ & Acide propanoïque \ce{CH3-CH2-COOH} & Méthanoate d'éthyle (ester) \ce{H-COO-CH2-CH3} \\
\hline
$C_2H_6O$ & Éthanol (alcool) \ce{CH3-CH2-OH} & Méthoxyméthane (éther) \ce{CH3-O-CH3} \\
\hline
\end{tabularx}
\end{center}

Retiens les couples à connaître au baccalauréat : \textbf{aldéhyde/cétone}, \textbf{acide carboxylique/ester}, \textbf{alcool/éther}.

\begin{popfigure}
\begin{center}
\begin{tabularx}{\linewidth}{|l|X|X|}
\hline
\rowcolor{popgreenL}
\textbf{Type d'isomérie} & \textbf{Isomère A} & \textbf{Isomère B} \\
\hline
Chaîne ($C_4H_{10}$) & \chemfig{CH_3-CH_2-CH_2-CH_3} butane & \chemfig{CH_3-CH(-[2]CH_3)-CH_3} 2-méthylpropane \\
\hline
Position ($C_4H_8$) & \chemfig{CH_2=CH-CH_2-CH_3} but-1-ène & \chemfig{CH_3-CH=CH-CH_3} but-2-ène \\
\hline
Fonction ($C_2H_6O$) & \chemfig{CH_3-CH_2-OH} éthanol (alcool) & \chemfig{CH_3-O-CH_3} méthoxyméthane (éther) \\
\hline
\end{tabularx}
\end{center}
\legende{Les trois types d'isomérie de constitution : même formule brute à chaque ligne, mais chaîne, position ou fonction différentes.}
\end{popfigure}

\courssub{Méthode systématique de recherche de tous les isomères}

Face à une formule brute donnée, comment être sûr de n'oublier aucun isomère et de n'en compter aucun deux fois ? Il faut procéder de manière \textbf{ordonnée}.

\begin{methode}[Trouver tous les isomères de constitution d'une formule brute]
\begin{enumerate}
\item \textbf{Varier la chaîne} : écrire d'abord la chaîne la plus longue possible, puis raccourcir progressivement la chaîne principale en greffant les carbones retirés comme ramifications (méthyle, éthyle...). Attention : une ramification ne peut pas être placée sur un carbone terminal, sinon on rallonge la chaîne !
\item \textbf{Varier la position} : pour chaque squelette obtenu, déplacer le groupe caractéristique (ou la double liaison) sur toutes les positions non équivalentes. Deux positions sont équivalentes si elles donnent le même composé après renumérotation de la chaîne.
\item \textbf{Varier la fonction} : chercher les autres familles compatibles avec la formule brute (alcool/éther, aldéhyde/cétone, acide/ester).
\item \textbf{Vérifier en nommant} : nommer chaque structure obtenue selon les règles de nomenclature. Deux structures qui portent le même nom sont le même isomère : il faut en éliminer une.
\end{enumerate}
\end{methode}

\begin{popfigure}
\begin{center}
\begin{tikzpicture}[
box/.style={draw=popblue, fill=popblueL, rounded corners=3pt, text width=3.1cm, align=center, minimum height=1.1cm, font=\small},
dec/.style={draw=poporange, fill=poporangeL, rounded corners=3pt, text width=2.9cm, align=center, minimum height=1cm, font=\small\bfseries},
arr/.style={-{Stealth[length=3mm]}, thick, popdark},
node distance=0.55cm and 1.1cm]
\node[dec] (start) {Formule brute donnée};
\node[box, below=of start] (chaine) {1. Varier la \textbf{chaîne} : linéaire puis ramifiée};
\node[box, below=of chaine] (pos) {2. Varier la \textbf{position} du groupe caractéristique};
\node[box, below=of pos] (fonc) {3. Varier la \textbf{fonction} (alcool/éther, aldéhyde/cétone...)};
\node[box, below=of fonc] (nom) {4. \textbf{Nommer} chaque structure pour éliminer les doublons};
\node[dec, right=of fonc] (fin) {Liste complète des isomères};
\draw[arr] (start) -- (chaine);
\draw[arr] (chaine) -- (pos);
\draw[arr] (pos) -- (fonc);
\draw[arr] (fonc) -- (nom);
\draw[arr] (nom.east) -| (fin.south);
\end{tikzpicture}
\end{center}
\legende{Arbre de décision : démarche systématique de recherche de tous les isomères de constitution d'une formule brute.}
\end{popfigure}

\begin{attention}[Le piège du dessin tourné]
L'erreur la plus fréquente consiste à dessiner deux fois le même isomère en le tournant ou en l'écrivant de droite à gauche. Par exemple, \ce{CH3-CH2-CH2-OH} et \ce{HO-CH2-CH2-CH3} sont \textbf{la même molécule} (propan-1-ol), simplement lue dans l'autre sens. De même, le butan-2-ol écrit \ce{CH3-CH(OH)-CH2-CH3} ou \ce{CH3-CH2-CH(OH)-CH3} est un seul et même composé. \textbf{Parade infaillible : nomme chaque structure.} Deux dessins portant le même nom sont un seul isomère.
\end{attention}

\begin{exoresolu}[Dénombrer les isomères de $C_4H_{10}O$]
Énoncé. Déterminer le nombre total d'isomères de constitution de formule brute $C_4H_{10}O$, en précisant leur nom et leur famille.
\tcblower
Solution détaillée. La formule $C_4H_{10}O$ est saturée (forme générale $C_nH_{2n+2}O$) : il s'agit donc d'\textbf{alcools} ou d'\textbf{éthers}.

\textbf{Étape 1 : les alcools.} On varie la chaîne, puis la position de \ce{-OH}.
\begin{itemize}
\item Chaîne butane : \ce{-OH} en position 1 $\Rightarrow$ \textbf{butan-1-ol} ; en position 2 $\Rightarrow$ \textbf{butan-2-ol} (les positions 3 et 4 redonnent les mêmes par renumérotation).
\item Chaîne 2-méthylpropane : \ce{-OH} sur un carbone terminal $\Rightarrow$ \textbf{2-méthylpropan-1-ol} ; sur le carbone central $\Rightarrow$ \textbf{2-méthylpropan-2-ol}.
\end{itemize}
Soit \textbf{4 alcools}.

\textbf{Étape 2 : les éthers.} On répartit les 4 carbones de part et d'autre de l'oxygène :
\begin{itemize}
\item \ce{CH3-O-CH2-CH2-CH3} : \textbf{1-méthoxypropane} ;
\item \ce{CH3-O-CH(CH3)-CH3} : \textbf{2-méthoxypropane} ;
\item \ce{CH3-CH2-O-CH2-CH3} : \textbf{éthoxyéthane}.
\end{itemize}
Soit \textbf{3 éthers}.

\textbf{Conclusion :} $C_4H_{10}O$ admet au total $4 + 3 = \boxed{7}$ isomères de constitution. Les alcools et les éthers sont entre eux des isomères de fonction ; au sein de chaque famille, on trouve des isomères de chaîne et de position.
\end{exoresolu}

\courssub{Conséquences sur les propriétés physiques et chimiques}

\begin{aretenir}[Les isomères sont des espèces différentes]
Des isomères de constitution sont des \textbf{espèces chimiques distinctes} :
\begin{itemize}
\item leurs \textbf{propriétés physiques} diffèrent : températures de fusion et d'ébullition, densité, solubilité (ex. butane $\SI{-0,5}{\celsius}$ contre 2-méthylpropane $\SI{-11,7}{\celsius}$ ; éthanol $\SI{78,4}{\celsius}$ contre méthoxyméthane $\SI{-24}{\celsius}$) ;
\item leurs \textbf{propriétés chimiques} diffèrent, surtout entre isomères de fonction : un alcool réagit avec le sodium en dégageant du dihydrogène, un éther non ; un aldéhyde réduit la liqueur de Fehling, une cétone non.
\end{itemize}
C'est pourquoi les tests d'identification des fonctions (Fehling, DNPH, sodium, pH) permettent de distinguer expérimentalement des isomères de fonction.
\end{aretenir}

\begin{savaistu}[Glucose et fructose : le sucre a ses isomères]
Le glucose et le fructose ont tous deux la formule brute $C_6H_{12}O_6$, mais le glucose est un \textbf{aldéhyde} (aldose) tandis que le fructose est une \textbf{cétone} (cétose) : ce sont des isomères de fonction ! Or le fructose, présent dans les fruits et le miel, est presque \textbf{deux fois plus sucré} que le glucose. C'est pourquoi les jus de bissap ou de foléré préparés avec des fruits mûrs semblent plus doux : la nature joue elle aussi avec l'isomérie. Ce simple changement de fonction suffit à modifier radicalement la perception de notre palais, preuve éclatante que la structure fait la propriété.
\end{savaistu}

\begin{exercice}[À toi de jouer]
\begin{enumerate}
\item Écrire les formules semi-développées et nommer tous les isomères de constitution de formule brute $C_4H_8$ (alcènes uniquement). Préciser le type d'isomérie qui les relie deux à deux.
\item L'acide butanoïque et le méthanoate de propyle ont-ils la même formule brute ? Sont-ils isomères ? Si oui, de quel type ?
\end{enumerate}
\end{exercice}

\begin{corrige}[Exercice]
\textbf{1.} Trois alcènes répondent à la formule $C_4H_8$ : le but-1-ène \ce{CH2=CH-CH2-CH3}, le but-2-ène \ce{CH3-CH=CH-CH3} et le 2-méthylpropène \ce{CH2=C(CH3)-CH3}. But-1-ène et but-2-ène sont des isomères de \textbf{position} ; le 2-méthylpropène est un isomère de \textbf{chaîne} des deux précédents.

\textbf{2.} Acide butanoïque \ce{CH3-CH2-CH2-COOH} : $C_4H_8O_2$. Méthanoate de propyle \ce{H-COO-CH2-CH2-CH3} : $C_4H_8O_2$. Même formule brute, fonctions différentes (acide carboxylique et ester) : ce sont des \textbf{isomères de fonction}.
\end{corrige}

\textbf{En résumé.} Une formule brute peut cacher plusieurs molécules : les isomères de constitution se classent en trois types — chaîne, position, fonction — que l'on explore systématiquement dans cet ordre, en nommant chaque structure pour éviter les doublons. Ces isomères ont des propriétés physiques et chimiques distinctes. Mais l'histoire de la thalidomide nous l'a rappelé : même des molécules de \emph{même constitution} peuvent différer dans l'espace. C'est l'objet de la section suivante, où nous découvrirons le carbone asymétrique et la chiralité.

\coursec{Carbone asymétrique, chiralité et énantiomères}

Dans la section précédente, nous avons découvert que deux molécules peuvent posséder la \emph{même formule brute} tout en différant par la \emph{constitution} : l'ordre d'enchaînement des atomes (isomérie de chaîne, de position ou de fonction). Mais il existe un niveau de différence encore plus subtil. Deux molécules peuvent avoir la même formule brute \emph{et} le même enchaînement d'atomes, et pourtant ne pas être identiques, parce que leurs atomes ne sont pas disposés de la même façon \emph{dans l'espace}. C'est le domaine de la \cle{stéréo-isomérie}.

Pour bien comprendre, prends tes deux mains. Elles ont exactement les mêmes « constituants » : un pouce, quatre doigts, une paume. Pourtant, tu ne pourras jamais superposer ta main gauche sur ta main droite paume contre paume : les pouces pointent dans des directions opposées. Tes deux mains sont \emph{images l'une de l'autre dans un miroir} et \emph{non superposables}. Certaines molécules se comportent exactement comme cela. Cette propriété, appelée \cle{chiralité}, est au cœur de cette section et de toute la chimie du vivant.

\courssub{Le carbone asymétrique}

\textbf{Intuition : un carbone tétraédrique}\par

Rappelons d'abord un résultat fondamental : un atome de carbone formant quatre liaisons simples est \cle{tétraédrique}. Les quatre liaisons partent du carbone vers les sommets d'un tétraèdre régulier, avec des angles d'environ $109,5\degres$. Le carbone n'est donc \emph{pas} plat : ses quatre substituants occupent des positions bien précises dans l'espace.

Imaginons maintenant que ces quatre substituants soient \emph{tous différents}. Alors il existe exactement \emph{deux façons} de les disposer autour du carbone, et ces deux arrangements sont images l'un de l'autre dans un miroir sans être superposables. Un tel carbone est dit \emph{asymétrique}.

\begin{definition}[Carbone asymétrique]
Un \cle{carbone asymétrique} est un atome de carbone \textbf{tétraédrique} (engagé dans quatre liaisons simples) lié à \textbf{quatre atomes ou groupes d'atomes tous différents}. On le repère par un astérisque : $C^*$.
\end{definition}

\begin{methode}[Repérer un carbone asymétrique]
Pour repérer un carbone asymétrique dans une molécule, procède en deux étapes :
\begin{enumerate}
\item \textbf{Vérifier que le carbone est tétraédrique} : il doit porter quatre liaisons \emph{simples}. Un carbone engagé dans une double liaison ($\ce{C=C}$, $\ce{C=O}$) ou une triple liaison est plan ou linéaire : il ne peut \emph{jamais} être asymétrique.
\item \textbf{Comparer les quatre substituants} : ils doivent être \emph{tous différents}. Dès que deux substituants sont identiques (par exemple deux groupes $\ce{CH3}$ ou deux atomes H), le carbone n'est \emph{pas} asymétrique.
\end{enumerate}
\end{methode}

\begin{exemplebox}[L'acide lactique : un carbone asymétrique]
L'\cle{acide lactique}, présent dans le lait caillé et dans les muscles après un effort, a pour formule semi-développée $\ce{CH3-CHOH-COOH}$, c'est-à-dire \chemfig{CH_3-\chembelow{C}{\scriptstyle *}(-[2]OH)(-[6]H)-COOH}. Le carbone central est lié à quatre groupes différents : $\ce{-H}$, $\ce{-OH}$, $\ce{-CH3}$ et $\ce{-COOH}$. C'est un carbone asymétrique $C^*$. En revanche, le carbone du groupe $\ce{-CH3}$ porte trois hydrogènes identiques, et celui du groupe $\ce{-COOH}$ porte une double liaison : aucun des deux n'est asymétrique.
\end{exemplebox}

\begin{attention}[Pièges fréquents]
\begin{itemize}
\item Un carbone portant une \textbf{double liaison} ($\ce{C=O}$ d'un acide carboxylique, d'un aldéhyde ou d'une cétone ; $\ce{C=C}$ d'un alcène) n'est \emph{jamais} asymétrique.
\item Attention aux groupes qui \emph{se ressemblent} sans être identiques : dans $\ce{CH3-CHBr-CH2-CH3}$, le carbone 2 porte $\ce{-H}$, $\ce{-Br}$, $\ce{-CH3}$ et $\ce{-CH2-CH3}$ : quatre groupes \emph{différents}, donc $C^*$. Ne conclus pas trop vite que $\ce{-CH3}$ et $\ce{-CH2-CH3}$ seraient « pareils » !
\item La présence d'un $C^*$ suffit en général à rendre la molécule chirale, mais il existe des cas particuliers (molécules \emph{méso}) qui seront évoqués plus loin.
\end{itemize}
\end{attention}

\courssub{Chiralité d'une molécule}

\textbf{L'analogie des mains}\par

Le mot \emph{chiral} vient du grec \emph{kheir}, « la main ». Une molécule est dite \cle{chirale} lorsqu'elle n'est \emph{pas superposable à son image dans un miroir}, exactement comme ta main gauche n'est pas superposable à ta main droite. À l'inverse, un objet superposable à son image miroir (une cuillère, une balle de football) est dit \emph{achiral}.

\begin{definition}[Chiralité]
Une molécule est \cle{chirale} si elle n'est \textbf{pas superposable} à son image dans un miroir plan. En pratique, au niveau de la Terminale TI, une molécule possédant \textbf{un seul carbone asymétrique} est toujours chirale.
\end{definition}

\begin{popfigure}
\begin{center}
\begin{tikzpicture}[scale=0.9, line cap=round, line join=round]
% ---- Main gauche ----
\begin{scope}[xshift=-3.2cm, popblue]
\fill[popblueL, draw=popblue, thick]
  (0,0) .. controls (-0.9,0.1) and (-1.0,1.0) .. (-0.75,1.2) % petit doigt
  .. controls (-0.55,1.35) .. (-0.5,1.1)
  -- (-0.35,1.75) .. controls (-0.3,2.0) .. (-0.12,1.95) .. controls (-0.05,1.9) .. (-0.08,1.7)
  -- (0.05,2.25) .. controls (0.08,2.5) .. (0.25,2.45) .. controls (0.32,2.4) .. (0.3,2.2)
  -- (0.4,1.85) .. controls (0.45,2.05) .. (0.6,2.0) .. controls (0.68,1.95) .. (0.65,1.7)
  -- (0.7,0.9) .. controls (1.15,0.75) and (1.25,0.35) .. (1.05,0.15) % pouce
  .. controls (0.95,0.0) .. (0.7,0.1)
  .. controls (0.6,-0.5) and (0.3,-0.9) .. (0,-0.95)
  .. controls (-0.4,-0.9) and (-0.7,-0.5) .. (0,0);
\end{scope}
\node[popdark] at (-3.2,-1.5) {\small Main gauche};
% ---- Miroir ----
\draw[popmuted, thick, decorate, decoration={coil, aspect=0.4, segment length=4pt}] (0,-1.6) -- (0,2.7);
\node[popmuted, rotate=90] at (0.25,0.5) {\scriptsize miroir};
% ---- Main droite (image miroir) ----
\begin{scope}[xshift=3.2cm, xscale=-1, poporange]
\fill[poporangeL, draw=poporange, thick]
  (0,0) .. controls (-0.9,0.1) and (-1.0,1.0) .. (-0.75,1.2)
  .. controls (-0.55,1.35) .. (-0.5,1.1)
  -- (-0.35,1.75) .. controls (-0.3,2.0) .. (-0.12,1.95) .. controls (-0.05,1.9) .. (-0.08,1.7)
  -- (0.05,2.25) .. controls (0.08,2.5) .. (0.25,2.45) .. controls (0.32,2.4) .. (0.3,2.2)
  -- (0.4,1.85) .. controls (0.45,2.05) .. (0.6,2.0) .. controls (0.68,1.95) .. (0.65,1.7)
  -- (0.7,0.9) .. controls (1.15,0.75) and (1.25,0.35) .. (1.05,0.15)
  .. controls (0.95,0.0) .. (0.7,0.1)
  .. controls (0.6,-0.5) and (0.3,-0.9) .. (0,-0.95)
  .. controls (-0.4,-0.9) and (-0.7,-0.5) .. (0,0);
\end{scope}
\node[popdark] at (3.2,-1.5) {\small Main droite (image miroir)};
% ---- Mention non superposable ----
\node[align=center, popdark] at (0,-2.3) {\small Les deux mains sont \textbf{images miroir} et \textbf{non superposables} :};
\node[align=center, popdark] at (0,-2.8) {\small impossible de faire coïncider paumes, doigts et pouces à la fois.};
\end{tikzpicture}
\legende{Analogie fondamentale : la main gauche et la main droite ont la même « constitution » mais ne sont pas superposables. Une molécule chirale se comporte de la même manière vis-à-vis de son image miroir.}
\end{center}
\end{popfigure}

\courssub{Les énantiomères}

Lorsqu'une molécule est chirale, elle existe sous \emph{deux formes} : la molécule elle-même et son image miroir, qui est une molécule \emph{différente} puisqu'on ne peut pas les superposer. Ces deux formes constituent un couple d'\emph{énantiomères}.

\begin{definition}[Énantiomères]
Deux stéréo-isomères sont des \cle{énantiomères} s'ils sont \textbf{images l'un de l'autre dans un miroir} et \textbf{non superposables}. Ils ont la même formule brute, le même enchaînement d'atomes, mais des arrangements spatiaux différents.
\end{definition}

\begin{popfigure}
\begin{center}
\setchemfig{atom sep=2.2em}
\chemfig{C(-[2]H)(-[6]CH_3)(<:[4]HO)(<OH)-[0]C(=[1]O)-[7]OH}
\hspace{1.2cm}
\begin{tikzpicture}[baseline={(current bounding box.center)}]
\draw[popmuted, thick, decorate, decoration={coil, aspect=0.4, segment length=4pt}] (0,-1.4) -- (0,1.6);
\node[popmuted, rotate=90] at (0.28,0.1) {\scriptsize miroir};
\end{tikzpicture}
\hspace{1.2cm}
\chemfig{C(-[2]H)(-[6]CH_3)(<:[4]OH)(<HO)-[0]C(=[1]O)-[7]OH}
\legende{Les deux énantiomères de l'acide lactique $\ce{CH3-CHOH-COOH}$ en représentation de Cram : le carbone central est le $C^*$ ; les liaisons en triangle plein sortent du plan vers l'avant, les liaisons hachurées partent vers l'arrière. Les deux formes sont images miroir et non superposables.}
\end{center}
\end{popfigure}

\begin{aretenir}[Nombre maximal de stéréo-isomères]
Une molécule possédant $n$ carbones asymétriques admet \textbf{au maximum} $2^n$ stéréo-isomères :
\[ N_{\max} = 2^n \]
Ce nombre est un \emph{maximum} : il peut être inférieur si la molécule possède des éléments de symétrie (cas des composés \emph{méso}).
\end{aretenir}

\begin{exemplebox}[Exemple chiffré : le 2-bromobutane]
Le 2-bromobutane $\ce{CH3-CHBr-CH2-CH3}$ possède \textbf{un seul} carbone asymétrique (le carbone 2, lié à $\ce{-H}$, $\ce{-Br}$, $\ce{-CH3}$ et $\ce{-CH2-CH3}$). Il admet donc :
\[ N_{\max} = 2^1 = 2 \ \text{stéréo-isomères} \]
Ces deux stéréo-isomères sont deux énantiomères, images miroir non superposables. De même, l'acide tartrique $\ce{HOOC-CHOH-CHOH-COOH}$ possède $n = 2$ carbones asymétriques, donc au maximum $2^2 = 4$ stéréo-isomères (en réalité 3, car l'un d'eux est un composé méso, achiral par symétrie interne).
\end{exemplebox}

\begin{exoresolu}[Dénombrer les stéréo-isomères du 3-méthylhexane]
Énoncé : on considère le 3-méthylhexane de formule $\ce{CH3-CH2-CH(CH3)-CH2-CH2-CH3}$. Identifier le(s) carbone(s) asymétrique(s) et déterminer le nombre maximal de stéréo-isomères.
\tcblower
Solution détaillée :
\begin{enumerate}
\item \textbf{Repérage des carbones tétraédriques} : tous les carbones de cet alcane portent quatre liaisons simples, ils sont tous candidats.
\item \textbf{Examen de chaque carbone} :
\begin{itemize}
\item C1, C2, C5, C6 et le $\ce{CH3}$ en position 3 portent au moins deux hydrogènes identiques : pas asymétriques.
\item C4 est un groupe $\ce{CH2}$ : il porte deux atomes d'hydrogène identiques, donc pas asymétrique.
\item C3 est lié à : $\ce{-H}$, $\ce{-CH3}$ (le substituant méthyle), $\ce{-CH2-CH3}$ (vers C2-C1) et $\ce{-CH2-CH2-CH3}$ (vers C4-C5-C6). Ces quatre groupes sont \emph{tous différents} : C3 est un carbone asymétrique $C^*$.
\end{itemize}
\item \textbf{Conclusion} : $n = 1$, donc $N_{\max} = 2^1 = 2$ stéréo-isomères, qui forment un couple d'énantiomères.
\end{enumerate}
\end{exoresolu}

\courssub{Propriétés des énantiomères}

Voici le point le plus déroutant pour un élève : deux énantiomères sont des molécules \emph{différentes}, et pourtant elles se ressemblent énormément.

\begin{propriete}[Propriétés des énantiomères]
Deux énantiomères possèdent :
\begin{itemize}
\item les \textbf{mêmes propriétés physiques} : température d'ébullition $T_{\text{éb}}$, température de fusion $T_{\text{fus}}$, densité, solubilité, indice de réfraction ;
\item les \textbf{mêmes propriétés chimiques} vis-à-vis de réactifs \emph{achiraux} (ils réagissent de la même façon, avec la même vitesse) ;
\item des \textbf{comportements différents} dans deux situations seulement :
\begin{enumerate}
\item vis-à-vis de la \textbf{lumière polarisée} : ils dévient le plan de polarisation en sens opposés (activité optique, étudiée dans la section suivante) ;
\item vis-à-vis de \textbf{réactifs ou d'environnements chiraux} : enzymes, récepteurs biologiques, autres molécules chirales.
\end{enumerate}
\end{itemize}
\end{propriete}

C'est cette dernière différence qui explique l'importance capitale de la stéréochimie en biologie et en pharmacie : une enzyme étant elle-même chirale, elle « reconnaît » un énantiomère et pas l'autre, comme une main droite ne rentre correctement que dans un gant droit.

\courssub{Représentation conventionnelle en perspective (représentation de Cram)}

Pour dessiner une molécule tridimensionnelle sur une feuille plane, on utilise la \cle{représentation de Cram}, qui code la profondeur par trois types de liaisons :

\begin{aretenir}[Conventions de la représentation de Cram]
\begin{itemize}
\item un \textbf{trait plein fin} : la liaison est \emph{dans le plan} de la feuille ;
\item un \textbf{triangle plein} (coin solide) : la liaison sort du plan, \emph{vers l'avant}, vers l'observateur ;
\item des \textbf{hachures} (triangle pointillé) : la liaison part \emph{vers l'arrière}, derrière le plan de la feuille.
\end{itemize}
Autour d'un carbone asymétrique, on place en général deux liaisons dans le plan, une en avant (triangle plein) et une en arrière (hachures).
\end{aretenir}

Avec chemfig, on écrit par exemple pour l'acide lactique : \verb|\chemfig{C(-[2]H)(-[6]CH_3)(<:[4]HO)(<OH)-COOH}| où \texttt{<} désigne le triangle plein et \texttt{<:} les hachures. La figure des deux énantiomères de l'acide lactique ci-dessus applique exactement ces conventions : observe que passer d'un énantiomère à l'autre revient à \emph{échanger} les positions avant/arrière de deux substituants ($\ce{-OH}$ et $\ce{-H}$).

\courssub{La projection de Fischer}

La représentation de Cram est parlante mais parfois lourde à dessiner, surtout pour les molécules à plusieurs $C^*$ comme les sucres. Le chimiste allemand Emil Fischer a proposé en 1891 une projection plane très pratique, à condition de respecter des règles strictes.

\begin{aretenir}[Règles de construction d'une projection de Fischer]
\begin{enumerate}
\item Le carbone asymétrique $C^*$ se trouve à l'\textbf{intersection} des deux traits (vertical et horizontal) ; on ne l'écrit généralement pas.
\item La \textbf{chaîne carbonée principale} est placée \textbf{verticalement}, et les liaisons verticales partent \emph{vers l'arrière} du plan (elles s'éloignent de l'observateur).
\item Les \textbf{substituants horizontaux} (à gauche et à droite) viennent \emph{vers l'avant}, vers l'observateur.
\item Le groupe le plus oxydé (souvent $\ce{-COOH}$, $\ce{-CHO}$) est placé par convention \emph{en haut} de la chaîne verticale.
\end{enumerate}
\end{aretenir}

\begin{popfigure}
\begin{center}
\begin{tikzpicture}[scale=0.9]
% Premier énantiomère (L-alanine)
\begin{scope}[xshift=-3.5cm]
\draw[thick, popblue] (0,1.2) -- (0,-1.2); % vertical
\draw[thick, popblue] (-1.2,0) -- (1.2,0); % horizontal
\node[popblue, above] at (0,1.2) {\ce{COOH}};
\node[popblue, below] at (0,-1.2) {\ce{CH3}};
\node[popblue, left] at (-1.2,0) {\ce{NH2}};
\node[popblue, right] at (1.2,0) {\ce{H}};
\node[popdark] at (0,-1.8) {\small (S)-Alanine (série L)};
\end{scope}
% Miroir
\draw[popmuted, thick, decorate, decoration={coil, aspect=0.4, segment length=4pt}] (0,-2.2) -- (0,1.8);
\node[popmuted, rotate=90] at (0.25,0) {\scriptsize miroir};
% Second énantiomère (D-alanine)
\begin{scope}[xshift=3.5cm]
\draw[thick, poporange] (0,1.2) -- (0,-1.2);
\draw[thick, poporange] (-1.2,0) -- (1.2,0);
\node[poporange, above] at (0,1.2) {\ce{COOH}};
\node[poporange, below] at (0,-1.2) {\ce{CH3}};
\node[poporange, left] at (-1.2,0) {\ce{H}};
\node[poporange, right] at (1.2,0) {\ce{NH2}};
\node[popdark] at (0,-1.8) {\small (R)-Alanine (série D)};
\end{scope}
\end{tikzpicture}
\legende{Projections de Fischer des deux énantiomères de l'alanine $\ce{CH3-CH(NH2)-COOH}$. Le $C^*$ est à l'intersection des traits ; $\ce{-COOH}$ en haut et $\ce{-CH3}$ en bas partent vers l'arrière ; $\ce{-NH2}$ et $\ce{-H}$, horizontaux, viennent vers l'avant. Les deux projections sont images miroir.}
\end{center}
\end{popfigure}

\begin{attention}[Erreur fréquente : la rotation d'une projection de Fischer]
Une projection de Fischer n'est \emph{pas} un dessin ordinaire : elle code une géométrie précise.
\begin{itemize}
\item Une \textbf{rotation de $180\degres$ dans le plan} de la feuille \textbf{conserve} la configuration : la molécule représentée est la même.
\item Une \textbf{rotation de $90\degres$} (ou de $270\degres$) \textbf{change la configuration} : on obtient l'\emph{énantiomère} ! En effet, une telle rotation ferait passer des liaisons « vers l'avant » en position verticale, ce qui violerait les conventions.
\item De même, \emph{sortir} la projection du plan (la retourner comme une crêpe) est interdit.
\item Échanger \emph{deux} substituants sur le $C^*$ inverse la configuration ; effectuer \emph{deux} échanges successifs (ou une permutation circulaire de trois groupes) la conserve.
\end{itemize}
\end{attention}

\courssub{Passage de Cram à Fischer et inversement}

\begin{methode}[Convertir une représentation de Cram en projection de Fischer]
Prenons l'exemple de l'acide lactique $\ce{CH3-CHOH-COOH}$.
\begin{enumerate}
\item \textbf{Orienter la molécule} de sorte que la chaîne principale ($\ce{HOOC-C^*-CH3}$) soit verticale et \emph{derrière} le plan, avec le groupe le plus oxydé ($\ce{-COOH}$) en haut.
\item \textbf{Faire pivoter} mentalement la molécule autour du $C^*$ jusqu'à ce que les deux substituants restants ($\ce{-OH}$ et $\ce{-H}$) pointent \emph{vers l'avant}, l'un à gauche, l'autre à droite.
\item \textbf{Aplatir} : remplacer le tétraèdre par une croix ; le $C^*$ est à l'intersection ; écrire les quatre groupes aux quatre extrémités.
\item \textbf{Vérifier} : en Fischer, les liaisons verticales vont vers l'arrière, les horizontales vers l'avant.
\end{enumerate}
Pour le passage inverse (Fischer $\rightarrow$ Cram), on lit la croix en se rappelant que l'horizontal vient vers soi (triangle plein) et que le vertical s'éloigne (hachures), puis on redessine le tétraèdre.
\end{methode}

\begin{exemplebox}[Application à l'alanine]
L'alanine $\ce{CH3-CH(NH2)-COOH}$ est l'acide aminé le plus simple après la glycine. Son $C^*$ porte $\ce{-H}$, $\ce{-NH2}$, $\ce{-CH3}$ et $\ce{-COOH}$. En plaçant $\ce{-COOH}$ en haut et $\ce{-CH3}$ en bas (chaîne verticale vers l'arrière), il reste deux possibilités : $\ce{-NH2}$ à gauche et $\ce{-H}$ à droite, ou l'inverse. Ce sont exactement les deux projections de Fischer de la figure ci-dessus : deux énantiomères. L'alanine des protéines naturelles correspond toujours à la même configuration (série L).
\end{exemplebox}

\begin{savaistu}[Le vivant a choisi son camp]
Les acides aminés constituant les protéines — à l'exception de la \emph{glycine} $\ce{H2N-CH2-COOH}$, dont le carbone porte deux hydrogènes identiques et n'est donc pas asymétrique — sont \textbf{tous chiraux}. Or, sur Terre, le vivant n'utilise pratiquement \emph{qu'une seule} des deux formes : les acides aminés L. Les sucres naturels, eux, appartiennent presque tous à la série D. Cette « préférence » du vivant, appelée \emph{homochiralité}, reste l'un des grands mystères de l'origine de la vie. Conséquence dramatique : un médicament et son énantiomère peuvent avoir des effets totalement différents. La thalidomide, prescrite dans les années 1950 contre les nausées des femmes enceintes, en est le triste exemple : un énantiomère était calmant, l'autre tératogène. Depuis, l'industrie pharmaceutique contrôle rigoureusement la stéréochimie de chaque principe actif.
\end{savaistu}

\begin{exercice}[À toi de jouer]
On considère les molécules suivantes : (a) $\ce{CH3-CHCl-CH3}$ ; (b) $\ce{CH3-CHOH-CH2-CH3}$ (butan-2-ol) ; (c) $\ce{CH3-CO-CH2-CH3}$ ; (d) $\ce{CH3-CH(NH2)-COOH}$.
\begin{enumerate}
\item Pour chaque molécule, indiquer si elle possède un carbone asymétrique.
\item En déduire le nombre maximal de stéréo-isomères dans chaque cas.
\item Représenter en projection de Fischer les deux énantiomères du butan-2-ol.
\end{enumerate}
\end{exercice}

\begin{corrige}[Exercice : carbones asymétriques et énantiomères]
\begin{enumerate}
\item (a) $\ce{CH3-CHCl-CH3}$ : le carbone central porte \emph{deux} groupes $\ce{-CH3}$ identiques : \textbf{pas de $C^*$}. (b) Le carbone 2 du butan-2-ol porte $\ce{-H}$, $\ce{-OH}$, $\ce{-CH3}$ et $\ce{-CH2-CH3}$, tous différents : \textbf{un $C^*$}. (c) La butanone possède un $\ce{C=O}$ : le carbone 2 est trigonal plan, et les autres portent des hydrogènes identiques : \textbf{pas de $C^*$}. (d) L'alanine : \textbf{un $C^*$} ($\ce{-H}$, $\ce{-NH2}$, $\ce{-CH3}$, $\ce{-COOH}$).
\item (a) $2^0 = 1$ (pas de stéréo-isomère de ce type) ; (b) $2^1 = 2$ énantiomères ; (c) $1$ ; (d) $2$ énantiomères.
\item Fischer du butan-2-ol : chaîne verticale $\ce{-CH3}$ en haut (C1) et $\ce{-CH2CH3}$ en bas (C3-C4) ; premier énantiomère : $\ce{-OH}$ à gauche, $\ce{-H}$ à droite ; second énantiomère : $\ce{-H}$ à gauche, $\ce{-OH}$ à droite. Les deux croix sont images miroir.
\end{enumerate}
\end{corrige}

En résumé : un carbone tétraédrique portant quatre groupes différents est un \cle{carbone asymétrique} ; il rend la molécule \cle{chirale}, c'est-à-dire non superposable à son image miroir ; les deux formes obtenues sont des \cle{énantiomères}, aux propriétés physiques et chimiques identiques, sauf vis-à-vis de la lumière polarisée et des réactifs chiraux. C'est précisément cette action sur la lumière polarisée — l'\emph{activité optique} — et la notion de \emph{diastéréoisomères} que nous allons maintenant étudier dans la section suivante.

\coursec{Activité optique et diastéréoisomères}

Dans la section précédente, nous avons découvert que deux énantiomères sont comme deux gants : identiques en tout point, mais non superposables. Une question se pose alors immédiatement : \emph{comment les distinguer en pratique ?} Leur masse molaire, leur température d'ébullition, leur densité sont rigoureusement identiques\ldots{} Il existe pourtant un test physique remarquable, découvert au XIX\textsuperscript{e} siècle par Louis Pasteur et Jean-Baptiste Biot : le comportement de ces molécules vis-à-vis de la \cle{lumière polarisée}. C'est cette propriété, l'\cle{activité optique}, qui va nous permettre de « voir » la chiralité. Nous verrons ensuite que les stéréo-isomères ne sont pas tous des énantiomères : certains, appelés \cle{diastéréoisomères}, ont des propriétés différentes, ce qui change tout en chimie.

\courssub{La lumière polarisée : un rappel de physique}

\courssub{De la lumière naturelle à la lumière polarisée}

La lumière est une onde électromagnétique : elle comporte un champ électrique qui vibre perpendiculairement à la direction de propagation. Dans la \cle{lumière naturelle} (soleil, lampe), cette vibration se fait dans \textbf{tous les plans} contenant la direction de propagation : on dit que la lumière n'est pas polarisée.

Lorsqu'on fait traverser cette lumière à un cristal spécial appelé \cle{polariseur} (prisme de Nicol, filtre polaroïd), seule la composante vibrant dans \textbf{un seul plan} est transmise. On obtient une \cle{lumière polarisée rectilignement} : le champ électrique vibre dans un plan unique, appelé \cle{plan de polarisation}.

\begin{definition}[Lumière polarisée]
Une lumière est dite \cle{polarisée rectilignement} lorsque son champ électrique vibre dans un seul plan, le \cle{plan de polarisation}. On l'obtient en faisant passer la lumière naturelle à travers un polariseur.
\end{definition}

\begin{savaistu}[Les lunettes de soleil polarisées]
Les verres « polarisants » des lunettes de soleil vendues à Douala ou Yaoundé sont des polariseurs : ils éliminent la lumière réfléchie horizontalement par la chaussée ou par l'eau, ce qui supprime l'éblouissement. Le même principe physique sert en chimie à analyser les molécules chirales !
\end{savaistu}

\courssub{Le pouvoir rotatoire : la signature des molécules chirales}

\textbf{Intuition.} Imagine une hélice : si tu la fais avancer dans un écrou, elle tourne. Une molécule chirale, qui a une structure « en hélice » dans l'espace, va de même « visser » le plan de polarisation de la lumière qui la traverse.

\begin{propriete}[Pouvoir rotatoire]
Une substance \cle{optiquement active} fait tourner le plan de polarisation de la lumière polarisée qui la traverse d'un angle $\alpha$, appelé \cle{pouvoir rotatoire} (ou angle de déviation). Toutes les substances chirales pures (un seul énantiomère) sont optiquement actives ; les substances achirales ne le sont pas.
\end{propriete}

L'observateur qui regarde la lumière sortir de la solution distingue deux cas :

\begin{itemize}
\item le plan de polarisation a tourné vers la \textbf{droite} (sens des aiguilles d'une montre) : la substance est \cle{dextrogyre}, notée $(+)$ ;
\item le plan de polarisation a tourné vers la \textbf{gauche} : la substance est \cle{lévogyre}, notée $(-)$.
\end{itemize}

\begin{propriete}[Propriété fondamentale des énantiomères]
Deux énantiomères ont des pouvoirs rotatoires de \textbf{même valeur absolue mais de signes opposés} : si l'un est dextrogyre $(+\alpha)$, l'autre est lévogyre $(-\alpha)$. C'est la \textbf{seule} propriété physique qui les distingue (hors interaction avec d'autres objets chiraux).
\end{propriete}

\begin{attention}[Signe et configuration ne coïncident pas forcément]
Le signe $(+)$ ou $(-)$ est une propriété \emph{expérimentale} : il ne se déduit pas de la configuration absolue ($R$ ou $S$) de la molécule. Un composé ($R$) peut très bien être lévogyre. Ne confonds jamais « configuration » (géométrie de la molécule) et « sens de déviation » (mesure au polarimètre).
\end{attention}

\courssub{Le polarimètre : mesurer l'angle de déviation}

Le \cle{polarimètre} est l'appareil qui mesure le pouvoir rotatoire. Son principe est simple : on compare la direction du plan de polarisation avant et après la traversée de la solution étudiée.

\begin{popfigure}
\begin{center}
\begin{tikzpicture}[scale=0.92, every node/.style={font=\small}]
% Source
\draw[fill=popgold!60, draw=popdark] (0,0) circle (0.45);
\node[below=2pt] at (0,-0.45) {Source (Na)};
\foreach \a in {30,90,150,210,270,330} \draw[popgold, -{Stealth}] (\a:0.5) -- (\a:0.85);
% Lumière naturelle
\draw[popmuted, -{Stealth}, thick] (0.6,0) -- (1.5,0);
\node[popmuted, above] at (1.05,0.05) {\scriptsize lumière naturelle};
% Polariseur
\draw[fill=popblueL, draw=popblue, thick] (1.6,-0.9) rectangle (1.9,0.9);
\node[below=2pt, align=center] at (1.75,-0.9) {Polariseur};
\draw[popblue, -{Stealth}, thick] (2.0,0) -- (2.9,0);
% Plan de polarisation vertical
\draw[popblue, thick, dashed] (3.6,-0.8) -- (3.6,0.8);
\node[popblue, above, font=\scriptsize] at (3.6,0.8) {plan de polarisation};
% Tube
\draw[fill=popgreenL, draw=popgreen!70!black, thick, rounded corners=3pt] (4.2,-0.55) rectangle (7.4,0.55);
\node at (5.8,0) {\scriptsize solution de substance active};
\node[below=2pt] at (5.8,-0.55) {Tube de longueur $\ell$};
\draw[popgreen!70!black, -{Stealth}, thick] (7.5,0) -- (8.3,0);
% Plan dévié
\draw[poppink, thick, dashed] (8.9,-0.7) -- (9.35,0.7);
\draw[popblue, thick, dashed] (9.12,-0.8) -- (9.12,0.8);
\draw[poppink, -{Stealth}] (9.12,0.5) arc (90:58:0.5);
\node[poppink, right, font=\scriptsize] at (9.45,0.55) {$\alpha$};
% Analyseur
\draw[fill=popblueL, draw=popblue, thick] (10.1,-0.9) rectangle (10.4,0.9);
\node[below=2pt] at (10.25,-0.9) {Analyseur};
% Oeil
\draw[fill=popcream, draw=popdark] (11.3,0) circle (0.35);
\fill[popdark] (11.3,0) circle (0.12);
\node[below=2pt] at (11.3,-0.35) {Observateur};
\end{tikzpicture}
\end{center}
\legende{Schéma de principe du polarimètre : la lumière polarisée traverse la solution ; le plan de polarisation tourne d'un angle $\alpha$, mesuré en tournant l'analyseur jusqu'au rétablissement de l'intensité maximale.}
\end{popfigure}

\textbf{Protocole de mesure.} On règle d'abord l'analyseur (tube vide ou rempli de solvant pur) pour obtenir l'intensité lumineuse maximale. On place ensuite la solution : l'intensité diminue car le plan a tourné. On fait pivoter l'analyseur jusqu'à retrouver le maximum : l'angle de rotation lu est $\alpha$.

\begin{propriete}[Loi de Biot]
Pour une solution de substance active, l'angle de déviation est proportionnel à la concentration et à la longueur du tube :
\[ \alpha = [\alpha] \times \ell \times C \]
où $[\alpha]$ est la \cle{rotation spécifique} (caractéristique de la substance, à température et longueur d'onde données), $\ell$ la longueur du tube en décimètres et $C$ la concentration massique en \SI{}{g.mL^{-1}}.
\end{propriete}

\begin{exoresolu}[Mesure du pouvoir rotatoire d'une solution de glucose]
Une solution de glucose de concentration massique $C = \SI{0,20}{g.mL^{-1}}$ est placée dans un tube polarimétrique de longueur $\ell = \SI{2,0}{dm}$. La rotation spécifique du glucose est $[\alpha] = +52,5$ degrés. Calculer l'angle de déviation $\alpha$ mesuré et préciser si la solution est dextrogyre ou lévogyre.
\tcblower
\textbf{Solution.} On applique la loi de Biot :
\[ \alpha = [\alpha] \times \ell \times C = 52,5 \times 2,0 \times 0,20 = +21 \text{ degrés}. \]
Le plan de polarisation tourne de $21$ degrés vers la droite : la solution est \textbf{dextrogyre}. C'est d'ailleurs pour cela que le glucose est aussi appelé \emph{dextrose}.
\end{exoresolu}

\courssub{Le mélange racémique : quand les effets s'annulent}

Que se passe-t-il si l'on dissout dans le tube \emph{autant} d'énantiomère $(+)$ que d'énantiomère $(-)$ ? Chaque molécule $(+)$ fait tourner le plan vers la droite, chaque molécule $(-)$ le fait tourner du même angle vers la gauche. Les deux effets se compensent exactement : le plan de polarisation ne tourne pas.

\begin{definition}[Mélange racémique]
Un \cle{mélange racémique} (ou \cle{racémate}) est un mélange \textbf{équimolaire} (50 \% / 50 \%) de deux énantiomères. Il est \textbf{optiquement inactif} : son pouvoir rotatoire global est nul, car les déviations opposées se compensent. On le note $(+/-)$ ou $[\alpha] = 0$.
\end{definition}

\begin{exemplebox}[L'acide lactique du lait caillé]
L'acide lactique \ce{CH3-CHOH-COOH} possède un carbone asymétrique. L'acide lactique produit par les muscles est dextrogyre ; celui obtenu par fermentation du lait (lait caillé, très apprécié au Nord-Cameroun) est un mélange racémique $(+/-)$, optiquement inactif, car la fermentation produit les deux énantiomères en quantités égales.
\end{exemplebox}

\begin{attention}[Inactif ne veut pas dire achiral]
Un mélange racémique est optiquement inactif, mais il \emph{contient} des molécules chirales ! L'inactivité résulte d'une compensation statistique, non de l'absence de chiralité. Ne confonds pas un racémate avec une substance achirale comme l'éthanol.
\end{attention}

\courssub{Les diastéréoisomères : des stéréo-isomères différents}

\courssub{Définition et reconnaissance}

Jusqu'ici, nous n'avons rencontré que des paires d'énantiomères. Mais dès qu'une molécule possède \textbf{plusieurs} carbones asymétriques, la situation s'enrichit : on peut trouver des stéréo-isomères qui ne sont \emph{pas} images l'un de l'autre dans un miroir.

\begin{definition}[Diastéréoisomères]
Deux stéréo-isomères sont des \cle{diastéréoisomères} s'ils \textbf{ne sont pas} images l'un de l'autre dans un miroir. Contrairement aux énantiomères, les diastéréoisomères ont des \textbf{propriétés physiques différentes} (températures de fusion et d'ébullition, solubilité, densité, pouvoir rotatoire) et des \textbf{propriétés chimiques différentes} (vitesses de réaction, produits formés).
\end{definition}

\begin{methode}[Distinguer énantiomères et diastéréoisomères]
Pour deux molécules possédant plusieurs carbones asymétriques $C^*$, compare les configurations \textbf{carbone par carbone} :
\begin{enumerate}
\item Repère tous les $C^*$ de chaque molécule.
\item Pour chaque $C^*$, vérifie si la configuration est \emph{identique} ou \emph{inversée} entre les deux molécules.
\item Conclus :
\begin{itemize}
\item si \textbf{tous} les $C^*$ sont inversés : les deux molécules sont des \textbf{énantiomères} (images miroir non superposables) ;
\item si \textbf{certains} $C^*$ seulement sont inversés (et les autres identiques) : ce sont des \textbf{diastéréoisomères} ;
\item si \textbf{aucun} n'est inversé : c'est la même molécule.
\end{itemize}
\end{enumerate}
\end{methode}

\courssub{Exemple type : l'acide tartrique}

L'\cle{acide tartrique}, de formule \ce{HOOC-CHOH-CHOH-COOH}, est l'acide du raisin et du vin (il se dépose en cristaux, le « tartre », au fond des cuves). Il possède \textbf{deux} carbones asymétriques, $C_2$ et $C_3$. On pourrait s'attendre à $2^2 = 4$ stéréo-isomères\ldots{} mais il n'en existe que \textbf{trois} !

\begin{popfigure}
\begin{center}
\setchemfig{atom sep=1.6em}
\begin{tikzpicture}[node distance=0.4cm]
% (2R,3R)
\node (A) {\chemfig{COOH-[2](-[1]H-[3]HO)-[6](-[5]H-[7]OH)-[2]COOH}};
\node[below=of A, font=\small\bfseries, popblue] {(2R,3R) $(+)$};
% (2S,3S)
\node[right=2.6cm of A] (B) {\chemfig{COOH-[2](-[1]HO-[3]H)-[6](-[5]OH-[7]H)-[2]COOH}};
\node[below=of B, font=\small\bfseries, popblue] {(2S,3S) $(-)$};
% méso
\node[right=2.6cm of B] (C) {\chemfig{COOH-[2](-[1]H-[3]HO)-[6](-[5]OH-[7]H)-[2]COOH}};
\node[below=of C, font=\small\bfseries, poppink] {méso (2R,3S), inactif};
% relations
\draw[poppurple, thick, {Stealth}-{Stealth}] ($(A.south east)+(0.1,-0.9)$) -- ($(B.south west)+(-0.1,-0.9)$) node[midway, below, font=\scriptsize, poppurple] {énantiomères};
\draw[popgreen!70!black, thick, {Stealth}-{Stealth}] ($(B.south east)+(0.1,-0.9)$) -- ($(C.south west)+(-0.1,-0.9)$) node[midway, below, font=\scriptsize, popgreen!70!black] {diastéréoisomères};
\draw[popgreen!70!black, thick, {Stealth}-{Stealth}] ($(A.south)+(0,-1.15)$) to[out=-25,in=-155] node[midway, below, font=\scriptsize, popgreen!70!black] {diastéréoisomères} ($(C.south)+(0,-1.15)$);
\end{tikzpicture}
\end{center}
\legende{Les trois stéréo-isomères de l'acide tartrique en représentation de Fischer (chaîne verticale, \ce{COOH} en haut et en bas). Les formes $(+)$ et $(-)$ sont énantiomères ; la forme \emph{méso} possède un plan de symétrie interne (entre $C_2$ et $C_3$) : elle est achirale et optiquement inactive.}
\end{popfigure}

Analysons ces trois formes :

\begin{itemize}
\item \textbf{Forme (2R,3R)} : dextrogyre $(+)$, c'est l'acide tartrique naturel du raisin ;
\item \textbf{Forme (2S,3S)} : lévogyre $(-)$, image miroir de la précédente : les deux $C^*$ sont inversés, ce sont bien des \textbf{énantiomères} ;
\item \textbf{Forme méso (2R,3S)} : les deux moitiés de la molécule sont images l'une de l'autre \emph{à l'intérieur} de la molécule. Elle possède un \cle{plan de symétrie interne} (horizontal, entre $C_2$ et $C_3$) : elle est donc \textbf{superposable à son image miroir}, c'est-à-dire \textbf{achirale}, malgré ses deux $C^*$ ! Elle est optiquement inactive.
\end{itemize}

La forme méso est diastéréoisomère de chacune des deux formes actives : comparée à (2R,3R), un seul $C^*$ ($C_3$) est inversé.

\begin{attention}[Piège classique : $2^n$ n'est pas toujours la bonne réponse]
On croit souvent qu'une molécule à $n$ carbones asymétriques possède toujours $2^n$ stéréo-isomères. C'est \textbf{faux} dès que la molécule présente une symétrie interne ! Pour l'acide tartrique ($n = 2$), il n'y a que \textbf{3} stéréo-isomères et non 4, car la forme (2R,3S) et la forme (2S,3R) sont \emph{une seule et même molécule} : la forme \cle{mésomère}, achirale grâce à son plan de symétrie. Avant d'appliquer $2^n$, cherche toujours un éventuel plan de symétrie.
\end{attention}

\courssub{Pourquoi les diastéréoisomères sont-ils si importants ?}

Contrairement aux énantiomères (séparables seulement par des méthodes chirales délicates), les diastéréoisomères ont des propriétés physiques \emph{différentes} : on peut donc les \textbf{séparer par les techniques classiques} (distillation, cristallisation, chromatographie). C'est d'ailleurs la stratégie historique de Pasteur pour séparer les énantiomères : on les transforme d'abord en diastéréoisomères (en les faisant réagir avec un réactif chiral), on les sépare, puis on régénère les énantiomères purs.

\begin{center}
\rowcolors{1}{popcream}{white}
\begin{tabularx}{\linewidth}{|l|X|X|}
\hline
\rowcolor{popblueL} \textbf{Critère} & \textbf{Énantiomères} & \textbf{Diastéréoisomères} \\
\hline
Image miroir ? & Oui, non superposables & Non \\
Propriétés physiques & Identiques (sauf $\alpha$) & Différentes \\
Propriétés chimiques & Identiques (réactif achiral) & Différentes \\
Pouvoir rotatoire & Opposés ($\pm\alpha$) & Quelconques \\
Séparation & Méthodes chirales & Techniques classiques \\
\hline
\end{tabularx}
\end{center}

\begin{savaistu}[Le sucre inverti, star des pâtissiers]
Le saccharose (le sucre de canne ou de betterave) est \textbf{dextrogyre} : $[\alpha] = +66,5$ degrés. Son hydrolyse, catalysée par un acide ou par l'enzyme invertase, coupe la molécule en deux :
\[ \ce{C12H22O11 + H2O ->[H^+\ \text{ou invertase}] C6H12O6\ (glucose) + C6H12O6\ (fructose)} \]
Or le mélange glucose $+$ fructose obtenu est \textbf{lévogyre} ($[\alpha] \approx -20$ degrés), car le fructose, très lévogyre ($-92$ degrés), l'emporte sur le glucose dextrogyre ($+52,5$ degrés). Le signe du pouvoir rotatoire s'est \emph{inversé} au cours de la réaction : c'est pourquoi ce produit est appelé \cle{sucre inverti}. Les pâtissiers l'adorent car il retient l'humidité et garde les gâteaux moelleux. Le polarimètre permet ainsi de suivre l'avancement de l'hydrolyse : c'est la \emph{saccharimétrie}, utilisée dans les sucreries comme la SOSUCAM !
\end{savaistu}

\begin{exercice}[Énantiomères ou diastéréoisomères ?]
On considère les molécules A (2R,3R), B (2S,3S), C (2R,3S) et D (2S,3R) d'un composé à deux $C^*$ sans symétrie interne. Préciser la relation (énantiomères, diastéréoisomères ou même molécule) entre : A et B ; A et C ; B et D ; C et D.
\end{exercice}

\begin{corrige}[Exercice 1]
On applique la méthode de comparaison carbone par carbone.
\begin{itemize}
\item \textbf{A et B} : les deux $C^*$ sont inversés (R$\to$S et R$\to$S) : \textbf{énantiomères}.
\item \textbf{A et C} : $C_2$ identique (R,R), $C_3$ inversé (R$\to$S) : \textbf{diastéréoisomères}.
\item \textbf{B et D} : $C_2$ identique (S,S), $C_3$ inversé (S$\to$R) : \textbf{diastéréoisomères}.
\item \textbf{C et D} : les deux $C^*$ sont inversés : \textbf{énantiomères}.
\end{itemize}
Remarque : sans symétrie interne, les quatre formes sont bien distinctes ($2^2 = 4$), contrairement au cas de l'acide tartrique.
\end{corrige}

\begin{aretenir}[L'essentiel de la section]
\begin{itemize}
\item Une substance chirale pure est \cle{optiquement active} : elle fait tourner le plan de polarisation de la lumière polarisée d'un angle $\alpha$, mesuré au \cle{polarimètre} (loi de Biot : $\alpha = [\alpha]\,\ell\,C$).
\item Deux énantiomères ont des pouvoirs rotatoires \textbf{opposés} : l'un est dextrogyre $(+)$, l'autre lévogyre $(-)$.
\item Un \cle{mélange racémique} (50/50 de deux énantiomères) est optiquement \textbf{inactif} par compensation.
\item Les \cle{diastéréoisomères} sont des stéréo-isomères non images miroir : certains $C^*$ inversés seulement ; leurs propriétés physiques et chimiques sont \textbf{différentes}.
\item Une molécule à $n$ $C^*$ possède au plus $2^n$ stéréo-isomères ; une \cle{forme méso} (plan de symétrie interne) est achirale et réduit ce nombre : l'acide tartrique n'a que 3 stéréo-isomères.
\end{itemize}
\end{aretenir}

\coursec{Isomérie Z/E des alcènes}

Dans la section précédente, nous avons découvert les \cle{diastéréoisomères} : des stéréo-isomères qui ne sont \emph{pas} images l'un de l'autre dans un miroir. Nous les avions rencontrés chez des molécules possédant plusieurs carbones asymétriques. Mais il existe une autre famille de diastéréoisomères, sans aucun carbone asymétrique, qui joue un rôle considérable en chimie et en biologie : ce sont les isomères \cle{Z} et \cle{E} des \cle{alcènes}. Tu en connais déjà sans le savoir : c'est grâce à une isomérisation Z/E que ton œil perçoit la lumière en ce moment même !

\courssub{Origine : la double liaison C=C bloque la rotation}

\textbf{Intuition d'abord.} Prends deux cure-dents plantés dans une pomme : tu peux faire tourner librement l'un par rapport à l'autre autour de l'axe. C'est ce qui se passe autour d'une liaison simple C--C (liaison $\sigma$) : la rotation est \emph{libre}. Maintenant, relie les deux cure-dents par un élastique tendu en plus : impossible de tourner sans déformer ou casser l'ensemble. C'est exactement la situation d'une double liaison \ce{C=C}.

\begin{propriete}[Structure de la double liaison C=C]
Une double liaison \ce{C=C} est constituée de :
\begin{itemize}
\item une \cle{liaison $\sigma$}, formée par recouvrement axial des orbitales hybrides $sp^2$ ; elle autoriserait la rotation ;
\item une \cle{liaison $\pi$}, formée par recouvrement \emph{latéral} de deux orbitales $p$ parallèles, situées au-dessus et au-dessous du plan de la molécule.
\end{itemize}
Faire tourner un carbone par rapport à l'autre imposerait de \emph{rompre le recouvrement latéral} des orbitales $p$, donc de casser la liaison $\pi$. L'énergie nécessaire (environ \SI{270}{kJ.mol^{-1}}) est bien trop élevée à température ordinaire : la rotation autour de \ce{C=C} est \cle{bloquée}.
\end{propriete}

\begin{popfigure}
\begin{center}
\begin{tikzpicture}[scale=1.0, >=Stealth]
% --- Liaison simple : rotation libre ---
\begin{scope}[xshift=-4.2cm]
\draw[line width=2.2pt, popblue] (0,0) -- (2.2,0);
\fill[popblue] (0,0) circle (0.16) node[below=4pt, popdark]{\footnotesize C};
\fill[popblue] (2.2,0) circle (0.16) node[below=4pt, popdark]{\footnotesize C};
\draw[->, popgreen, line width=1.1pt] (2.2,0.75) arc[start angle=90, end angle=-30, radius=0.75];
\node[popgreen, font=\footnotesize] at (3.15,0.55) {rotation libre};
\node[popdark, font=\small\bfseries] at (1.1,-1.15) {Liaison simple $\sigma$};
\end{scope}
% --- Double liaison : rotation bloquée ---
\begin{scope}[xshift=2.2cm]
% lobes pi
\fill[poporange!35] (0,0.55) ellipse (0.42 and 0.30);
\fill[poporange!35] (2.2,0.55) ellipse (0.42 and 0.30);
\fill[poporange!35] (0,-0.55) ellipse (0.42 and 0.30);
\fill[poporange!35] (2.2,-0.55) ellipse (0.42 and 0.30);
\draw[poporange, dashed, line width=0.9pt] (0.35,0.55) -- (1.85,0.55);
\draw[poporange, dashed, line width=0.9pt] (0.35,-0.55) -- (1.85,-0.55);
% liaison sigma
\draw[line width=2.2pt, popblue] (0,0) -- (2.2,0);
\fill[popblue] (0,0) circle (0.16);
\fill[popblue] (2.2,0) circle (0.16);
% croix de blocage
\draw[poppink, line width=1.6pt] (3.05,0.35) -- (3.55,0.85);
\draw[poppink, line width=1.6pt] (3.05,0.85) -- (3.55,0.35);
\node[poppink, font=\footnotesize, align=center] at (3.3,-0.35) {rotation\\bloquée};
\node[poporange, font=\footnotesize] at (1.1,1.15) {liaison $\pi$ (recouvrement latéral)};
\node[popdark, font=\small\bfseries] at (1.1,-1.15) {Double liaison $\sigma + \pi$};
\end{scope}
\end{tikzpicture}
\legende{Autour d'une liaison simple $\sigma$, la rotation est libre. Autour d'une double liaison \ce{C=C}, la liaison $\pi$ (recouvrement latéral des orbitales $p$, en orange) bloque toute rotation : les positions relatives des substituants sont figées.}
\end{center}
\end{popfigure}

\textbf{Conséquence capitale.} Puisque la rotation est bloquée, deux substituants placés du même côté de la double liaison y \emph{restent} : la géométrie de la molécule est figée. Si l'on peut écrire deux géométries différentes et non superposables, on obtient deux \emph{molécules distinctes} : ce sont les isomères Z et E.

\begin{definition}[Isomérie Z/E]
L'\cle{isomérie Z/E} (ou isomérie géométrique) est une \cle{stéréo-isomérie} due à l'absence de rotation libre autour d'une double liaison \ce{C=C}. Deux isomères Z/E ont la \emph{même} formule brute et la \emph{même} chaîne de liaisons, mais diffèrent par la \emph{position dans l'espace} des groupes fixés sur les deux carbones de la double liaison.
\end{definition}

\courssub{Condition d'existence de l'isomérie Z/E}

Tous les alcènes ne possèdent pas d'isomères Z/E. Il faut examiner \emph{chaque} carbone de la double liaison.

\begin{propriete}[Condition d'existence]
Un alcène présente l'isomérie Z/E si et seulement si \textbf{chacun} des deux carbones de la double liaison porte \textbf{deux groupes différents}. Le motif minimal est donc :
\[ \text{a}\ce{C=C}\text{b} \quad\text{avec, sur chaque carbone, deux substituants distincts (motif abC=Cab au minimum).} \]
Si l'un des deux carbones porte deux groupes \emph{identiques}, échanger ces groupes ne change rien : il n'existe qu'\emph{une seule} molécule, pas d'isomérie Z/E.
\end{propriete}

\begin{exemplebox}[Alcènes sans isomérie Z/E]
\begin{itemize}
\item \textbf{Propène} \ce{CH2=CH-CH3} : le premier carbone porte \emph{deux} atomes d'hydrogène identiques $\Rightarrow$ pas d'isomérie Z/E.
\item \textbf{But-1-ène} \ce{CH2=CH-CH2-CH3} : le carbone 1 porte deux H identiques $\Rightarrow$ pas d'isomérie Z/E.
\item \textbf{2-méthylpropène} \ce{CH2=C(CH3)2} : un carbone porte deux H, l'autre porte deux \ce{CH3} $\Rightarrow$ pas d'isomérie Z/E.
\item \textbf{But-2-ène} \ce{CH3-CH=CH-CH3} : chaque carbone de \ce{C=C} porte un H et un \ce{CH3}, tous différents $\Rightarrow$ \textbf{deux isomères Z/E existent}.
\end{itemize}
\end{exemplebox}

\courssub{Configurations Z et E : la règle CIP}

Pour nommer sans ambiguïté les isomères, on utilise la règle de \cle{Cahn-Ingold-Prelog (CIP)}, fondée sur le numéro atomique.

\begin{definition}[Règle de priorité CIP et configurations Z/E]
Sur chaque carbone de la double liaison, on compare les numéros atomiques des atomes \emph{directement liés} au carbone :
\begin{itemize}
\item L'atome de numéro atomique le plus élevé a la \textbf{priorité la plus forte}.
\item En cas d'égalité, on compare les atomes suivants (liste dans l'ordre décroissant de Z).
\end{itemize}
Une fois les groupes prioritaires identifiés sur les deux carbones :
\begin{itemize}
\item configuration \cle{Z} (allemand \emph{zusammen}, « ensemble ») : les deux groupes prioritaires sont situés \textbf{du même côté} de la double liaison ;
\item configuration \cle{E} (allemand \emph{entgegen}, « opposé ») : les deux groupes prioritaires sont situés \textbf{de part et d'autre} de la double liaison.
\end{itemize}
\end{definition}

\begin{popfigure}
\begin{center}
\setchemfig{atom sep=2.1em, bond offset=1.5pt}
\chemfig{H_3C-[2](-[1]CH_3)=[2](-[3]H)-[2]CH_3}
\hspace{2.5cm}
\chemfig{H_3C-[2](-[3]CH_3)=[2](-[1]H)-[2]CH_3}
\legende{À gauche : \textbf{(Z)-but-2-ène} — les deux groupes \ce{CH3} (prioritaires sur H) sont du même côté. À droite : \textbf{(E)-but-2-ène} — les deux groupes \ce{CH3} sont de part et d'autre. Ces deux molécules ne sont pas superposables et ne s'interconvertissent pas à température ordinaire.}
\end{center}
\end{popfigure}

\begin{methode}[Attribuer la configuration Z ou E]
\begin{enumerate}
\item \textbf{Repérer} la double liaison \ce{C=C} dans la molécule.
\item \textbf{Vérifier la condition d'existence} : chaque carbone de \ce{C=C} porte-t-il deux groupes \emph{différents} ? Si non, il n'y a pas d'isomérie Z/E : on s'arrête là.
\item \textbf{Déterminer les priorités CIP} sur chaque carbone : comparer les numéros atomiques des atomes directement liés au carbone de la double liaison (ex. : Cl > O > N > C > H). Le groupe de plus haut numéro atomique est prioritaire.
\item \textbf{Observer la position relative} des deux groupes prioritaires : même côté $\Rightarrow$ \cle{Z} ; côtés opposés $\Rightarrow$ \cle{E}.
\item \textbf{Nommer} la molécule en préfixant (Z)- ou (E)- au nom de l'alcène : par exemple \textbf{(E)-but-2-ène}.
\end{enumerate}
\end{methode}

\courssub{Z et E sont des diastéréoisomères}

\begin{propriete}[Nature de la relation Z/E]
Deux isomères Z et E sont des \cle{diastéréoisomères} : ce sont des stéréo-isomères qui \textbf{ne sont pas} images l'un de l'autre dans un miroir. Par conséquent :
\begin{itemize}
\item ils ont des \textbf{propriétés physiques différentes} : températures de fusion et d'ébullition, densité, solubilité, indice de réfraction ;
\item ils ont souvent des \textbf{propriétés chimiques et biologiques différentes} (réactivité stéréospecifique, reconnaissance enzymatique) ;
\item ils sont \textbf{séparables} par les techniques classiques (distillation fractionnée, cristallisation, chromatographie).
\end{itemize}
\end{propriete}

\begin{exemplebox}[Acide maléique et acide fumarique : le même motif, deux molécules]
Les acides but-2-ènedioïques de formule \ce{HOOC-CH=CH-COOH} illustrent spectaculairement la différence Z/E :
\begin{center}
\begin{tabularx}{\linewidth}{|l|X|X|}
\hline
\rowcolor{popblueL}
\textbf{Propriété} & \textbf{Acide maléique (Z)} & \textbf{Acide fumarique (E)} \\
\hline
Position des \ce{COOH} & même côté de \ce{C=C} & côtés opposés \\
\hline
Température de fusion & \SI{131}{\celsius} & \SI{287}{\celsius} \\
\hline
Solubilité dans l'eau (\SI{25}{\celsius}) & élevée ($\approx$ \SI{790}{g.L^{-1}}) & faible ($\approx$ \SI{7}{g.L^{-1}}) \\
\hline
Présence dans le vivant & quasi absente & intermédiaire du cycle de Krebs \\
\hline
\end{tabularx}
\end{center}
L'écart de fusion de plus de \SI{150}{\celsius} s'explique : la molécule (E), plus symétrique et plus plane, s'empile mieux dans le cristal, ce qui renforce la cohésion du solide. L'acide maléique (Z) peut former un cycle intramoléculaire par liaison hydrogène, ce qui explique sa plus grande solubilité et sa réactivité particulière (formation d'anhydride maléique par chauffage).
\end{exemplebox}

\begin{attention}[Pièges fréquents à éviter]
\begin{itemize}
\item \textbf{Confondre Z/E et énantiomères} : les isomères Z/E \textbf{ne sont pas images dans un miroir} (ce sont des diastéréoisomères). Ils n'ont généralement \textbf{aucun carbone asymétrique} et sont \textbf{inactifs} sur la lumière polarisée (sauf cas rares de chiralité axiale). Leurs propriétés physiques sont \emph{différentes}, alors que deux énantiomères ont des propriétés physiques \emph{identiques} (sauf le pouvoir rotatoire).
\item \textbf{Oublier la condition d'existence} : attribuer Z/E à un alcène dont un carbone porte deux groupes identiques (propène, 2-méthylpropène, but-1-ène) est une erreur grave. Vérifier \emph{systématiquement} les deux carbones.
\item \textbf{Mauvaise priorité CIP} : ne pas confondre « groupe le plus gros » et « groupe prioritaire ». C'est le \textbf{numéro atomique du premier atome} qui décide (ex. : \ce{-CH2Cl} > \ce{-CH3} car Cl > H ; \ce{-COOH} > \ce{-CH2OH} car O > C).
\end{itemize}
\end{attention}

\begin{exoresolu}[Isomérie Z/E du pent-2-ène]
\textbf{Énoncé.} On considère le pent-2-ène de formule semi-développée \ce{CH3-CH=CH-CH2-CH3}.
\begin{enumerate}
\item Montrer que cet alcène présente l'isomérie Z/E.
\item Écrire les représentations des deux stéréo-isomères et les nommer.
\item Quelle relation d'isomérie les lie ? Sont-ils séparables par distillation ?
\end{enumerate}
\tcblower
\textbf{Solution détaillée.}
\begin{enumerate}
\item \textbf{Condition d'existence.} Le carbone 2 de la double liaison porte H (Z=1) et \ce{CH3} (C, Z=6) : groupes différents. Le carbone 3 porte H (Z=1) et \ce{CH2-CH3} (C, Z=6) : groupes différents. Chaque carbone de \ce{C=C} porte donc deux groupes différents : \textbf{l'isomérie Z/E existe}.
\item \textbf{Priorités CIP.} Sur C2 : \ce{CH3} (C) > H. Sur C3 : \ce{CH2CH3} (C) > H. Les groupes prioritaires sont les groupes alkyles.
\item \textbf{Représentations et noms.}
\begin{itemize}
\item \textbf{(Z)-pent-2-ène} : les deux groupes alkyles (\ce{CH3} et \ce{C2H5}) sont du \emph{même côté} de la double liaison.
\item \textbf{(E)-pent-2-ène} : les deux groupes alkyles sont \emph{de part et d'autre} de la double liaison.
\end{itemize}
\item \textbf{Relation d'isomérie.} Ces deux stéréo-isomères ne sont pas images l'un de l'autre dans un miroir : ce sont des \textbf{diastéréoisomères}. Leurs propriétés physiques (températures d'ébullition voisines mais distinctes : \SI{36}{\celsius} pour Z, \SI{37}{\celsius} pour E ; densités différentes) diffèrent : ils sont donc \textbf{séparables} par distillation fractionnée ou chromatographie.
\end{enumerate}
\end{exoresolu}

\begin{experience}[Interconversion Z/E : une réaction qui demande de l'énergie]
\textbf{Protocole.} On chauffe l'acide maléique (Z) en solution aqueuse acide vers \SI{100}{\celsius}, ou on l'irradie en présence d'une trace de brome (catalyse par radicalaires).
\textbf{Observation.} On obtient progressivement l'acide fumarique (E), qui cristallise à la refroidissement :
\[ \ce{(Z)-HOOC-CH=CH-COOH ->[\text{chaleur ou } h\nu][\text{catalyseur}] (E)-HOOC-CH=CH-COOH} \]
\textbf{Interprétation.} L'apport d'énergie (thermique ou lumineuse) affaiblit momentanément la liaison $\pi$, autorisant une rotation autour de l'axe C--C ; la forme (E), plus stable thermodynamiquement (encombrement stérique moindre entre les deux \ce{COOH}), se forme préférentiellement. Cette réaction est \emph{lente} et nécessite un catalyseur ou un rayonnement : elle confirme que la rotation autour de \ce{C=C} est bloquée dans les conditions ordinaires.
\end{experience}

\begin{savaistu}[La vision : une isomérisation Z/E dans l'œil]
Comment vois-tu ce texte ? Grâce à une isomérisation Z/E ultrarapide ! Dans les cellules de ta rétine (cônes et bâtonnets) se trouve un pigment visuel, la \emph{rhodopsine}, composé d'une protéine (l'opsine) et d'un chromophore : le \emph{11-(Z)-rétinal} (aussi appelé \emph{rétinal cis}). Lorsqu'un photon lumineux frappe le rétinal, sa double liaison en position 11 s'isomérise en une fraction de picoseconde (\SI{e-12}{\second}) : la molécule passe de la configuration \cle{Z} (forme coudée, adaptée à la poche de l'opsine) à la configuration \cle{E} (forme étirée, \emph{all-trans}). Ce changement brutal de géométrie déforme la protéine, ce qui déclenche une cascade de signaux électriques vers le cerveau : c'est la vision. Le rétinal dérive de la vitamine A (rétinol), abondante dans la carotte, la patate douce, le foie et l'\emph{huile de palme rouge} — un argument de plus pour la cuisine camerounaise ! Par ailleurs, l'industrie pharmaceutique accorde une attention extrême à la configuration Z/E des médicaments : l'isomère (Z) et l'isomère (E) d'un même principe actif peuvent avoir des effets biologiques très différents, voire opposés, car les enzymes et récepteurs de notre organisme reconnaissent la \emph{forme tridimensionnelle} précise des molécules (ex. : tamoxifène, certains anti-inflammatoires).
\end{savaistu}

\begin{aretenir}[L'essentiel sur l'isomérie Z/E]
\begin{itemize}
\item La double liaison \ce{C=C} ($\sigma + \pi$) \textbf{bloque la rotation} : la géométrie de l'alcène est figée.
\item L'isomérie Z/E n'existe que si \textbf{chaque} carbone de \ce{C=C} porte \textbf{deux groupes différents}.
\item La configuration s'attribue selon la \textbf{règle CIP} (numéro atomique du premier atome lié) : \cle{Z} (\emph{zusammen}) = groupes prioritaires du \textbf{même côté} ; \cle{E} (\emph{entgegen}) = de \textbf{part et d'autre}.
\item Z et E sont des \textbf{diastéréoisomères} : propriétés physiques différentes (ex. acide maléique (Z) $T_{\text{fus}} = \SI{131}{\celsius}$ ; acide fumarique (E) $T_{\text{fus}} = \SI{287}{\celsius}$), séparables, jamais énantiomères.
\item Le propène, but-1-ène, 2-méthylpropène (un carbone porte deux groupes identiques) ne présentent \textbf{pas} d'isomérie Z/E.
\end{itemize}
\end{aretenir}

\begin{exercice}[À toi de jouer]
Pour chacun des composés suivants, dire s'il présente l'isomérie Z/E ; si oui, représenter les deux stéréo-isomères (représentation de Cram ou perspective) et les nommer :
\begin{enumerate}
\item \ce{CH3-CH=CH-Cl} (1-chloroprop-1-ène) ;
\item \ce{CH3-C(CH3)=CH-CH3} (2-méthylbut-2-ène) ;
\item \ce{Cl-CH=CH-Cl} (1,2-dichloroéthène).
\end{enumerate}
\end{exercice}

\begin{corrige}[Exercice : isomérie Z/E]
\begin{enumerate}
\item \textbf{1-chloroprop-1-ène} (numérotation IUPAC : double liaison en position 1 $\rightarrow$ \ce{ClCH=CH-CH3}).
Carbone 1 (de la double liaison) : porte Cl (Z=17) et H (Z=1) $\rightarrow$ différents, priorité : Cl.
Carbone 2 : porte H (Z=1) et \ce{CH3} (C, Z=6) $\rightarrow$ différents, priorité : \ce{CH3}.
Condition remplie $\rightarrow$ \textbf{isomérie Z/E existe}.
\begin{itemize}
\item \textbf{(Z)-1-chloroprop-1-ène} : Cl et \ce{CH3} du même côté.
\item \textbf{(E)-1-chloroprop-1-ène} : Cl et \ce{CH3} de part et d'autre.
\end{itemize}
\item \textbf{2-méthylbut-2-ène} : \ce{CH3-C(CH3)=CH-CH3}.
Carbone 2 (de la double liaison) : porte \ce{CH3} et \ce{CH3} $\rightarrow$ \textbf{deux groupes identiques}.
\underline{Pas d'isomérie Z/E} (condition non remplie).
\item \textbf{1,2-dichloroéthène} : \ce{Cl-CH=CH-Cl}.
Chaque carbone porte Cl (Z=17) et H (Z=1) $\rightarrow$ différents, priorité : Cl.
Condition remplie $\rightarrow$ \textbf{isomérie Z/E existe}.
\begin{itemize}
\item \textbf{(Z)-1,2-dichloroéthène} : les deux Cl du même côté (moment dipolaire non nul).
\item \textbf{(E)-1,2-dichloroéthène} : les deux Cl de part et d'autre (moment dipolaire nul par symétrie).
\end{itemize}
\end{enumerate}
\end{corrige}

\coursec{Stéréo-isomérie de conformation}

Dans la section précédente, nous avons découvert l'isomérie Z/E des alcènes : la présence d'une double liaison \ce{C=C} \emph{bloque} la rotation et fige la molécule dans deux géométries distinctes, isolables, appelées diastéréoisomères. Mais que se passe-t-il autour d'une liaison \emph{simple} \ce{C-C} ? Ici, la rotation est libre : les groupes fixés sur les deux carbones peuvent tourner l'un par rapport à l'autre, comme les deux moitiés d'une toupie. La molécule adopte alors une infinité de dispositions spatiales, appelées \cle{conformations}. C'est l'objet de cette section : comprendre ces formes, les représenter, et classer leur stabilité.

\courssub{Notion de conformation : la rotation autour des liaisons simples}

\begin{definition}[Conformation]
On appelle \cle{conformation} d'une molécule toute disposition spatiale des atomes résultant de la \textbf{rotation autour des liaisons simples} $\sigma$ (\ce{C-C} ou \ce{C-O}), \textbf{sans rupture de liaison}. Deux conformations différentes d'une même molécule sont appelées \cle{conformères} (ou conformères rotationnels).
\end{definition}

L'idée intuitive est simple : prends un modèle moléculaire de l'éthane \ce{CH3-CH3}. Les deux groupes méthyle \ce{-CH3} sont reliés par une liaison simple qui joue le rôle d'un \emph{axe}. Tu peux faire tourner un méthyle par rapport à l'autre : à chaque instant, la molécule adopte une conformation différente, mais sa formule, sa connectivité et sa configuration restent identiques.

\begin{propriete}[Libre rotation et interconversion rapide]
La barrière énergétique à franchir pour passer d'une conformation à une autre est \textbf{faible} (quelques kJ/mol). À température ordinaire, l'agitation thermique suffit largement : les conformères s'\textbf{interconvertissent très rapidement} (de l'ordre de $10^{11}$ fois par seconde pour l'éthane). Il est donc \textbf{impossible d'isoler} un conformère à température ambiante : on ne distingue pas deux « molécules différentes », contrairement aux énantiomères ou aux isomères Z/E.
\end{propriete}

\begin{attention}[Ne pas confondre conformation et configuration]
Erreur classique au Baccalauréat ! Les \textbf{conformères} se transforment les uns en les autres par \emph{simple rotation} autour d'une liaison simple, sans casser aucune liaison : ce ne sont \textbf{pas} des isomères au sens strict. Les \textbf{énantiomères} et les \textbf{diastéréoisomères}, eux, sont des \emph{isomères de configuration} : pour passer de l'un à l'autre, il faut \emph{rompre et reformer des liaisons}. Ils sont donc isolables et possèdent des propriétés distinctes. Retiens : \textbf{conformation = rotation libre ; configuration = liaisons à casser}.
\end{attention}

\courssub{La représentation de Newman}

Pour comparer les conformations, le chimiste Melvin Newman a imaginé en 1952 une représentation ingénieuse : on \emph{regarde la molécule le long de l'axe de la liaison \ce{C-C}} qui nous intéresse.

\begin{methode}[Dessiner une représentation de Newman]
\begin{enumerate}
\item Placer l'\oe il \textbf{dans l'axe} de la liaison \ce{C-C} étudiée, de sorte que le carbone arrière soit caché derrière le carbone avant.
\item Représenter le \textbf{carbone avant} par un \textbf{point} (le centre), d'où partent ses trois liaisons disposées à \ang{120} les unes des autres.
\item Représenter le \textbf{carbone arrière} par un \textbf{cercle}, d'où partent ses trois liaisons, également à \ang{120}.
\item Orienter les liaisons du carbone arrière selon la conformation voulue : \emph{éclipsée} si elles se cachent derrière celles de l'avant, \emph{décalée} si elles s'intercalent entre elles.
\end{enumerate}
L'\cle{angle dièdre} $\theta$ est l'angle formé par deux liaisons (une de l'avant, une de l'arrière) vues en projection. Il varie de \ang{0} (éclipsé) à \ang{360}.
\end{methode}

\courssub{Conformations de l'éthane}

L'éthane \ce{CH3-CH3} est le cas le plus simple : chaque carbone porte trois atomes d'hydrogène. En tournant autour de la liaison \ce{C-C}, on rencontre deux conformations remarquables.

\begin{popfigure}
\begin{center}
\begin{tikzpicture}[scale=1.05]
% --- Éthane éclipsé ---
\begin{scope}
\draw[thick,popblue] (0,0) circle (0.95);
\foreach \a in {90,210,330}{
  \draw[thick,popdark] (0,0) -- (\a:1.05);
  \node[popdark] at (\a:1.35) {\small H};
}
\foreach \a in {90,210,330}{
  \draw[thick,poporange,dashed] (\a:0.95) -- (\a:1.15);
}
\node[popmuted] at (0,-1.7) {\small \textbf{Éclipsée} : $\theta = \ang{0}$};
\node[popmuted] at (0,-2.05) {\small gêne maximale, énergie maximale};
\end{scope}
% --- Éthane décalé ---
\begin{scope}[xshift=5.2cm]
\draw[thick,popblue] (0,0) circle (0.95);
\foreach \a in {90,210,330}{
  \draw[thick,popdark] (0,0) -- (\a:1.05);
  \node[popdark] at (\a:1.35) {\small H};
}
\foreach \a in {30,150,270}{
  \draw[thick,poporange] (\a:0.95) -- (\a:1.2);
  \node[poporange] at (\a:1.5) {\small H};
}
\draw[<->,popgreen,thick] (90:0.55) arc (90:30:0.55);
\node[popgreen] at (60:0.35) {\scriptsize $\theta$};
\node[popmuted] at (0,-1.7) {\small \textbf{Décalée} : $\theta = \ang{60}$};
\node[popmuted] at (0,-2.05) {\small gêne minimale, stabilité maximale};
\end{scope}
\end{tikzpicture}
\end{center}
\legende{Représentations de Newman de l'éthane : conformation éclipsée (les H de l'arrière, en pointillés, se cachent derrière ceux de l'avant) et conformation décalée (les H s'intercalent).}
\end{popfigure}

\begin{definition}[Conformations éclipsée et décalée]
\begin{itemize}
\item \textbf{Conformation éclipsée} ($\theta = \ang{0}$) : les liaisons \ce{C-H} du carbone arrière se trouvent exactement derrière celles du carbone avant. Les doublets d'électrons des liaisons se repoussent au maximum : c'est la conformation de \textbf{gêne maximale} et d'\textbf{énergie maximale}.
\item \textbf{Conformation décalée} ($\theta = \ang{60}$) : les liaisons de l'arrière s'intercalent entre celles de l'avant. Les répulsions sont minimales : c'est la conformation la \textbf{plus stable}.
\end{itemize}
\end{definition}

\begin{exemplebox}[Énergie de rotation de l'éthane]
L'écart d'énergie entre la conformation éclipsée et la conformation décalée de l'éthane vaut environ \SI{12}{kJ.mol^{-1}}. Cette valeur est \textbf{faible} : à \SI{25}{\celsius}, l'énergie d'agitation thermique disponible est de l'ordre de \SI{2,5}{kJ.mol^{-1}} par degré de liberté, et la distribution des énergies des chocs moléculaires permet de franchir sans peine cette barrière. Conclusion : la rotation autour de la liaison \ce{C-C} de l'éthane reste \textbf{quasiment libre} à température ambiante, et l'on ne peut pas isoler les conformères.
\end{exemplebox}

\courssub{Conformations du butane}

Le butane \ce{CH3-CH2-CH2-CH3} est plus riche : on observe la rotation autour de la liaison centrale \ce{C2-C3}. Chaque carbone de cette liaison porte alors un groupe méthyle \ce{CH3}, beaucoup plus encombrant qu'un hydrogène. On suit l'angle dièdre $\theta$ entre les deux groupes \ce{CH3}.

\begin{definition}[Conformations remarquables du butane]
\begin{itemize}
\item \textbf{Conformation anti} ($\theta = \ang{180}$) : les deux groupes \ce{CH3} sont à l'opposé l'un de l'autre. L'encombrement est minimal : c'est la conformation \textbf{la plus stable}.
\item \textbf{Conformation gauche} ($\theta = \ang{60}$ ou $\ang{300}$) : les deux \ce{CH3} sont décalés mais voisins ; il subsiste une interaction appelée \emph{gêne gauche}. Stabilité intermédiaire.
\item \textbf{Conformations éclipsées} ($\theta = \ang{0}$, $\ang{120}$, $\ang{240}$) : maxima d'énergie. Le cas $\theta = \ang{0}$, où les deux \ce{CH3} s'éclipsent mutuellement, est le \textbf{plus défavorable}.
\end{itemize}
\end{definition}

\begin{popfigure}
\begin{center}
\begin{tikzpicture}
\begin{axis}[
  width=0.95\linewidth, height=6.2cm,
  xlabel={Angle dièdre $\theta$ (degrés)},
  ylabel={Énergie potentielle relative (\si{kJ.mol^{-1}})},
  xmin=0, xmax=360, ymin=0, ymax=22,
  xtick={0,60,120,180,240,300,360},
  ytick={0,5,10,15,20},
  grid=both, grid style={popline!40},
  axis line style={popmuted},
  label style={popdark},
]
\addplot[popcurve,domain=0:360,samples=300]
  {9.5 + 2.53*cos(x) + 6.97*cos(3*x)};
\node[popmuted,font=\scriptsize] at (axis cs:0,19.5) {éclipsée Me/Me};
\node[popmuted,font=\scriptsize] at (axis cs:60,4.5) {gauche};
\node[popmuted,font=\scriptsize] at (axis cs:120,15.5) {éclipsée Me/H};
\node[popmuted,font=\scriptsize] at (axis cs:180,0.5) {anti};
\node[popmuted,font=\scriptsize] at (axis cs:240,15.5) {éclipsée Me/H};
\node[popmuted,font=\scriptsize] at (axis cs:300,4.5) {gauche};
\node[popmuted,font=\scriptsize] at (axis cs:360,19.5) {éclipsée Me/Me};
\end{axis}
\end{tikzpicture}
\end{center}
\legende{Énergie potentielle du butane en fonction de l'angle dièdre $\theta$ entre les deux groupes méthyle. Les minima correspondent aux conformations anti (la plus stable) et gauche ; les maxima aux conformations éclipsées.}
\end{popfigure}

\begin{exoresolu}[Lecture du diagramme d'énergie du butane]
Énoncé. Le diagramme ci-dessus donne l'énergie potentielle du butane en fonction de l'angle dièdre. (1) Quelle est la conformation la plus stable ? Justifier. (2) L'écart d'énergie entre la conformation anti et la conformation gauche est d'environ \SI{3,8}{kJ.mol^{-1}}, et la barrière maximale (éclipsée totale, $\theta = \ang{0}$) d'environ \SI{19}{kJ.mol^{-1}}. Peut-on isoler la conformation anti à température ambiante ?
\tcblower
Solution détaillée.
\begin{enumerate}
\item La conformation la plus stable est celle d'\textbf{énergie minimale}, c'est-à-dire la conformation \textbf{anti} ($\theta = \ang{180}$) : les deux groupes \ce{CH3}, les plus encombrants, sont le plus éloignés possible, ce qui minimise les répulsions stériques.
\item Pour isoler une conformation, il faudrait que la barrière d'énergie à franchir soit grande devant l'énergie d'agitation thermique. Or la barrière maximale ne vaut que \SI{19}{kJ.mol^{-1}}, valeur comparable à celle de l'éthane (\SI{12}{kJ.mol^{-1}}) et très inférieure à l'énergie d'une liaison covalente (environ \SI{350}{kJ.mol^{-1}} pour \ce{C-C}). Les interconversions sont donc extrêmement rapides à \SI{25}{\celsius} : \textbf{il est impossible d'isoler} la conformation anti. Le butane est un mélange dynamique où la forme anti est simplement la \emph{plus peuplée} (environ 70 \% des molécules à un instant donné).
\end{enumerate}
\end{exoresolu}

\courssub{Conformations du cyclohexane : chaise et bateau}

Le cyclohexane \ce{C6H12} est un cycle à six atomes de carbone. S'il était \emph{plan}, ses angles \ce{C-C-C} vaudraient \ang{120}, très éloignés de l'angle tétraédrique idéal de \ang{109,5} : le cycle serait fortement tendu. En réalité, le cycle se \emph{plisse} grâce aux rotations autour des liaisons simples, ce qui lui permet de retrouver des angles proches de \ang{109}.

\begin{definition}[Conformations chaise et bateau du cyclohexane]
\begin{itemize}
\item \textbf{Conformation chaise} : quatre carbones sont sensiblement coplanaires, les deux autres se situant de part et d'autre de ce plan. Tous les angles sont proches de \ang{109} et toutes les liaisons sont \emph{décalées} : c'est la conformation \textbf{la plus stable}, pratiquement sans tension.
\item \textbf{Conformation bateau} : deux carbones « proue » et « poupe » sont relevés du même côté. Des liaisons y sont \emph{éclipsées} et deux hydrogènes (dits « de proue ») se font face : la gêne est importante, la conformation est \textbf{moins stable} d'environ \SI{25}{kJ.mol^{-1}}.
\end{itemize}
\end{definition}

\begin{popfigure}
\begin{center}
\begin{tikzpicture}[scale=0.9,line join=round,line cap=round]
% --- Chaise ---
\begin{scope}
% Carbones : C1(Up), C2(Down), C3(Up), C4(Down), C5(Up), C6(Down)
\coordinate (C1) at (0, 1.2);
\coordinate (C2) at (1.0, 1.6);
\coordinate (C3) at (2.0, 1.2);
\coordinate (C4) at (2.5, 0.2);
\coordinate (C5) at (1.5, -0.2);
\coordinate (C6) at (0.5, 0.2);
% Squelette
\draw[thick,popblue] (C1)--(C2)--(C3)--(C4)--(C5)--(C6)--cycle;
% Liaisons Axiales (verticales) et Équatoriales (sortantes)
% C1 (Up) : Axiale UP, Eq DOWN-RIGHT
\draw[thick,poporange] (C1) -- ++(0,0.7) node[above,poporange,font=\scriptsize] {H};
\draw[thick,popgreen] (C1) -- ++(0.6,-0.5) node[right,popgreen,font=\scriptsize] {H};
% C2 (Down) : Axiale DOWN, Eq UP-RIGHT
\draw[thick,poporange] (C2) -- ++(0,-0.7) node[below,poporange,font=\scriptsize] {H};
\draw[thick,popgreen] (C2) -- ++(0.6,0.5) node[right,popgreen,font=\scriptsize] {H};
% C3 (Up) : Axiale UP, Eq DOWN-RIGHT
\draw[thick,poporange] (C3) -- ++(0,0.7) node[above,poporange,font=\scriptsize] {H};
\draw[thick,popgreen] (C3) -- ++(0.6,-0.5) node[right,popgreen,font=\scriptsize] {H};
% C4 (Down) : Axiale DOWN, Eq UP-LEFT
\draw[thick,poporange] (C4) -- ++(0,-0.7) node[below,poporange,font=\scriptsize] {H};
\draw[thick,popgreen] (C4) -- ++(-0.6,0.5) node[left,popgreen,font=\scriptsize] {H};
% C5 (Up) : Axiale UP, Eq DOWN-LEFT
\draw[thick,poporange] (C5) -- ++(0,0.7) node[above,poporange,font=\scriptsize] {H};
\draw[thick,popgreen] (C5) -- ++(-0.6,-0.5) node[left,popgreen,font=\scriptsize] {H};
% C6 (Down) : Axiale DOWN, Eq UP-LEFT
\draw[thick,poporange] (C6) -- ++(0,-0.7) node[below,poporange,font=\scriptsize] {H};
\draw[thick,popgreen] (C6) -- ++(-0.6,0.5) node[left,popgreen,font=\scriptsize] {H};

\node[popmuted] at (1.25,-1.5) {\small \textbf{Chaise} (la plus stable)};
\node[poporange,font=\scriptsize] at (1.25,2.5) {Axiales (verticales)};
\node[popgreen,font=\scriptsize] at (3.2,0.5) {Équatoriales};
\end{scope}
% --- Bateau ---
\begin{scope}[xshift=7.5cm]
\coordinate (P1) at (0,0.6); \coordinate (P2) at (1.0,-0.1);
\coordinate (P3) at (2.0,-0.1); \coordinate (P4) at (3.0,0.6);
\coordinate (P5) at (2.0,-0.75); \coordinate (P6) at (1.0,-0.75);
\draw[thick,poppurple] (P1)--(P2)--(P3)--(P4)--(P5)--(P6)--cycle;
\draw[thick,poporange] (P1)--(0,1.4) node[above,poporange,font=\scriptsize] {H};
\draw[thick,poporange] (P4)--(3.0,1.4) node[above,poporange,font=\scriptsize] {H};
\draw[<->,popink,thick] (0.15,1.15) -- (2.85,1.15);
\node[popink,font=\scriptsize] at (1.5,1.6) {Gêne entre H de « proue »};
\node[popmuted] at (1.5,-1.5) {\small \textbf{Bateau} (moins stable)};
\node[popmuted,font=\scriptsize] at (1.5,-1.9) {Liaisons éclipsées sur les flancs};
\end{scope}
\end{tikzpicture}
\end{center}
\legende{Conformations chaise et bateau du cyclohexane. Dans la chaise, chaque carbone porte une liaison \textbf{axiale} (verticale, en orange) et une liaison \textbf{équatoriale} (dans le plan moyen du cycle, en vert).}
\end{popfigure}

\begin{propriete}[Positions axiales et équatoriales]
Dans la conformation chaise, les douze liaisons \ce{C-H} se répartissent en deux familles :
\begin{itemize}
\item six liaisons \cle{axiales}, parallèles à l'axe du cycle (verticales), dirigées alternativement vers le haut et vers le bas ;
\item six liaisons \cle{équatoriales}, dirigées vers l'extérieur, sensiblement dans le plan moyen du cycle.
\end{itemize}
Un substituant volumineux (par exemple \ce{CH3}) préfère la position \textbf{équatoriale}, où la gêne stérique est moindre. Par \emph{inversion de cycle} (la chaise « bascule »), les positions axiales et équatoriales s'échangent.
\end{propriete}

\begin{attention}[Pièges fréquents]
\begin{itemize}
\item Ne dis jamais que le cyclohexane est \emph{plan} : un cycle plan à six carbones aurait des angles de \ang{120} et serait très instable.
\item Dans la forme chaise, ne place pas toutes les liaisons axiales du même côté : elles alternent haut/bas autour du cycle.
\item La conformation bateau n'est \emph{pas} un isomère de la chaise : ce sont deux conformères du même cyclohexane, qui s'interconvertissent par rotation autour des liaisons simples.
\end{itemize}
\end{attention}

\begin{savaistu}[Le glucose aussi s'assoit sur une chaise]
La forme chaise du cyclohexane est si stable que la nature l'a adoptée partout : la plupart des sucres, comme le \textbf{glucose} présent dans le miel, les jus de fruits et le jus de bissap sucré, possèdent un cycle à six atomes (cinq carbones et un oxygène) qui adopte précisément cette géométrie chaise. Mieux : c'est parce que tous ses gros groupes \ce{-OH} et \ce{-CH2OH} occupent des positions \emph{équatoriales} que le glucose est le sucre le plus répandu du vivant. La stéréochimie n'est pas un jeu de salon : elle décide de la structure des molécules de la vie !
\end{savaistu}

\begin{aretenir}[L'essentiel de la section]
\begin{itemize}
\item Une \textbf{conformation} est une disposition spatiale obtenue par \textbf{rotation autour des liaisons simples}, sans rupture de liaison ; les conformères s'interconvertissent rapidement et ne sont \textbf{pas isolables}.
\item La \textbf{représentation de Newman} observe la molécule dans l'axe d'une liaison \ce{C-C} : carbone avant = point, carbone arrière = cercle.
\item Éthane : conformation \textbf{éclipsée} ($\theta = \ang{0}$, énergie maximale) et \textbf{décalée} ($\theta = \ang{60}$, la plus stable) ; barrière d'environ \SI{12}{kJ.mol^{-1}}.
\item Butane : conformation \textbf{anti} ($\theta = \ang{180}$, la plus stable), \textbf{gauche} ($\ang{60}$), et éclipsées (maxima d'énergie).
\item Cyclohexane : conformation \textbf{chaise} (la plus stable, angles proches de \ang{109}) et \textbf{bateau} (moins stable) ; dans la chaise, on distingue les positions \textbf{axiales} et \textbf{équatoriales}.
\end{itemize}
\end{aretenir}

\coursec{Importance de la stéréo-isomérie et synthèse}

Dans la section précédente, nous avons vu qu'une même molécule peut adopter plusieurs \cle{conformations} par simple rotation autour de ses liaisons simples, et que certaines de ces formes spatiales sont plus stables que d'autres. Nous avons ainsi achevé notre tour d'horizon des différentes façons dont des molécules de même formule brute peuvent différer. Mais une question fondamentale demeure : \emph{à quoi cela sert-il, concrètement ?} Pourquoi un industriel de la SONARA, un pharmacien de Yaoundé ou un chimiste des arômes devrait-il se soucier de l'arrangement des atomes \emph{dans l'espace} ? La réponse tient en une phrase : la vie elle-même est chirale. C'est ce que nous allons découvrir dans cette dernière section, avant de dresser le bilan complet de la leçon.

\courssub{Pourquoi la stéréochimie est-elle si importante ?}

\begin{definition}[Stéréochimie]
La \cle{stéréochimie} est la partie de la chimie qui étudie l'\textbf{arrangement des atomes dans l'espace} au sein des molécules, ainsi que les \textbf{conséquences de cet arrangement} sur les propriétés physiques, chimiques et biologiques de ces molécules.
\end{definition}

Commençons par une intuition simple. Ta main droite et ta main gauche sont des objets \emph{chiraux} : ce sont des images miroir non superposables. Imagine maintenant un gant droit : seule ta main droite peut y entrer confortablement. La main gauche, pourtant composée des mêmes « éléments » (doigts, paume, pouce), ne convient pas. Le gant « reconnaît » la chiralité de la main.

Il en va exactement de même dans le monde vivant. Les molécules du vivant — protéines, enzymes, récepteurs, ADN — sont elles-mêmes \textbf{chirales}, car elles sont construites presque exclusivement à partir d'acides aminés de configuration L et de sucres de série D. Par conséquent, un récepteur biologique se comporte comme un « gant » tridimensionnel : il peut reconnaître un énantiomère et ignorer (ou mal reconnaître) son image miroir.

\begin{propriete}[Reconnaissance chirale dans le vivant]
Les \cle{récepteurs biologiques} (enzymes, récepteurs olfactifs, récepteurs membranaires) étant chiraux, deux \textbf{énantiomères} d'une même molécule peuvent provoquer des \textbf{effets biologiques totalement différents} : odeurs différentes, activité thérapeutique différente, voire toxicité opposée.
\end{propriete}

\begin{popfigure}
\begin{center}
\begin{tikzpicture}[scale=1.0, every node/.style={font=\small}]
\% ---- Serrure (récepteur) de gauche : bonne clé ----
\fill[popblueL, draw=popblue, thick, rounded corners=6pt] (0,0) rectangle (3.2,3.4);
\node[popdark, font=\small\bfseries] at (1.6,3.7) {Récepteur biologique (chiral)};
\% empreinte en L
\fill[white, draw=popblue, thick] (0.9,0.6) rectangle (1.7,2.0);
\fill[white, draw=popblue, thick] (0.9,1.5) rectangle (2.5,2.3);
\% clé S qui entre
\fill[popgreen, draw=popgreen!60!black, thick] (4.1,1.0) rectangle (4.9,2.0);
\fill[popgreen, draw=popgreen!60!black, thick] (4.1,1.6) rectangle (5.7,2.3);
\draw[popgreen!60!black, thick] (5.7,1.95) circle (0.35);
\node[popgreen!50!black, font=\small\bfseries] at (5.0,0.5) {Énantiomère adapté};
\draw[-{Stealth[length=3mm]}, popgreen!60!black, very thick] (3.9,1.6) -- (2.7,1.6);
\node[popgreen!50!black] at (1.6,-0.4) {\textbf{Fixation : effet biologique}};

\% ---- Serrure de droite : mauvaise clé ----
\begin{scope}[xshift=9.2cm]
\fill[popblueL, draw=popblue, thick, rounded corners=6pt] (0,0) rectangle (3.2,3.4);
\node[popdark, font=\small\bfseries] at (1.6,3.7) {Même récepteur};
\fill[white, draw=popblue, thick] (0.9,0.6) rectangle (1.7,2.0);
\fill[white, draw=popblue, thick] (0.9,1.5) rectangle (2.5,2.3);
\% clé R (image miroir) bloquée
\fill[poporange, draw=poporange!70!black, thick] (4.1,1.0) rectangle (4.9,2.0);
\fill[poporange, draw=poporange!70!black, thick] (3.3,1.6) rectangle (4.9,2.3);
\draw[poporange!70!black, thick] (4.9,1.95) circle (0.35);
\node[poporange!70!black, font=\small\bfseries] at (4.2,0.5) {Énantiomère image miroir};
\draw[-{Stealth[length=3mm]}, poporange!70!black, very thick] (3.1,1.6) -- (2.75,1.6);
\draw[red, very thick] (2.45,1.2) -- (3.05,2.0);
\draw[red, very thick] (2.45,2.0) -- (3.05,1.2);
\node[red] at (1.6,-0.4) {\textbf{Pas de fixation : pas d'effet}};
\end{scope}
\end{tikzpicture}
\end{center}
\legende{Modèle « clé--serrure » : le récepteur biologique est chiral. Un seul des deux énantiomères (la bonne « clé ») épouse la forme du site actif et déclenche l'effet biologique ; son image miroir ne se fixe pas.}
\end{popfigure}

\courssub{Des énantiomères aux effets différents : exemples}

\courssub{Les arômes : limonène et carvone}

Nos récepteurs olfactifs sont chiraux : ils distinguent les énantiomères comme notre nez distingue une orange d'un citron.

\begin{exemplebox}[Limonène et carvone : même formule brute, odeurs différentes]
Le \textbf{limonène} ($\ce{C10H16}$) possède un carbone asymétrique :
\begin{itemize}
\item l'énantiomère \textbf{(R)-limonène} a une odeur d'\textbf{orange} ;
\item l'énantiomère \textbf{(S)-limonène} a une odeur de \textbf{citron}.
\end{itemize}
De même, la \textbf{carvone} ($\ce{C10H14O}$) :
\begin{itemize}
\item la \textbf{(R)-carvone} a une odeur de \textbf{menthe} verte ;
\item la \textbf{(S)-carvone} a une odeur de \textbf{carvi} (graine épice).
\end{itemize}
Les deux énantiomères ont \emph{exactement} les mêmes propriétés physiques (même température d'ébullition, même densité, même indice de réfraction) : seul un environnement chiral — ici le récepteur olfactif — les distingue.
\end{exemplebox}

\courssub{Les médicaments : l'ibuprofène}

L'\textbf{ibuprofène} est un anti-inflammatoire très répandu, vendu dans toutes les pharmacies de Douala ou de Bamenda. Sa molécule possède un carbone asymétrique ; elle existe donc sous deux formes énantiomères, (R) et (S).

\begin{exemplebox}[L'ibuprofène : un racémique dont un seul énantiomère est actif]
Le médicament vendu en pharmacie est un \textbf{mélange racémique} : il contient \SI{50}{\percent} d'énantiomère (S) et \SI{50}{\percent} d'énantiomère (R). Or \textbf{seul l'énantiomère (S) est actif} : il inhibe l'enzyme responsable de la synthèse des prostaglandines (médiateurs de la douleur et de l'inflammation). L'énantiomère (R) est pratiquement \textbf{inactif} à ce niveau — l'organisme en convertit d'ailleurs lentement une partie en forme (S).

\medskip
\textbf{Conséquence chiffrée :} dans un comprimé dosé à \SI{400}{mg} d'ibuprofène racémique, seuls \SI{200}{mg} (la forme (S)) exercent réellement l'effet thérapeutique recherché. Pour obtenir le même effet avec l'énantiomère pur (S) (commercialisé sous le nom de \emph{dexibuprofène}), une dose de \SI{200}{mg} suffit.
\end{exemplebox}

\begin{attention}[Erreur fréquente : confondre propriétés physiques et effets biologiques]
Deux énantiomères ont les \textbf{mêmes propriétés physiques} (températures de fusion et d'ébullition, densité, solubilité, indice de réfraction) et les \textbf{mêmes propriétés chimiques} vis-à-vis de réactifs \emph{achiraux}. Ils ne se distinguent que par :
\begin{itemize}
\item leur \textbf{pouvoir rotatoire} (déviations de la lumière polarisée \emph{opposées}) ;
\item leur comportement vis-à-vis d'un environnement \textbf{chiral} (réactif chiral, enzyme, récepteur).
\end{itemize}
Il est donc \textbf{faux} d'écrire : « les deux énantiomères ont les mêmes propriétés, donc les mêmes effets sur l'organisme ». Leurs effets biologiques peuvent être \emph{totalement différents}, comme le montre tragiquement l'exemple suivant.
\end{attention}

\courssub{Le drame de la thalidomide : une leçon historique}

La \textbf{thalidomide} est un médicament prescrit à la fin des années 1950 comme sédatif et contre les nausées des femmes enceintes. Sa molécule possède un carbone asymétrique :
\begin{itemize}
\item l'énantiomère \textbf{(R)} est le \textbf{sédatif} recherché ;
\item l'énantiomère \textbf{(S)} est \textbf{tératogène} : il provoque de graves malformations du fœtus (enfants nés avec des membres atrophiés ou absents).
\end{itemize}
Le médicament était vendu sous forme racémique. Pire : même si l'on avait administré le seul énantiomère (R), l'organisme le \emph{racémise} partiellement (les deux formes se convertissent l'une en l'autre dans le sang). Entre 1957 et 1961, on estime que plus de dix mille enfants dans le monde ont été victimes de malformations. C'est l'une des plus grandes catastrophes sanitaires de l'histoire du médicament.

\begin{savaistu}[La thalidomide a révolutionné le contrôle des médicaments]
Après le scandale de la thalidomide (années 1960), les agences du médicament (FDA aux États-Unis, puis leurs homologues européennes et mondiales) ont durci leurs exigences : avant toute mise sur le marché d'un médicament chiral, il faut désormais \textbf{étudier séparément chaque énantiomère} — son activité, sa toxicité, son devenir dans l'organisme. De nombreux médicaments modernes sont ainsi produits sous forme d'\textbf{un seul énantiomère} (on parle de synthèse \emph{énantiosélective}), ce qui a valu le prix Nobel de chimie 2001 à Knowles, Noyori et Sharpless pour leurs catalyseurs chiraux. La stéréochimie est devenue un enjeu de santé publique — y compris au Cameroun, où les médicaments importés sont contrôlés selon ces mêmes standards internationaux.
\end{savaistu}

\courssub{Stéréochimie et alimentation}

La chiralité gouverne aussi ce que nous mangeons :

\begin{itemize}
\item \textbf{Les sucres naturels} (glucose, fructose du miel ou du jus de bissap sucré) appartiennent presque tous à la série \textbf{D}. Nos enzymes ne savent métaboliser que ces formes : le glucose « image miroir » (série L) aurait le même goût sucré mais ne serait \emph{pas assimilé} — c'est d'ailleurs le principe de certains édulcorants.
\item \textbf{Les acides aminés} naturels, constituants des protéines (viande, poisson fumé, niébé), sont presque tous de configuration \textbf{L}. Les protéines que nous consommons sont donc chirales, et nos enzymes digestives, chirales elles aussi, les reconnaissent parfaitement.
\item \textbf{Les arômes alimentaires} : l'industrie agroalimentaire distingue les énantiomères pour reproduire fidèlement les goûts naturels. Un arôme « menthe » synthétique contenant le mauvais énantiomère de carvone sentirait le carvi !
\end{itemize}

\courssub{Bilan de la leçon : classification des isoméries}

Nous pouvons maintenant rassembler tout ce que nous avons appris dans cette leçon. Deux composés sont \textbf{isomères} lorsqu'ils ont la \emph{même formule brute} mais des structures différentes. On distingue deux grandes familles.

\begin{popfigure}
\begin{center}
\begin{tikzpicture}[
  grow=right,
  level 1/.style={sibling distance=52mm, level distance=34mm},
  level 2/.style={sibling distance=13mm, level distance=40mm},
  every node/.style={font=\small},
  edge from parent/.style={draw=popmuted, thick}
]
\node[fill=popink, text=white, rounded corners=4pt, font=\small\bfseries, align=center] {ISOMÈRES\\ (même formule brute)}
  child { node[fill=poppurpleL, draw=poppurple, rounded corners=4pt, font=\small\bfseries, align=center] {Stéréo-isomères\\ (même enchaînement)}
    child { node[fill=white, draw=poppurple, rounded corners=3pt, align=center] {Conformations\\ (éclipsée, décalée,\\ chaise, bateau)}
      edge from parent node[above, font=\scriptsize, popmuted] {} }
    child { node[fill=white, draw=poppurple, rounded corners=3pt, align=center] {Z / E\\ \ce{(E)}- et \ce{(Z)}-but-2-ène} }
    child { node[fill=white, draw=poppurple, rounded corners=3pt, align=center] {Diastéréoisomères\\ (2 $C^*$ non images\\ miroir)} }
    child { node[fill=white, draw=poppurple, rounded corners=3pt, align=center] {Énantiomères\\ (R)- et (S)-acide\\ lactique} }
  }
  child { node[fill=poporangeL, draw=poporange, rounded corners=4pt, font=\small\bfseries, align=center] {Isomères de\\ constitution}
    child { node[fill=white, draw=poporange, rounded corners=3pt, align=center] {Fonction\\ \ce{C2H5OH} / \ce{CH3-O-CH3}} }
    child { node[fill=white, draw=poporange, rounded corners=3pt, align=center] {Position\\ butan-1-ol /\\ butan-2-ol} }
    child { node[fill=white, draw=poporange, rounded corners=3pt, align=center] {Chaîne\\ butane /\\ 2-méthylpropane} }
  };
\end{tikzpicture}
\end{center}
\legende{Carte de classification complète des isoméries, avec un exemple par type. Les stéréo-isomères se répartissent en isomères de \emph{configuration} (énantiomères, diastéréoisomères, Z/E : interconversion impossible sans rompre une liaison) et isomères de \emph{conformation} (interconversion par simple rotation).}
\end{popfigure}

\begin{aretenir}[Tableau récapitulatif des types d'isomérie]
\begin{center}
\begin{tabularx}{\linewidth}{|l|X|X|}
\hline
\rowcolor{popblueL} \textbf{Type d'isomérie} & \textbf{Critère distinctif} & \textbf{Exemple ($\ce{C4H10O}$, $\ce{C4H10}$, $\ce{C4H8}$...)} \\
\hline
Chaîne & Squelette carboné différent & butane / 2-méthylpropane \\
\hline
Position & Position différente du groupe caractéristique & butan-1-ol / butan-2-ol \\
\hline
Fonction & Fonctions chimiques différentes & éthanol \ce{C2H5OH} / méthoxyméthane \ce{CH3-O-CH3} \\
\hline
Énantiomères & Images miroir non superposables (1 $C^*$) & (R)- et (S)-acide lactique \\
\hline
Diastéréoisomères & Stéréo-isomères non images miroir (2 $C^*$ ou plus) & acides tartriques (2R,3R) et (2R,3S) \\
\hline
Z / E & Pas de rotation autour de \ce{C=C} & (Z)- et (E)-but-2-ène \\
\hline
Conformations & Rotation autour d'une liaison simple & éthane éclipsé / décalé ; cyclohexane chaise / bateau \\
\hline
\end{tabularx}
\end{center}
\end{aretenir}

\courssub{Guide de choix de la représentation}

Face à une molécule, quelle représentation adopter ? Retiens ce guide pratique :

\begin{itemize}
\item \textbf{Représentation de Cram} (perspective : traits pleins dans le plan, triangle plein vers l'avant, hachures vers l'arrière) : c'est la représentation \emph{universelle} pour montrer la géométrie d'un carbone asymétrique et comparer deux énantiomères. Exemple : \chemfig{C(-[2]H)(<[5]CH_3)(<:[7]OH)-COOH} pour l'acide lactique.
\item \textbf{Projection de Fischer} : idéale pour les molécules possédant \emph{un ou plusieurs} carbones asymétriques alignés (sucres, acides aminés, acide tartrique). La chaîne principale est verticale, les liaisons horizontales sortent du plan vers le lecteur. Elle permet de comparer rapidement énantiomères et diastéréoisomères.
\item \textbf{Représentation de Newman} : réservée à l'étude des \emph{conformations} — on regarde la molécule le long d'une liaison \ce{C-C} pour distinguer formes éclipsées et décalées, et analyser les encombrements stériques.
\end{itemize}

\begin{methode}[Classer un couple d'isomères en trois questions]
Devant deux formules proposées à l'examen, pose-toi dans l'ordre :
\begin{enumerate}
\item \textbf{Même enchaînement d'atomes ?} Si les atomes ne sont pas reliés dans le même ordre (squelette, position ou fonction différents), ce sont des \textbf{isomères de constitution} : précise chaîne, position ou fonction.
\item \textbf{Si oui : sont-ils images l'un de l'autre dans un miroir ?} Si oui et non superposables : ce sont des \textbf{énantiomères}. Si non : ce sont des \textbf{diastéréoisomères} (dont les couples Z/E autour d'une double liaison).
\item \textbf{Peut-on passer de l'un à l'autre par simple rotation autour d'une liaison simple ?} Si oui : ce ne sont pas des isomères de configuration mais des \textbf{conformations} différentes de la même molécule.
\end{enumerate}
\end{methode}

\begin{exoresolu}[Identifier la relation d'isomérie entre deux composés]
\textbf{Énoncé.} On considère les quatre composés suivants : (A) butan-1-ol \ce{CH3-CH2-CH2-CH2OH} ; (B) butan-2-ol \ce{CH3-CH2-CHOH-CH3} ; (C) l'énantiomère (R) du butan-2-ol ; (D) l'énantiomère (S) du butan-2-ol. Préciser la relation d'isomérie entre : 1) A et B ; 2) C et D ; 3) B et le mélange équimolaire de C et D.
\tcblower
\textbf{Solution détaillée.}
\begin{enumerate}
\item \textbf{A et B} : même formule brute $\ce{C4H10O}$, même fonction alcool, mais le groupe \ce{-OH} est en position 1 dans A et en position 2 dans B. Ce sont des \textbf{isomères de position}.
\item \textbf{C et D} : le carbone 2 du butan-2-ol porte quatre groupes différents (\ce{H}, \ce{OH}, \ce{CH3}, \ce{C2H5}) : c'est un carbone asymétrique. (R) et (S) sont images miroir non superposables : ce sont des \textbf{énantiomères}.
\item \textbf{B et le mélange C+D} : le mélange équimolaire des deux énantiomères est un \textbf{mélange racémique}, optiquement inactif (pouvoir rotatoire nul). Le composé B, désigné par le nom « butan-2-ol » sans précision stéréochimique, correspond usuellement à ce mélange racémique. B et le mélange C+D sont donc \textbf{identiques} (même mélange racémique). Si B désignait un énantiomère pur, il serait optiquement actif, ce qui n'est pas le cas standard.
\end{enumerate}
\end{exoresolu}

\begin{exercice}[À toi de jouer]
1) L'acide tartrique \ce{HOOC-CHOH-CHOH-COOH} possède deux carbones asymétriques. Combien de stéréo-isomères de configuration comporte-t-il au maximum ? En réalité, on n'en compte que trois : expliquer pourquoi (penser à un plan de symétrie).
2) Le (S)-limonène sent le citron. Que peut-on dire de son pouvoir rotatoire par rapport à celui du (R)-limonène ?
3) Pourquoi un mélange racémique d'ibuprofène est-il optiquement inactif alors que chaque énantiomère dévie la lumière polarisée ?
\end{exercice}

\begin{corrige}[Exercice]
1) Avec deux $C^*$, le maximum est $2^2 = 4$. Mais les formes (2R,3S) et (2S,3R) sont \emph{identiques} : la molécule possède un plan de symétrie interne, c'est la forme \emph{méso}, achirale. Restent donc trois stéréo-isomères : (2R,3R), (2S,3S) — énantiomères entre eux — et la forme méso, diastéréoisomère des deux précédentes.
2) Les pouvoirs rotatoires de deux énantiomères sont \textbf{opposés} : même valeur absolue, signes contraires. Si le (R)-limonène est dextrogyre, le (S)-limonène est lévogyre de la même amplitude.
3) Le mélange racémique contient autant de molécules dextrogyres que de lévogyres : les déviations de la lumière polarisée se compensent exactement, d'où un pouvoir rotatoire global nul.
\end{corrige}

\begin{aretenir}[L'essentiel de la section]
\begin{itemize}
\item La stéréochimie étudie l'arrangement des atomes \textbf{dans l'espace} et ses conséquences.
\item Le vivant est chiral : récepteurs, enzymes et récepteurs olfactifs distinguent les énantiomères (limonène, carvone, ibuprofène, thalidomide).
\item Deux énantiomères ont les mêmes propriétés physiques mais des \textbf{effets biologiques potentiellement très différents} : le contrôle stéréochimique des médicaments est obligatoire.
\item Isomérie de constitution (chaîne, position, fonction) $\neq$ stéréo-isomérie de configuration (énantiomères, diastéréoisomères, Z/E) $\neq$ conformations.
\item Cram pour la perspective, Fischer pour les molécules à $C^*$, Newman pour les conformations.
\end{itemize}
\end{aretenir}

\newpage

\coursec{Travaux pratiques}

\courssub{Objectifs du TP}

À l'issue de cette séance de travaux pratiques, tu seras capable de :

\begin{itemize}
\item mettre en évidence expérimentalement le \cle{pouvoir rotatoire} d'une solution de saccharose à l'aide de deux filtres polarisants ;
\item observer qualitativement la \cle{déviation du plan de polarisation} de la lumière et en déterminer le sens (dextrogyre ou lévogyre) ;
\item construire, à l'aide de modèles moléculaires, les deux \cle{énantiomères} de l'acide lactique et vérifier qu'ils sont images l'un de l'autre dans un miroir et \cle{non superposables} ;
\item construire les conformations \cle{éclipsée} et \cle{décalée} de l'éthane et les relier à la représentation de Newman ;
\item relier l'activité optique d'une substance à la présence d'un \cle{carbone asymétrique} dans sa molécule.
\end{itemize}

Ce TP illustre directement la famille de situations \og comparaison de deux énantiomères \fg{} : tu vas \emph{voir} de tes propres yeux l'effet de la chiralité sur la lumière, puis le \emph{toucher} du bout des doigts avec les modèles moléculaires.

\courssub{Matériel et produits}

\begin{center}
\rowcolors{2}{popblueL}{white}
\begin{tabularx}{\linewidth}{|l|X|l|}
\hline
\rowcolor{popblue!40} \textbf{Matériel / produit} & \textbf{Utilisation} & \textbf{Alternative locale} \\
\hline
2 filtres polarisants (polaroïds) & Polariseur et analyseur & 2 verres de lunettes de soleil polarisées, ou 2 écrans de téléphone cassés (dalle polarisante récupérée) \\
\hline
Lampe de poche ou lampe de bureau & Source de lumière blanche & Lampe torche de téléphone portable \\
\hline
Cristallisoir (ou cuve en verre) & Contenir la solution de sucre & Verre à boire transparent à fond plat, bocal propre \\
\hline
Sucre de canne en poudre (\SI{150}{g}) & Substance optiquement active (saccharose) & Miel clair, sirop de glucose, jus de canne concentré \\
\hline
Eau distillée ou eau bouillie refroidie (\SI{100}{mL}) & Solvant & Eau de pluie bouillie et filtrée \\
\hline
Balance, spatule, agitateur en verre & Pesée et dissolution & Cuillère à café (1 cuillère $\approx$ \SI{10}{g}) \\
\hline
Rapporteur et feuille blanche & Mesure de l'angle de rotation & Rapporteur d'écolier \\
\hline
Boîte de modèles moléculaires (boules + tiges) & Construction des stéréo-isomères & Pâte à modeler de 4 couleurs + cure-dents \\
\hline
Petit miroir de poche & Vérification image miroir & Miroir de sac, écran de téléphone éteint \\
\hline
\end{tabularx}
\end{center}

\courssub{Sécurité}

\begin{attention}[Consignes de sécurité]
Ce TP est peu dangereux, mais la rigueur du chimiste s'acquiert dès les manipulations simples.
\begin{itemize}
\item \textbf{Pictogramme \og verrerie \fg{} (risque de coupure)} : le cristallisoir et l'agitateur sont en verre. Manipule-les à deux mains, loin du bord de la paillasse. En cas de casse, ne ramasse jamais les éclats à mains nues : utilise balayette et pelle.
\item \textbf{Pictogramme \og danger électrique \fg{}} : si tu utilises une lampe branchée sur secteur, garde les mains sèches et éloigne tout récipient rempli de liquide des prises.
\item \textbf{Hygiène} : le sucre et le miel utilisés au laboratoire ne doivent \textbf{jamais} être goûtés, même s'ils sont comestibles à l'origine. C'est une règle absolue : aucun produit ne se porte à la bouche dans un laboratoire.
\item La pâte à modeler ne se mange pas non plus ; lave-toi les mains en fin de séance.
\item Ne regarde jamais directement une source lumineuse puissante à travers les polariseurs.
\end{itemize}
\end{attention}

\courssub{Schéma du dispositif}

\begin{popfigure}
\begin{center}
\begin{tikzpicture}[scale=1.0, >=Stealth]
% source lumineuse
\draw[fill=popgoldL, draw=popdark] (-5.4,0.4) rectangle (-4.6,-0.4);
\node[below, align=center, font=\small] at (-5,-0.5) {Lampe\\(lumière naturelle)};
\draw[popgold, thick] (-4.6,0) -- (-3.7,0);
\draw[popgold, thick] (-4.6,0.15) -- (-3.7,0.05);
\draw[popgold, thick] (-4.6,-0.15) -- (-3.7,-0.05);
% polariseur
\draw[fill=popblue!50, draw=popdark] (-3.6,-1.1) rectangle (-3.3,1.1);
\node[above, font=\small] at (-3.45,1.15) {Polariseur};
\draw[->, popblue, thick] (-3.0,0) -- (-2.2,0) node[midway, above, font=\small]{lumière polarisée};
% cuve
\draw[fill=popgreenL, draw=popdark, thick] (-1.6,-1.0) -- (-1.6,1.0) -- (1.6,1.0) -- (1.6,-1.0);
\draw[popdark, thick] (-1.6,-1.0) -- (1.6,-1.0);
\draw[popgreen!60!black] (-1.55,0.55) -- (1.55,0.55);
\node[font=\small, align=center] at (0,0) {Solution de saccharose\\(cristallisoir)};
\node[below, font=\small] at (0,-1.05) {Cuve de longueur $\ell$};
% spirale (rotation)
\draw[->, poppurple, thick, domain=0:300, samples=60] plot ({-2.6+0.004*(\x+300)*0.5}, {0.35*cos(\x)});
% analyseur
\draw[fill=poporange!60, draw=popdark] (2.3,-1.1) rectangle (2.6,1.1);
\node[above, font=\small] at (2.45,1.15) {Analyseur};
\draw[->, poporange, thick] (2.8,0) -- (3.7,0);
% oeil
\draw[fill=white, draw=popdark] (4.1,0) ellipse (0.35 and 0.22);
\fill[popdark] (4.15,0) circle (0.07);
\node[below, font=\small] at (4.1,-0.3) {Œil};
% angle
\draw[<->, poppurple, thick] (2.45,-1.35) arc[start angle=-90, end angle=-35, radius=1.0];
\node[below right, font=\small, poppurple] at (2.7,-1.5) {angle $\alpha$};
\end{tikzpicture}
\end{center}
\legende{Dispositif de polarimétrie simplifiée : la lumière traversant le polariseur devient polarisée rectilignement ; la solution de saccharose fait tourner son plan de polarisation d'un angle $\alpha$ ; on le mesure en tournant l'analyseur jusqu'au retour de l'extinction.}
\end{popfigure}

\courssub{Protocole}

\textbf{Partie 1 : Mise en évidence de l'activité optique du saccharose}

\begin{enumerate}
\item Prépare la solution de saccharose : pèse environ \SI{150}{g} de sucre en poudre (ou 15 cuillères à café rases) et dissous-les dans \SI{100}{mL} d'eau tiède en agitant jusqu'à dissolution complète. Laisse refroidir. Tu obtiens une solution très concentrée (masse volumique $\rho \approx \SI{1,5}{g.mL^{-1}}$, soit $c \approx \SI{4,4}{mol.L^{-1}}$).
\item Pose la feuille blanche sur la paillasse et installe la lampe de façon à éclairer horizontalement.
\item Place le premier filtre polarisant (le \cle{polariseur}) devant la lampe, puis le second (l'\cle{analyseur}) à environ \SI{20}{cm} plus loin. Regarde à travers l'analyseur.
\item Tourne lentement l'analyseur jusqu'à obtenir l'\cle{extinction} : le champ devient noir (ou très sombre). Les plans de polarisation des deux filtres sont alors perpendiculaires. Note la position angulaire de repère de l'analyseur : $\theta_1$.
\item Sans déplacer les filtres, interpose le cristallisoir contenant la solution de saccharose entre les deux filtres. Observe : la lumière \textbf{réapparaît}.
\item Tourne à nouveau l'analyseur jusqu'à retrouver l'extinction. Note la nouvelle position $\theta_2$ et le \textbf{sens} dans lequel tu as dû tourner (sens horaire ou antihoraire, vu du côté de l'observateur).
\item Mesure l'angle de rotation $\alpha = |\theta_2 - \theta_1|$ à l'aide du rapporteur.
\item Recommence avec une solution deux fois moins concentrée (dilue la moitié de la solution avec un volume égal d'eau), puis avec une cuve deux fois moins large si possible. Note les angles obtenus.
\end{enumerate}

\textbf{Partie 2 : Construction de modèles moléculaires chiraux}

\begin{enumerate}
\item Construis une molécule d'\cle{acide lactique} \ce{CH3-CH(OH)-COOH} : une boule noire pour le carbone central (asymétrique), reliée à quatre groupes différents : \ce{H} (blanc), \ce{OH} (rouge), \ce{CH3} (gris), \ce{COOH} (vert). Avec la pâte à modeler : quatre couleurs différentes et des cure-dents pour les liaisons.
\item Construis une seconde molécule, \textbf{image de la première dans le miroir} : place le miroir à côté du premier modèle et reproduis exactement ce que tu y vois.
\item Essaie de \textbf{superposer} les deux modèles : fais coïncider les groupes \ce{COOH} et \ce{CH3}, puis regarde si \ce{H} et \ce{OH} coïncident aussi. Conclus.
\item Vérifie que le carbone central porte bien quatre groupes \textbf{tous différents} : c'est un \cle{carbone asymétrique}, noté $C^*$.
\item Construis maintenant l'éthane \ce{C2H6} avec deux boules noires et six blanches. Fais tourner les deux groupes \ce{CH3} l'un par rapport à l'autre autour de la liaison \ce{C-C} : réalise la conformation \cle{éclipsée} (les hydrogènes se font face) puis la conformation \cle{décalée} (les hydrogènes sont intercalés, rotation de \SI{60}{\degree}).
\item Compare tes modèles aux représentations de Newman dessinées au tableau et fais correspondre chaque modèle à sa projection.
\end{enumerate}

\courssub{Observations et résultats attendus}

\textbf{Partie 1 : polarimétrie}

\begin{center}
\rowcolors{2}{popgreenL}{white}
\begin{tabularx}{\linewidth}{|X|c|c|c|}
\hline
\rowcolor{popgreen!40} \textbf{Expérience} & \textbf{Extinction sans solution ?} & \textbf{Lumière avec solution ?} & \textbf{Angle $\alpha$ mesuré} \\
\hline
Polariseurs croisés, cuve d'eau pure & Oui & Non (extinction conservée) & \SI{0}{\degree} \\
\hline
Solution de saccharose concentrée (\SI{1,5}{g.mL^{-1}}), cuve de \SI{10}{cm} & Oui & Oui, la lumière réapparaît & $\approx$ \SI{100}{\degree} (sens horaire) \\
\hline
Solution diluée 2 fois, même cuve & Oui & Oui, plus faiblement & $\approx$ \SI{50}{\degree} \\
\hline
Solution concentrée, cuve de \SI{5}{cm} & Oui & Oui & $\approx$ \SI{50}{\degree} \\
\hline
\end{tabularx}
\end{center}

\begin{itemize}
\item L'eau pure ne rétablit pas la lumière : elle est \cle{optiquement inactive}.
\item La solution de saccharose fait réapparaître la lumière : elle possède un \cle{pouvoir rotatoire}. Pour retrouver l'extinction, on doit tourner l'analyseur dans le \textbf{sens horaire} (vu de l'observateur) : le saccharose est \cle{dextrogyre} (notation $(+)$). \textbf{Attention} : le terme \og dextrose \fg{} désigne historiquement le \textbf{glucose} (D-glucose), pas le saccharose, bien que les deux soient déxtrogyres.
\item L'angle $\alpha$ est \textbf{proportionnel} à la concentration et à la longueur de la cuve : c'est la loi de Biot, $\alpha = [\alpha] \times \ell \times c$ (avec $[\alpha]$ le pouvoir rotatoire massique, $\ell$ en dm, $c$ en \SI{}{g.mL^{-1}}).
\end{itemize}

\textbf{Partie 2 : modèles moléculaires}

\begin{center}
\rowcolors{2}{poporangeL}{white}
\begin{tabularx}{\linewidth}{|X|X|}
\hline
\rowcolor{poporange!40} \textbf{Manipulation} & \textbf{Observation attendue} \\
\hline
Superposition des deux acides lactiques & Impossible : quand \ce{COOH} et \ce{CH3} coïncident, \ce{H} et \ce{OH} sont inversés. Les deux modèles sont \textbf{non superposables} : ce sont des énantiomères. \\
\hline
Modèle dans le miroir & L'image du modèle A est identique au modèle B : les deux énantiomères sont images l'un de l'autre dans un miroir. \\
\hline
Rotation autour de \ce{C-C} de l'éthane & La rotation est libre : éclipsée et décalée se transforment l'une en l'autre sans casser de liaison. Ce ne sont \textbf{pas} des isomères différents mais des conformations (conformères). \\
\hline
\end{tabularx}
\end{center}

\courssub{Exploitation}

\begin{exoresolu}[Questions guidées d'exploitation]
\textbf{1.} Pourquoi la lumière réapparaît-elle quand on place la solution de saccharose entre les polariseurs croisés ?

\textbf{2.} Le saccharose est déxtrogyre. Que signifie ce terme et quelle notation lui associe-t-on ?

\textbf{3.} D'après tes mesures et la loi de Biot, que devient l'angle de rotation si l'on double la concentration de la solution ? Et si l'on double la longueur de la cuve ?

\textbf{4.} Pourquoi les deux modèles d'acide lactique ne sont-ils pas superposables ? Quelle caractéristique structurale de la molécule explique cette propriété ?

\textbf{5.} Les conformations éclipsée et décalée de l'éthane sont-elles des énantiomères ? Justifie ta réponse en précisant la nature de l'isomérie concernée.
\tcblower
\textbf{1.} Entre polariseurs croisés, les axes de polarisation sont perpendiculaires : aucune lumière ne passe (extinction). La solution de saccharose, optiquement active, fait \textbf{tourner le plan de polarisation} de la lumière polarisée d'un angle $\alpha$. Le plan de vibration de la lumière sortante n'est donc plus perpendiculaire à l'axe de l'analyseur : une composante de l'amplitude électrique le traverse et le champ s'éclaire.

\textbf{2.} \emph{Dextrogyre} (noté $(+)$) signifie que la substance fait tourner le plan de polarisation vers la \textbf{droite} (sens horaire) pour un observateur qui regarde la lumière venir vers lui. Le saccharose est donc le $(+)$-saccharose. \emph{Précision} : l'ancien nom \og dextrose \fg{} correspond au D-glucose, à ne pas confondre avec le saccharose (sucre de table).

\textbf{3.} D'après la loi de Biot $\alpha = [\alpha]\,\ell\,c$, l'angle est proportionnel à la concentration $c$ et à la longueur $\ell$ : si $c$ double, $\alpha$ double ; si $\ell$ double, $\alpha$ double. Nos mesures le confirment : \SI{100}{\degree} pour la solution concentrée dans la cuve de \SI{10}{cm}, \SI{50}{\degree} pour la solution diluée deux fois (même $\ell$), et \SI{50}{\degree} pour la solution concentrée dans une cuve de \SI{5}{cm} (même $c$, $\ell$ divisé par 2).

\textbf{4.} Le carbone central de l'acide lactique (le $C_2$) porte quatre groupes différents (\ce{H}, \ce{OH}, \ce{CH3}, \ce{COOH}) : c'est un \textbf{carbone asymétrique} ($C^*$). La molécule est donc \textbf{chirale} : elle n'a pas de plan de symétrie ni de centre d'inversion. Elle existe sous deux formes images dans un miroir et non superposables, appelées \textbf{énantiomères}, exactement comme nos deux mains (chiralité = \og main \fg{} en grec).

\textbf{5.} Non. Les conformations éclipsée et décalée de l'éthane s'interconvertissent par simple \textbf{rotation autour de la liaison simple} \ce{C-C}, sans rupture de liaison. Elles ne sont pas images miroir non superposables (l'éthane n'a pas de carbone asymétrique). Ce sont des \textbf{conformères} (stéréo-isomères de conformation), distincts des énantiomères qui sont des stéréo-isomères de configuration stables à température ambiante.
\end{exoresolu}

\courssub{Conclusion}

Ce TP t'a permis de relier deux facettes complémentaires de la stéréochimie. D'une part, la \textbf{polarimétrie} montre qu'une solution de saccharose dévie le plan de la lumière polarisée : cette \cle{activité optique} est la signature macroscopique de la \cle{chiralité} des molécules. L'angle de rotation dépend de la nature de la substance (pouvoir rotatoire propre $[\alpha]$), de la concentration et de la longueur traversée, conformément à la loi de Biot : c'est le principe du \emph{saccharimètre}, utilisé dans les sucreries (SOSUCAM) et pour contrôler la qualité du miel sur les marchés. D'autre part, la manipulation des \textbf{modèles moléculaires} a rendu concrète la notion d'énantiomérie : deux molécules contenant un carbone asymétrique sont images dans un miroir et rigoureusement \textbf{non superposables}. Enfin, tu as distingué les énantiomères (isomères de configuration, stables, non interconvertibles sans rupture de liaison) des conformations (éclipsée, décalée), qui s'interconvertissent par simple rotation autour d'une liaison simple. Retiens l'idée directrice : \emph{la structure tridimensionnelle d'une molécule détermine ses propriétés}, qu'il s'agisse de dévier la lumière ou, comme tu le verras en biochimie et pharmacologie, d'agir comme médicament ou substrat enzymatique.

\newpage

\coursec{Méthodes et exercices résolus}

\courssub{Méthode 1 : Dénombrer et classer les isomères de constitution}

\begin{methode}[Trouver et classer les isomères de constitution]
Pour trouver tous les isomères de constitution d'une formule brute donnée (ex. \ce{C4H10O}) :
\begin{enumerate}
\item \textbf{Vérifier la formule brute} : compter le nombre d'insaturations si nécessaire (cycles, doubles liaisons).
\item \textbf{Varier la chaîne carbonée} : écrire d'abord la chaîne linéaire, puis raccourcir la chaîne principale en déplaçant des ramifications (isomérie de \cle{chaîne}).
\item \textbf{Déplacer le groupe caractéristique} ou l'insaturation le long de la chaîne (isomérie de \cle{position}).
\item \textbf{Changer de fonction} : pour une même formule brute, plusieurs familles sont possibles (ex. \ce{C2H6O} : alcool ou éther ; \ce{C3H6O} : aldéhyde ou cétone) — isomérie de \cle{fonction}.
\item \textbf{Nommer chaque composé} obtenu pour vérifier qu'il n'y a pas de doublon : deux structures qui portent le même nom sont identiques !
\item Conclure en indiquant, pour chaque couple d'isomères, le \textbf{type d'isomérie} (chaîne, position ou fonction).
\end{enumerate}
\textbf{Réflexe Bac} : présente les résultats dans un tableau (formule semi-développée + nom + type d'isomérie).
\end{methode}

\begin{exoresolu}[Isomères de formule brute \ce{C4H10O}]
\textbf{Énoncé.} Un composé organique $A$, présent dans une boisson artisanale analysée au laboratoire, a pour formule brute \ce{C4H10O}.
\begin{enumerate}
\item Donner les formules semi-développées et les noms de tous les isomères de constitution de $A$ qui sont des alcools.
\item Donner un exemple d'isomère de fonction de ces alcools.
\item Préciser le type d'isomérie qui lie le butan-1-ol au butan-2-ol, puis le butan-1-ol au 2-méthylpropan-1-ol.
\end{enumerate}
\tcblower
\textbf{Solution détaillée.}

\textbf{1. Alcools de formule \ce{C4H10O}.} On applique la méthode : on fait varier la chaîne (4 C linéaires, puis 3 C + ramification) et la position du groupe $-OH$.

\begin{center}
\begin{tabularx}{\linewidth}{|l|X|l|}
\hline
\rowcolor{popblueL} \textbf{N°} & \textbf{Formule semi-développée} & \textbf{Nom} \\
\hline
1 & \ce{CH3-CH2-CH2-CH2-OH} & butan-1-ol \\
\hline
2 & \ce{CH3-CH2-CH(OH)-CH3} & butan-2-ol \\
\hline
3 & \ce{(CH3)2CH-CH2-OH} & 2-méthylpropan-1-ol \\
\hline
4 & \ce{(CH3)3C-OH} & 2-méthylpropan-2-ol \\
\hline
\end{tabularx}
\end{center}

Il existe donc \textbf{4 alcools} isomères de formule \ce{C4H10O}.

\textbf{2. Isomère de fonction.} La formule \ce{C4H10O} correspond aussi à des \textbf{éthers}, par exemple :
\[ \ce{CH3-CH2-O-CH2-CH3} \quad \text{éthoxyéthane (éther diéthylique)} \]
ou encore \ce{CH3-O-CH2-CH2-CH3} (méthoxypropane).

\textbf{3. Types d'isomérie.}
\begin{itemize}
\item Butan-1-ol \ce{CH3-CH2-CH2-CH2-OH} et butan-2-ol \ce{CH3-CH2-CH(OH)-CH3} : même chaîne carbonée, même fonction alcool, mais le groupe $-OH$ n'occupe pas la même position $\Rightarrow$ \textbf{isomérie de position}.
\item Butan-1-ol et 2-méthylpropan-1-ol \ce{(CH3)2CH-CH2-OH} : le groupe $-OH$ est en position 1 dans les deux cas, mais les squelettes carbonés diffèrent (chaîne linéaire / chaîne ramifiée) $\Rightarrow$ \textbf{isomérie de chaîne}.
\end{itemize}

\textbf{Conclusion.} \ce{C4H10O} compte 4 isomères alcools (chaîne et position) et admet des isomères de fonction de la famille des éthers.
\end{exoresolu}

\courssub{Méthode 2 : Représenter une molécule en perspective (représentation de Cram)}

\begin{methode}[Construire la représentation de Cram]
Pour représenter en perspective une molécule possédant un \cle{carbone asymétrique} (carbone tétraédrique portant 4 groupes différents) :
\begin{enumerate}
\item \textbf{Identifier le carbone asymétrique} : repérer le carbone lié à 4 atomes ou groupes \textbf{tous différents} ; le marquer d'un astérisque $C^*$.
\item Placer le carbone asymétrique au centre ; dessiner \textbf{deux liaisons dans le plan} de la feuille (traits pleins fins).
\item Dessiner \textbf{une liaison en avant du plan} : trait plein épais (triangle plein) pointant vers l'observateur.
\item Dessiner \textbf{une liaison en arrière du plan} : trait hachuré (triangle pointillé) s'éloignant de l'observateur.
\item Répartir les 4 substituants sur ces 4 liaisons.
\item \textbf{Vérification} : l'image dans un miroir ne doit pas être superposable à la molécule initiale ; si deux groupes identiques sont échangés, on obtient l'énantiomère.
\end{enumerate}
\textbf{Réflexes} : l'échange de \textbf{deux} substituants transforme l'énantiomère en son image ; l'échange de \textbf{trois} substituants (permutation circulaire) laisse la configuration inchangée.
\end{methode}

\begin{exoresolu}[Carbone asymétrique et représentation de Cram de l'acide lactique]
\textbf{Énoncé.} L'acide lactique, présent dans le lait caillé et dans les muscles après un effort, a pour formule semi-développée \ce{CH3-CH(OH)-COOH}.
\begin{enumerate}
\item Justifier l'existence d'un carbone asymétrique et le repérer.
\item Représenter en perspective (Cram) les deux stéréo-isomères de cette molécule.
\item Comment nomme-t-on ces deux stéréo-isomères ? Quelle relation les lie ?
\end{enumerate}
\tcblower
\textbf{Solution détaillée.}

\textbf{1. Carbone asymétrique.} Le carbone porteur du groupe $-OH$ est lié à quatre groupes tous différents :
\begin{itemize}
\item un atome d'hydrogène \ce{-H} ;
\item un groupe hydroxyle \ce{-OH} ;
\item un groupe méthyle \ce{-CH3} ;
\item un groupe carboxyle \ce{-COOH}.
\end{itemize}
C'est donc un \textbf{carbone asymétrique}, noté $C^*$ : \ce{CH3-\overset{*}{C}H(OH)-COOH}. La molécule est \textbf{chirale} : elle n'est pas superposable à son image dans un miroir.

\textbf{2. Représentations de Cram des deux énantiomères.}

\begin{popfigure}
\begin{center}
\begin{tikzpicture}[scale=1.0, baseline=(current bounding box.center)]
% Enantiomère 1 : OH en avant (plein), COOH en arrière (hachuré)
\node at (0,0) {\chemfig{CH_3-[:30]C(-[:90]H)(<[:150]OH)(<:[-90]COOH)}};
\node at (0,-2.5) {\small énantiomère (a)};
% Miroir
\draw[popmuted, dashed, thick] (3.5,-2.5) -- (3.5,1.5);
\node[popmuted] at (3.5,1.8) {\small miroir};
% Enantiomère 2 : COOH en avant (plein), OH en arrière (hachuré)
\node at (7,0) {\chemfig{CH_3-[:30]C(-[:90]H)(<[:150]COOH)(<:[-90]OH)}};
\node at (7,-2.5) {\small énantiomère (b)};
\end{tikzpicture}
\end{center}
\legende{Les deux énantiomères de l'acide lactique en représentation de Cram : le trait plein épais (coin) pointe vers l'avant, le trait hachuré vers l'arrière.}
\end{popfigure}

Dans (a), le groupe \ce{-OH} est en avant du plan (liaison pleine) et \ce{-COOH} en arrière (liaison hachurée) ; dans (b), on a échangé ces deux groupes : on obtient l'image miroir.

\textbf{3. Relation entre les deux stéréo-isomères.} Ces deux molécules sont \textbf{images l'une de l'autre dans un miroir et non superposables} : ce sont des \textbf{énantiomères}. Elles ont les mêmes propriétés physiques (température d'ébullition, densité...) et chimiques, sauf leur action sur la lumière polarisée et leur comportement vis-à-vis des réactifs chiraux (enzymes).

\textbf{Conclusion.} L'acide lactique possède un carbone asymétrique, donc deux énantiomères, représentés en Cram en inversant les liaisons avant/arrière.
\end{exoresolu}

\courssub{Méthode 3 : Représenter les énantiomères en projection de Fischer}

\begin{methode}[Construire et lire une projection de Fischer]
La projection de \cle{Fischer} est une représentation plane conventionnelle d'une molécule chirale :
\begin{enumerate}
\item Placer le \textbf{carbone asymétrique à l'intersection} d'une croix (il n'est pas écrit).
\item Disposer la \textbf{chaîne carbonée principale verticalement}, le groupe le plus oxydé (ex. \ce{-COOH}, \ce{-CHO}) \textbf{en haut}.
\item Convention : les liaisons \textbf{verticales} vont \textbf{vers l'arrière} du plan ; les liaisons \textbf{horizontales} viennent \textbf{vers l'avant}.
\item Placer les deux autres groupes (\ce{-H}, \ce{-OH}, \ce{-NH2}...) sur la ligne horizontale.
\item Pour obtenir l'\textbf{énantiomère} : échanger les deux groupes horizontaux (ou effectuer une symétrie miroir gauche/droite).
\end{enumerate}
\textbf{Interdictions} : ne jamais tourner la projection de \SI{90}{\degree} dans le plan (cela change la configuration !) ; seule une rotation de \SI{180}{\degree} dans le plan laisse la molécule identique.
\end{methode}

\begin{exoresolu}[Projection de Fischer de la glycéraldéhyde et de l'alanine]
\textbf{Énoncé.} 
\begin{enumerate}
\item Le glycéraldéhyde \ce{HOCH2-CH(OH)-CHO} est le plus simple des sucres chiraux. Représenter ses deux énantiomères en projection de Fischer.
\item L'alanine \ce{CH3-CH(NH2)-COOH} est un acide aminé. Représenter en Fischer l'énantiomère dont le groupe \ce{-NH2} est à gauche, puis son image miroir.
\item Rappeler la convention spatiale associée à la projection de Fischer.
\end{enumerate}
\tcblower
\textbf{Solution détaillée.}

\textbf{1. Énantiomères du glycéraldéhyde.} Le carbone central est lié à \ce{-H}, \ce{-OH}, \ce{-CHO} et \ce{-CH2OH} : il est asymétrique. On place la chaîne verticale avec le groupe le plus oxydé \ce{-CHO} en haut :

\begin{popfigure}
\begin{center}
\begin{tikzpicture}[scale=0.9]
% D-glycéraldéhyde (OH à droite)
\draw[thick] (0,1.0) -- (0,-1.0); % vertical
\draw[thick] (-0.7,0) -- (0.7,0); % horizontal
\node at (0,1.5) {\ce{CHO}};
\node at (-1.2,0) {\ce{H}};
\node at (1.2,0) {\ce{OH}};
\node at (0,-1.5) {\ce{CH2OH}};
\node at (0,-2.3) {\small D-glycéraldéhyde};
% Miroir
\draw[popmuted, dashed, thick] (2.5,-2) -- (2.5,2);
\node[popmuted] at (2.5,2.3) {\small miroir};
% L-glycéraldéhyde (OH à gauche)
\draw[thick] (5,1.0) -- (5,-1.0);
\draw[thick] (4.3,0) -- (5.7,0);
\node at (5,1.5) {\ce{CHO}};
\node at (3.8,0) {\ce{HO}};
\node at (6.2,0) {\ce{H}};
\node at (5,-1.5) {\ce{CH2OH}};
\node at (5,-2.3) {\small L-glycéraldéhyde};
\end{tikzpicture}
\end{center}
\legende{Projections de Fischer des deux énantiomères du glycéraldéhyde.}
\end{popfigure}

\textbf{2. Énantiomères de l'alanine.} Le carbone central est lié à \ce{-H}, \ce{-NH2}, \ce{-CH3}, \ce{-COOH} : il est asymétrique. Chaîne verticale, \ce{-COOH} en haut, \ce{-CH3} en bas :

\begin{popfigure}
\begin{center}
\begin{tikzpicture}[scale=0.9]
% L-alanine (NH2 à gauche)
\draw[thick] (0,1.0) -- (0,-1.0);
\draw[thick] (-0.7,0) -- (0.7,0);
\node at (0,1.5) {\ce{COOH}};
\node at (-1.2,0) {\ce{H2N}};
\node at (1.2,0) {\ce{H}};
\node at (0,-1.5) {\ce{CH3}};
\node at (0,-2.3) {\small L-alanine (\ce{NH2} à gauche)};
% Miroir
\draw[popmuted, dashed, thick] (2.5,-2) -- (2.5,2);
\node[popmuted] at (2.5,2.3) {\small miroir};
% D-alanine (NH2 à droite)
\draw[thick] (5,1.0) -- (5,-1.0);
\draw[thick] (4.3,0) -- (5.7,0);
\node at (5,1.5) {\ce{COOH}};
\node at (3.8,0) {\ce{H}};
\node at (6.2,0) {\ce{NH2}};
\node at (5,-1.5) {\ce{CH3}};
\node at (5,-2.3) {\small D-alanine (image miroir)};
\end{tikzpicture}
\end{center}
\legende{Projections de Fischer des deux énantiomères de l'alanine.}
\end{popfigure}

\textbf{3. Convention spatiale.} Dans une projection de Fischer :
\begin{itemize}
\item les liaisons \textbf{horizontales} (\ce{-H}, \ce{-NH2} ici) sortent du plan, \textbf{vers l'observateur} ;
\item les liaisons \textbf{verticales} (\ce{-COOH}, \ce{-CH3}) s'enfoncent \textbf{derrière le plan}.
\end{itemize}

\textbf{Conclusion.} La projection de Fischer permet de représenter rapidement les énantiomères : passer de l'un à l'autre revient à échanger les deux groupes horizontaux.
\end{exoresolu}

\courssub{Méthode 4 : Distinguer énantiomères et diastéréoisomères}

\begin{methode}[Comparer deux stéréo-isomères]
Pour deux stéréo-isomères donnés :
\begin{enumerate}
\item \textbf{Repérer tous les éléments de stéréochimie} : carbones asymétriques $C^*$ et doubles liaisons \ce{C=C} donnant Z/E.
\item \textbf{Comparer chaque élément} entre les deux molécules : la configuration est-elle identique ou inversée ?
\item Conclure :
\begin{itemize}
\item si \textbf{toutes} les configurations sont inversées $\Rightarrow$ \textbf{énantiomères} (images miroir non superposables) ;
\item si \textbf{certaines seulement} sont inversées (et au moins une identique) $\Rightarrow$ \textbf{diastéréoisomères} ;
\item si toutes sont identiques $\Rightarrow$ \textbf{même molécule}.
\end{itemize}
\item Cas particulier : une molécule possédant des $C^*$ mais un \textbf{plan de symétrie interne} (ou centre d'inversion) est \textbf{achirale} (composé \emph{méso}) : elle est superposable à son image et optiquement inactive.
\end{enumerate}
\textbf{Réflexe} : pour $n$ carbones asymétriques, il existe au maximum $2^n$ stéréo-isomères (moins si composés \emph{méso}).
\end{methode}

\begin{exoresolu}[Stéréo-isomères de l'acide tartrique]
\textbf{Énoncé.} L'acide tartrique, présent dans le raisin et le vin, a pour formule \ce{HOOC-CH(OH)-CH(OH)-COOH}.
\begin{enumerate}
\item Combien de carbones asymétriques possède cette molécule ? Combien de stéréo-isomères peut-on prévoir a priori ?
\item Représenter en projection de Fischer les différents stéréo-isomères.
\item Montrer que l'un d'eux est achiral (composé \emph{méso}). Conclure sur le nombre réel de stéréo-isomères et préciser les relations (énantiomères / diastéréoisomères) entre eux.
\end{enumerate}
\tcblower
\textbf{Solution détaillée.}

\textbf{1. Carbones asymétriques.} Les deux carbones centraux sont chacun liés à \ce{-H}, \ce{-OH}, \ce{-COOH} et \ce{-CH(OH)COOH} : ce sont \textbf{deux carbones asymétriques}. A priori : $2^2 = 4$ stéréo-isomères.

\textbf{2. Projections de Fischer.} Chaîne verticale, les deux \ce{-COOH} aux extrémités :

\begin{popfigure}
\begin{center}
\begin{tikzpicture}[scale=0.85]
% (I) : OH à droite sur les deux carbones
\draw[thick] (0,2.0) -- (0,-2.0);
\draw[thick] (-0.7,0.7) -- (0.7,0.7);
\draw[thick] (-0.7,-0.7) -- (0.7,-0.7);
\node at (0,2.5) {\ce{COOH}}; \node at (0,-2.5) {\ce{COOH}};
\node at (-1.2,0.7) {\ce{H}}; \node at (1.2,0.7) {\ce{OH}};
\node at (-1.2,-0.7) {\ce{H}}; \node at (1.2,-0.7) {\ce{OH}};
\node at (0,-3.2) {\small (I)};
% (II) : OH à gauche sur les deux carbones (image miroir de I)
\draw[thick] (4,2.0) -- (4,-2.0);
\draw[thick] (3.3,0.7) -- (4.7,0.7);
\draw[thick] (3.3,-0.7) -- (4.7,-0.7);
\node at (4,2.5) {\ce{COOH}}; \node at (4,-2.5) {\ce{COOH}};
\node at (2.8,0.7) {\ce{HO}}; \node at (5.2,0.7) {\ce{H}};
\node at (2.8,-0.7) {\ce{HO}}; \node at (5.2,-0.7) {\ce{H}};
\node at (4,-3.2) {\small (II)};
% (III) Méso : OH à droite en haut, à gauche en bas (symétrie)
\draw[thick] (8,2.0) -- (8,-2.0);
\draw[thick] (7.3,0.7) -- (8.7,0.7);
\draw[thick] (7.3,-0.7) -- (8.7,-0.7);
\node at (8,2.5) {\ce{COOH}}; \node at (8,-2.5) {\ce{COOH}};
\node at (6.8,0.7) {\ce{H}}; \node at (9.2,0.7) {\ce{OH}};
\node at (6.8,-0.7) {\ce{HO}}; \node at (9.2,-0.7) {\ce{H}};
\node at (8,-3.2) {\small (III) \emph{méso}};
% Plan de symétrie horizontal pour le méso
\draw[popmuted, dashed, thick] (7.0,0) -- (9.0,0);
\node[popmuted] at (8,0.4) {\small plan de symétrie};
\end{tikzpicture}
\end{center}
\legende{Les trois stéréo-isomères de l'acide tartrique en projection de Fischer.}
\end{popfigure}

\textbf{3. Analyse.}
\begin{itemize}
\item (I) et (II) sont images miroir l'un de l'autre, non superposables, et \textbf{les deux} configurations sont inversées $\Rightarrow$ ce sont des \textbf{énantiomères}.
\item (III) possède un \textbf{plan de symétrie} horizontal passant entre les deux carbones (représenté en pointillés) : la moitié haute est l'image de la moitié basse. La molécule est donc \textbf{superposable à son image miroir} : elle est \textbf{achirale}, c'est le composé \emph{méso}, optiquement \textbf{inactif}.
\item (III) a une configuration inversée et une identique par rapport à (I) ou (II) $\Rightarrow$ (III) est \textbf{diastéréoisomère} de (I) et de (II).
\end{itemize}

\textbf{Conclusion.} L'acide tartrique possède seulement \textbf{3 stéréo-isomères} (et non 4) : un couple d'énantiomères (I)/(II) et un composé \emph{méso} (III), diastéréoisomère des deux précédents.
\end{exoresolu}

\courssub{Méthode 5 : Attribuer la configuration Z/E d'un alcène}

\begin{methode}[Déterminer la configuration Z ou E]
\begin{enumerate}
\item \textbf{Vérifier l'existence de l'isomérie} : chaque carbone de la double liaison \ce{C=C} doit porter \textbf{deux substituants différents}. Si un carbone porte deux groupes identiques (ex. \ce{CH2=}), il n'y a ni Z ni E.
\item Sur \textbf{chaque carbone} de la double liaison, identifier le \textbf{groupe prioritaire} : en Terminale TI, on retient que le groupe de plus grande masse / le groupe carboné est prioritaire sur \ce{H} (règle de Cahn-Ingold-Prelog simplifiée).
\item Observer la position relative des deux groupes prioritaires :
\begin{itemize}
\item du \textbf{même côté} de la double liaison $\Rightarrow$ configuration \cle{Z} (\emph{zusammen}, « ensemble ») ;
\item de \textbf{part et d'autre} $\Rightarrow$ configuration \cle{E} (\emph{entgegen}, « opposé »).
\end{itemize}
\item Écrire le nom complet : (Z)-but-2-ène ou (E)-but-2-ène.
\end{enumerate}
\textbf{Réflexe} : les isomères Z et E sont des \textbf{diastéréoisomères} (pas des énantiomères) : ils ont des propriétés physiques différentes (températures de fusion et d'ébullition distinctes).
\end{methode}

\begin{exoresolu}[Isomérie Z/E du but-2-ène et d'un alcène de la chimie des arômes]
\textbf{Énoncé.}
\begin{enumerate}
\item Le but-1-ène présente-t-il une isomérie Z/E ? Justifier.
\item Représenter les deux stéréo-isomères du but-2-ène \ce{CH3-CH=CH-CH3} et les nommer.
\item L'acide oléique, constituant majoritaire de l'huile de palme et de l'huile d'olive, possède une double liaison de configuration Z. Quelle conséquence cette configuration a-t-elle sur la forme de la chaîne carbonée et sur l'état physique de l'huile à température ambiante ?
\end{enumerate}
\tcblower
\textbf{Solution détaillée.}

\textbf{1. But-1-ène \ce{CH2=CH-CH2-CH3}.} Le premier carbone de la double liaison porte \textbf{deux atomes d'hydrogène identiques} (\ce{CH2=}). La condition d'existence n'est pas remplie : le but-1-ène \textbf{ne présente pas} d'isomérie Z/E.

\textbf{2. Stéréo-isomères du but-2-ène.} Chaque carbone de \ce{C=C} porte \ce{-H} et \ce{-CH3} (différents) : l'isomérie Z/E existe. Le groupe \ce{-CH3} est prioritaire sur \ce{H} sur chaque carbone.

\begin{popfigure}
\begin{center}
\begin{tikzpicture}[scale=1.0]
% Z : les deux CH3 du même côté (ici "en haut" par rapport au plan de la double liaison)
\node at (0,0) {\chemfig{CH_3-[:0]C(-[2]H)=C(-[2]H)-[:0]CH_3}};
\node at (0,-2.0) {\small (Z)-but-2-ène : groupes prioritaires (\ce{CH3}) du même côté};
% E : les deux CH3 de part et d'autre
\node at (7,0) {\chemfig{CH_3-[:0]C(-[2]H)=C(-[-2]H)-[:0]CH_3}};
\node at (7,-2.0) {\small (E)-but-2-ène : groupes prioritaires (\ce{CH3}) de part et d'autre};
\end{tikzpicture}
\end{center}
\legende{Les deux diastéréoisomères du but-2-ène.}
\end{popfigure}

\begin{itemize}
\item \textbf{(Z)-but-2-ène} : les deux groupes prioritaires \ce{-CH3} sont du même côté de la double liaison (ici tous deux "au-dessus" du plan de la double liaison, les H étant "en dessous").
\item \textbf{(E)-but-2-ène} : les deux groupes \ce{-CH3} sont de part et d'autre (l'un "au-dessus", l'autre "en dessous").
\end{itemize}

\textbf{3. Conséquence de la configuration Z dans l'acide oléique.} La configuration Z impose un \textbf{coude} (un « pli ») dans la chaîne carbonée au niveau de la double liaison. Les molécules ne peuvent alors pas s'empiler régulièrement : les forces d'attraction intermoléculaires (de Van der Waals) sont affaiblies, la température de fusion est \textbf{abaissée}. C'est pourquoi les huiles riches en acides gras Z (insaturés), comme l'huile de palme rouge ou l'huile d'arachide, sont \textbf{liquides} à température ambiante, alors que les graisses saturées (chaînes rectilignes) sont solides.

\textbf{Conclusion.} L'isomérie Z/E n'existe que si chaque carbone de \ce{C=C} porte deux groupes différents ; Z = groupes prioritaires du même côté, E = de part et d'autre. Cette géométrie a des conséquences concrètes, comme l'état liquide des huiles.
\end{exoresolu}

\courssub{Méthode 6 : Représenter les conformations (Newman, chaise, bateau)}

\begin{methode}[Construire une projection de Newman]
Pour représenter les \cle{conformations} d'une molécule autour d'une liaison simple \ce{C-C} (ex. éthane, butane) :
\begin{enumerate}
\item Placer l'œil \textbf{dans l'axe} de la liaison \ce{C-C} étudiée.
\item Le \textbf{carbone avant} est représenté par un \textbf{point} (centre du cercle) d'où partent 3 liaisons à \ang{120}.
\item Le \textbf{carbone arrière} est représenté par le \textbf{cercle} ; ses 3 liaisons partent du bord du cercle.
\item Faire \textbf{tourner} le carbone arrière autour de la liaison :
\begin{itemize}
\item \textbf{conformation éclipsée} : les liaisons avant et arrière se recouvrent (angle de \ang{0}) — énergie \textbf{maximale}, gêne stérique forte ;
\item \textbf{conformation décalée} : les liaisons arrière sont intercalées (angle de \ang{60}) — énergie \textbf{minimale}, conformation la plus stable.
\end{itemize}
\item Pour le \textbf{cyclohexane}, retenir les deux conformations : \textbf{chaise} (la plus stable, sans tension) et \textbf{bateau} (moins stable, interactions entre les « proues »).
\end{enumerate}
\textbf{Point clé} : les conformations se transforment les unes en autres par simple \textbf{rotation autour d'une liaison simple}, sans rompre de liaison : ce ne sont \textbf{pas} des isomères distincts isolables.
\end{methode}

\begin{exoresolu}[Conformations de l'éthane et du butane]
\textbf{Énoncé.}
\begin{enumerate}
\item Représenter en projection de Newman les conformations éclipsée et décalée de l'éthane \ce{CH3-CH3}, en regardant suivant l'axe \ce{C-C}.
\item Laquelle est la plus stable ? Justifier.
\item On considère maintenant le butane \ce{CH3-CH2-CH2-CH3} observé suivant la liaison \ce{C2-C3}. Représenter la conformation décalée dite « anti » (les deux groupes \ce{CH3} à \ang{180}) et la conformation éclipsée où les deux \ce{CH3} se superposent. Comparer leurs stabilités.
\item Citer les deux conformations caractéristiques du cyclohexane et préciser la plus stable.
\end{enumerate}
\tcblower
\textbf{Solution détaillée.}

\textbf{1. Projections de Newman de l'éthane.}

\begin{popfigure}
\begin{center}
\begin{tikzpicture}[scale=1.1]
% Décalée (staggered)
\draw[thick] (0,0) circle (1.0);
% Liens avant (centre)
\foreach \a in {90,210,330} {\draw[thick] (0,0) -- (\a:1.0);}
\node at (90:1.35) {\small H}; \node at (210:1.35) {\small H}; \node at (330:1.35) {\small H};
% Liens arrière (bord du cercle, décalés de 60°)
\foreach \a in {30,150,270} {\draw[thick, popblue] (\a:0.6) -- (\a:1.0);}
\node[popblue] at (30:1.35) {\small H}; \node[popblue] at (150:1.35) {\small H}; \node[popblue] at (270:1.35) {\small H};
\node at (0,-1.9) {\small décalée (stable)};
% Éclipsée
\begin{scope}[xshift=5.5cm]
\draw[thick] (0,0) circle (1.0);
\foreach \a in {90,210,330} {\draw[thick] (0,0) -- (\a:1.0);}
\node at (90:1.35) {\small H}; \node at (210:1.35) {\small H}; \node at (330:1.35) {\small H};
% Liens arrière presque alignés (légèrement décalés pour visibilité)
\foreach \a in {85,205,325} {\draw[thick, poporange] (\a:0.6) -- (\a:1.0);}
\node[poporange] at (85:1.35) {\small H}; \node[poporange] at (205:1.35) {\small H}; \node[poporange] at (325:1.35) {\small H};
\node at (0,-1.9) {\small éclipsée (instable)};
\end{scope}
\end{tikzpicture}
\end{center}
\legende{Projections de Newman de l'éthane : conformations décalée (liaisons décalées de \ang{60}) et éclipsée (liaisons alignées).}
\end{popfigure}

\textbf{2. Stabilité.} La conformation \textbf{décalée} est la plus stable : les liaisons \ce{C-H} (et leurs doublets électroniques) sont le plus \textbf{éloignées} possible, ce qui minimise les répulsions électroniques et la gêne stérique. La conformation éclipsée correspond à un maximum d'énergie (environ \SI{12}{kJ.mol^{-1}} au-dessus).

\textbf{3. Conformations du butane suivant \ce{C2-C3}.}

\begin{popfigure}
\begin{center}
\begin{tikzpicture}[scale=1.1]
% Anti (décalée, CH3 à 180°)
\draw[thick] (0,0) circle (1.0);
% Carbone avant (C2) : CH3 en haut (90°), H en bas à gauche (210°), H en bas à droite (330°)
\foreach \a in {90,210,330} {\draw[thick] (0,0) -- (\a:1.0);}
\node at (90:1.4) {\small \ce{CH3}}; \node at (210:1.3) {\small H}; \node at (330:1.3) {\small H};
% Carbone arrière (C3) : CH3 en bas (270°), H en haut à droite (30°), H en haut à gauche (150°)
\foreach \a in {30,150,270} {\draw[thick, popgreen] (\a:0.6) -- (\a:1.0);}
\node[popgreen] at (270:1.4) {\small \ce{CH3}}; \node[popgreen] at (30:1.3) {\small H}; \node[popgreen] at (150:1.3) {\small H};
\node at (0,-2.1) {\small anti (la plus stable)};
% Éclipsée CH3/CH3
\begin{scope}[xshift=6cm]
\draw[thick] (0,0) circle (1.0);
% Carbone avant : CH3 à 90°
\foreach \a in {90,210,330} {\draw[thick] (0,0) -- (\a:1.0);}
\node at (90:1.45) {\small \ce{CH3}}; \node at (210:1.3) {\small H}; \node at (330:1.3) {\small H};
% Carbone arrière : CH3 à ~90° (éclipsé), H aux autres positions
\foreach \a in {85,205,325} {\draw[thick, poporange] (\a:0.6) -- (\a:1.0);}
\node[poporange] at (85:1.45) {\small \ce{CH3}}; \node[poporange] at (205:1.3) {\small H}; \node[poporange] at (325:1.3) {\small H};
\node at (0,-2.1) {\small éclipsée \ce{CH3}/\ce{CH3} (la moins stable)};
\end{scope}
\end{tikzpicture}
\end{center}
\legende{Conformations extrêmes du butane vu suivant la liaison \ce{C2-C3}.}
\end{popfigure}

La conformation \textbf{anti} (les deux \ce{CH3} à \ang{180}) est la plus stable car les deux groupes volumineux sont au plus éloignés. La conformation \textbf{éclipsée} où les deux \ce{CH3} se superposent est la \textbf{moins stable} : gêne stérique maximale.

\textbf{4. Cyclohexane.} Les deux conformations caractéristiques sont la forme \textbf{chaise} et la forme \textbf{bateau}. La conformation \textbf{chaise} est la plus stable : tous les angles de valence sont proches de \ang{109,5} et toutes les liaisons \ce{C-H} sont décalées, ce qui supprime les tensions d'angle et les interactions éclipsées.

\textbf{Conclusion.} Les conformations résultent de la rotation libre autour des liaisons simples ; la décalée (anti) est toujours plus stable que l'éclipsée, et le cyclohexane adopte préférentiellement la forme chaise.
\end{exoresolu}

\courssub{Méthode 7 : Exploiter l'activité optique d'une substance}

\begin{methode}[Activité optique et pouvoir rotatoire]
\begin{enumerate}
\item \textbf{Principe} : une solution d'une substance \textbf{chirale} (un seul énantiomère, ou excès d'un énantiomère) fait \textbf{tourner le plan de polarisation} de la lumière polarisée qui la traverse : on dit qu'elle est \textbf{optiquement active}.
\item Si la déviation s'effectue vers la droite (sens horaire) : énantiomère \textbf{dextrogyre} (noté $+$ ou \emph{d}) ; vers la gauche (sens anti-horaire) : \textbf{lévogyre} (noté $-$ ou \emph{l}).
\item Un \textbf{mélange racémique} (50 \% de chaque énantiomère) est \textbf{globalement inactif} : les effets se compensent exactement. De même, un composé \emph{méso} est inactif.
\item \textbf{Application} : la polarimétrie permet de mesurer la concentration d'une solution de substance chirale (ex. dosage du sucre par saccharimétrie).
\end{enumerate}
\textbf{Réflexe Bac} : « optiquement actif » $\Rightarrow$ molécule chirale (ou mélange non racémique) ; « inactif » $\Rightarrow$ molécule achirale, composé \emph{méso} ou mélange racémique.
\end{methode}

\begin{exoresolu}[Analyse polarimétrique d'un jus de fruit fermenté]
\textbf{Énoncé.} Au laboratoire d'une usine de jus de la région de l'Ouest, on étudie un acide $B$ de formule \ce{CH3-CH(OH)-COOH} isolé lors d'une fermentation.
\begin{enumerate}
\item La molécule $B$ est-elle chirale ? Justifier.
\item Un échantillon $E_1$ de $B$, obtenu par synthèse chimique en laboratoire, ne dévie pas le plan de la lumière polarisée. Un échantillon $E_2$, extrait d'un ferment naturel, dévie le plan de polarisation vers la gauche. Interpréter ces deux observations.
\item On réalise un mélange équimolaire des deux énantiomères de $B$ ($n = \SI{0,10}{mol}$ de chacun). Quel est l'effet de ce mélange sur la lumière polarisée ? Comment appelle-t-on ce mélange ?
\end{enumerate}
\tcblower
\textbf{Solution détaillée.}

\textbf{1. Chiralité de $B$.} Le carbone central de \ce{CH3-CH(OH)-COOH} porte quatre groupes différents (\ce{-H}, \ce{-OH}, \ce{-CH3}, \ce{-COOH}) : c'est un \textbf{carbone asymétrique}. La molécule $B$ (acide lactique) est donc \textbf{chirale} et possède deux énantiomères.

\textbf{2. Interprétation des observations.}
\begin{itemize}
\item $E_1$ est \textbf{optiquement inactif} alors que la molécule est chirale : la synthèse chimique classique produit les deux énantiomères \textbf{en quantités égales}. $E_1$ est donc un \textbf{mélange racémique} : les pouvoirs rotatoires opposés des deux énantiomères se compensent.
\item $E_2$ dévie le plan de polarisation vers la \textbf{gauche} : il contient un \textbf{seul énantiomère} (ou un fort excès de celui-ci), l'acide lactique \textbf{lévogyre}. Cela s'explique par le fait que les fermentations biologiques font intervenir des \textbf{enzymes}, elles-mêmes chirales, qui ne produisent qu'un seul stéréo-isomère.
\end{itemize}

\textbf{3. Mélange équimolaire.} Le mélange contient $n = \SI{0,10}{mol}$ de chaque énantiomère, soit des fractions molaires égales à \SI{50}{\%}. Les déviations, égales en valeur absolue et de sens opposés, \textbf{s'annulent} : le mélange est \textbf{optiquement inactif}. On l'appelle \textbf{mélange racémique} (ou racémate).

\textbf{Conclusion.} L'activité optique est la signature expérimentale de la chiralité : un seul énantiomère dévie la lumière polarisée, le mélange racémique est inactif. C'est un moyen simple de distinguer une synthèse chimique (racémique) d'une synthèse enzymatique (un seul énantiomère).
\end{exoresolu}

\begin{attention}[Les 5 erreurs qui coûtent des points au Bac]
\begin{enumerate}
\item \textbf{Oublier de vérifier la condition d'existence de l'isomérie Z/E.} Si un carbone de la double liaison porte deux groupes identiques (ex. \ce{CH2=} dans le but-1-ène), il n'y a \textbf{pas} d'isomérie Z/E. Toujours le justifier avant de répondre.
\item \textbf{Confondre énantiomères et diastéréoisomères.} Énantiomères = images miroir non superposables avec \textbf{toutes} les configurations inversées. Si une seule configuration change sur plusieurs, ce sont des diastéréoisomères. Les isomères Z/E sont des diastéréoisomères, jamais des énantiomères !
\item \textbf{Mal tracer la représentation de Cram.} Le coin plein épais pointe \textbf{vers l'avant}, le coin hachuré \textbf{vers l'arrière}. Inverser les deux conventions, ou dessiner quatre liaisons dans le plan, rend la représentation fausse.
\item \textbf{Tourner une projection de Fischer de \ang{90}.} Seule une rotation de \ang{180} dans le plan conserve la configuration. De même, dans Fischer, les liaisons horizontales viennent \textbf{vers l'avant} et les verticales vont \textbf{vers l'arrière} : ne jamais l'oublier dans l'interprétation.
\item \textbf{Croire qu'un mélange racémique ou un composé \emph{méso} est optiquement actif.} Un mélange 50/50 des deux énantiomères et une molécule \emph{méso} (plan de symétrie interne) sont \textbf{optiquement inactifs}. Et attention : conformations (éclipsée, décalée, chaise, bateau) ne sont \textbf{pas} des isomères distincts : elles s'interconvertissent par simple rotation autour d'une liaison simple.
\end{enumerate}
\end{attention}

\newpage

\coursec{Exercices}

\courssub{Je vérifie mes connaissances}

\begin{exercice}[QCM : à toi de choisir]
Pour chaque question, une seule réponse est exacte. Recopie la lettre correspondante sur ta feuille.
\begin{enumerate}
\item Deux composés qui ont la même formule brute mais des formules développées différentes sont :
\begin{enumerate}
\item des conformères ;
\item des \cle{isomères} (d constitution) ;
\item des énantiomères.
\end{enumerate}
\item Un carbone asymétrique est un carbone :
\begin{enumerate}
\item porteur d'une double liaison ;
\item tétraédrique lié à quatre atomes ou groupes d'atomes \textbf{différents} ;
\item lié à quatre atomes d'hydrogène.
\end{enumerate}
\item Deux énantiomères possèdent :
\begin{enumerate}
\item des propriétés chimiques et physiques toujours différentes ;
\item les \textbf{mêmes propriétés physiques} (point de fusion, ébullition, densité\ldots) \textbf{mais des pouvoirs rotatoires opposés} ;
\item des masses molaires différentes.
\end{enumerate}
\item La représentation de Newman sert à décrire :
\begin{enumerate}
\item l'isomérie de fonction ;
\item les \cle{conformations} autour d'une liaison simple C--C ;
\item les configurations Z et E.
\end{enumerate}
\end{enumerate}
\end{exercice}

\begin{exercice}[Vrai ou faux ? Justifie]
Réponds par \textbf{vrai} ou \textbf{faux}, puis justifie ta réponse en une ou deux phrases.
\begin{enumerate}
\item Le propène \ce{CH2=CH-CH3} présente l'isomérie Z/E.
\item Une molécule possédant un plan de symétrie est \cle{achirale}.
\item Deux diastéréoisomères sont des images l'un de l'autre dans un miroir.
\item Un mélange racémique dévie le plan de la lumière polarisée.
\item La conformation décalée de l'éthane est plus stable que la conformation éclipsée.
\end{enumerate}
\end{exercice}

\begin{exercice}[Texte à trous]
Recopie et complète le texte suivant avec les mots de la liste : \emph{chirale ; énantiomères ; polarimètre ; racémique ; diastéréoisomères ; asymétrique ; mésomère}.

Une molécule qui n'est pas superposable à son image dans un miroir est dite \textbf{\ldots\ldots\ldots\ldots}. Elle possède généralement au moins un carbone \textbf{\ldots\ldots\ldots\ldots}. Deux stéréo-isomères images l'un de l'autre dans un miroir et non superposables sont appelés \textbf{\ldots\ldots\ldots\ldots} ; les autres stéréo-isomères, non images dans un miroir, sont des \textbf{\ldots\ldots\ldots\ldots}. L'appareil qui mesure la déviation du plan de la lumière polarisée est le \textbf{\ldots\ldots\ldots\ldots}. Un mélange équimolaire de deux énantiomères, optiquement inactif, est un mélange \textbf{\ldots\ldots\ldots\ldots}. Une molécule qui possède plusieurs carbones asymétriques mais qui reste achirale grâce à un plan de symétrie interne est une forme \textbf{\ldots\ldots\ldots\ldots}.
\end{exercice}

\begin{exercice}[Définitions express]
Définis en une phrase claire chacun des termes suivants, en donnant un exemple quand c'est possible :
\begin{enumerate}
\item \cle{isomérie de chaîne} ;
\item \cle{isomérie de fonction} ;
\item \cle{pouvoir rotatoire} (ou rotation spécifique) ;
\item \cle{conformation chaise} du cyclohexane.
\end{enumerate}
\end{exercice}

\courssub{Je m'entraîne}

\begin{exercice}[Isomères de constitution du pentane et du pentanol]
\begin{enumerate}
\item Écris les formules semi-développées de \textbf{tous} les isomères de constitution de formule brute \ce{C5H12} et nomme-les (nomenclature IUPAC). Précise le type d'isomérie qui les relie.
\item Écris les formules semi-développées des \textbf{quatre} alcools primaires de formule brute \ce{C5H12O} et nomme-les.
\item Le pentan-1-ol et le 1-méthoxybutane \ce{CH3-O-CH2-CH2-CH2-CH3} sont-ils isomères ? Si oui, de quel type d'isomérie s'agit-il ?
\end{enumerate}
\end{exercice}

\begin{exercice}[Carbone asymétrique et représentation de Cram]
On considère le butan-2-ol de formule semi-développée \ce{CH3-CH(OH)-CH2-CH3}.
\begin{enumerate}
\item Repère le carbone asymétrique (marque-le d'un astérisque) et justifie ta réponse en citant les quatre groupes qui y sont fixés.
\item Représente en perspective de Cram les deux énantiomères du butan-2-ol.
\item Cite deux propriétés physiques identiques et une propriété qui différencie ces deux énantiomères.
\end{enumerate}
\end{exercice}

\begin{exercice}[Isomérie Z/E]
On considère les alcènes suivants : but-2-ène (A), propène (B), pent-2-ène (C), 2-méthylbut-2-ène (D).
\begin{enumerate}
\item Écris la formule semi-développée de chaque alcène.
\item Identifie, en justifiant, ceux qui présentent l'isomérie Z/E.
\item Représente en perspective de Cram les configurations Z et E du pent-2-ène et nomme chaque stéréo-isomère selon la nomenclature \cle{Z/E}.
\end{enumerate}
\end{exercice}

\begin{exercice}[Représentations de Newman]
\begin{enumerate}
\item Construis les représentations de Newman des conformations \textbf{éclipsée} et \textbf{décalée} (staggered) de l'éthane, en regardant le long de la liaison C--C.
\item Construis la représentation de Newman de la conformation \textbf{anti} du butane (regard le long de la liaison C2--C3), puis celle de la conformation \textbf{éclipsée syn} (où les deux groupes méthyle se font face).
\item Classe ces quatre conformations par ordre d'\textbf{énergie potentielle croissante} (de la moins stable à la plus stable) en justifiant ton classement par les encombrements stériques (interactions H/H, H/Me, Me/Me).
\end{enumerate}
\end{exercice}

\courssub{Je me prépare au Bac}

\begin{exercice}[Type Bac — Les isomères de \ce{C4H10O}]
\emph{(D'après Baccalauréat TI, session 2019)}

Au cours d'une séance de travaux dirigés, un professeur de lycée technique de Douala demande à ses élèves d'étudier les composés organiques de formule brute \ce{C4H10O}.

Données : masses molaires atomiques en \si{g.mol^{-1}} : $M(\ce{C}) = 12$ ; $M(\ce{H}) = 1$ ; $M(\ce{O}) = 16$.
\begin{enumerate}
\item Calcule la masse molaire moléculaire de \ce{C4H10O}.
\item Écris les formules semi-développées de \textbf{tous} les isomères de constitution de formule brute \ce{C4H10O} et nomme-les (on attend \textbf{sept} composés).
\item Classe ces isomères en deux familles fonctionnelles et donne le nom de chaque famille.
\item Parmi ces isomères, identifie :
\begin{enumerate}
\item deux couples d'isomères de chaîne ;
\item deux couples d'isomères de position ;
\item un couple d'isomères de fonction.
\end{enumerate}
\item L'un de ces isomères possède un carbone asymétrique. Lequel ? Représente ses deux énantiomères en perspective de Cram.
\item On fait réagir le butan-1-ol avec l'acide éthanoïque en présence d'acide sulfurique concentré (catalyseur). Écris l'équation-bilan de la réaction d'estérification et nomme le produit organique obtenu.
\end{enumerate}
\end{exercice}

\begin{exercice}[Type Bac — Le 2-chlorobutane et l'activité optique]
\emph{(D'après Baccalauréat TI, session 2021)}

Le 2-chlorobutane, de formule semi-développée \ce{CH3-CHCl-CH2-CH3}, est utilisé comme intermédiaire de synthèse dans l'industrie pharmaceutique. Un technicien de la SONARA dispose d'un échantillon de ce composé et d'un polarimètre.

Données : masses molaires atomiques en \si{g.mol^{-1}} : $M(\ce{C}) = 12$ ; $M(\ce{H}) = 1$ ; $M(\ce{Cl}) = 35{,}5$.
\begin{enumerate}
\item Calcule la masse molaire moléculaire du 2-chlorobutane.
\item Montre que la molécule de 2-chlorobutane possède un carbone asymétrique que tu repéreras par un astérisque sur la formule semi-développée.
\item Représente en perspective de Cram les deux énantiomères du 2-chlorobutane.
\item Représente ces deux mêmes énantiomères en projection de Fischer, en plaçant la chaîne carbonée principale verticalement (groupe le plus oxydé en haut).
\item Le technicien mesure le pouvoir rotatoire d'une solution préparée avec l'échantillon et trouve une déviation nulle du plan de la lumière polarisée, alors que l'échantillon est pur (pas de solvant chiral). Propose une explication en nommant le type de mélange concerné.
\item Cite deux propriétés physiques communes aux deux énantiomères et une propriété biologique qui peut les différencier.
\end{enumerate}
\end{exercice}

\begin{exercice}[Type Bac — L'acide tartrique, du vin de palme à la stéréochimie]
\emph{(D'après Baccalauréat TI, session 2022)}

L'acide tartrique, de formule semi-développée \ce{HOOC-CH(OH)-CH(OH)-COOH}, est un acide présent dans le raisin et intervenant dans la fermentation des boissons traditionnelles (vin de palme, bili-bili). Louis Pasteur établit sa célébrité en séparant à la main, en 1848, les cristaux de ses deux formes optiquement actives.

Données : masses molaires atomiques en \si{g.mol^{-1}} : $M(\ce{C}) = 12$ ; $M(\ce{H}) = 1$ ; $M(\ce{O}) = 16$.
\begin{enumerate}
\item Calcule la masse molaire moléculaire de l'acide tartrique.
\item Repère, par des astérisques, les atomes de carbone asymétriques de la molécule et justifie ta réponse.
\item Combien de stéréo-isomères de configuration l'acide tartrique possède-t-il ? Justifie en expliquant pourquoi le nombre théorique $2^n$ n'est pas atteint.
\item Représente en projection de Fischer \textbf{tous} les stéréo-isomères de l'acide tartrique.
\item Parmi ces représentations, identifie :
\begin{enumerate}
\item les couples d'énantiomères ;
\item un couple de diastéréoisomères ;
\item la forme mésomère, en justifiant son achiralité par la présence d'un élément de symétrie.
\end{enumerate}
\item Une solution de l'un des énantiomères purs dévie le plan de la lumière polarisée vers la droite : on dit qu'elle est dextrogyre.
\begin{enumerate}
\item Nomme l'appareil qui permet cette mesure.
\item Que vaut le pouvoir rotatoire de la forme mésomère ? Justifie.
\item Que vaut le pouvoir rotatoire d'un mélange équimolaire des deux énantiomères ? Comment appelle-t-on un tel mélange ?
\end{enumerate}
\end{enumerate}
\end{exercice}

\newpage

\coursec{Corrigés des exercices}

\begin{corrige}[Exercice 1 — QCM : à toi de choisir]
\begin{enumerate}
\item \textbf{Réponse b.} Deux composés de même formule brute mais de formules développées différentes (enchaînement des atomes différent) sont des \cle{isomères de constitution}. Les conformères ne diffèrent que par une rotation autour d'une liaison simple (même enchaînement) ; les énantiomères ont la même constitution mais des configurations opposées.
\item \textbf{Réponse b.} Un \cle{carbone asymétrique} est un atome de carbone tétraédrique ($sp^3$) lié à \textbf{quatre atomes ou groupes d'atomes différents}. Un carbone porteur d'une double liaison est plan ($sp^2$) et ne peut pas être asymétrique.
\item \textbf{Réponse b.} Deux énantiomères ont les \textbf{mêmes propriétés physiques} (températures de fusion et d'ébullition, densité, solubilité, indice de réfraction) et la même masse molaire ; seuls leurs \textbf{pouvoirs rotatoires} sont opposés (même valeur absolue, signes contraires).
\item \textbf{Réponse b.} La \cle{représentation de Newman} décrit les \cle{conformations} d'une molécule observées en regardant le long d'une liaison simple C--C (formes éclipsée et décalée).
\end{enumerate}
\end{corrige}

\begin{corrige}[Exercice 2 — Vrai ou faux ? Justifie]
\begin{enumerate}
\item \textbf{Faux.} Le propène \ce{CH2=CH-CH3} ne présente pas l'isomérie Z/E car l'un des deux carbones de la double liaison (celui de \ce{CH2}) porte \textbf{deux atomes identiques} (deux H). La condition d'existence est que chaque carbone de \ce{C=C} porte deux substituants différents.
\item \textbf{Vrai.} Une molécule possédant un plan de symétrie (ou un centre de symétrie) est superposable à son image dans un miroir : elle est \cle{achirale} et optiquement inactive, même si elle contient des carbones asymétriques (cas des formes mésomères).
\item \textbf{Faux.} Deux stéréo-isomères images l'un de l'autre dans un miroir et non superposables sont des \textbf{énantiomères}. Les \cle{diastéréoisomères} sont des stéréo-isomères qui \textbf{ne sont pas} images l'un de l'autre dans un miroir.
\item \textbf{Faux.} Un mélange racémique (équimolaire des deux énantiomères) est \textbf{optiquement inactif} : les pouvoirs rotatoires, égaux en valeur absolue et de signes opposés, se compensent exactement ; la déviation mesurée est nulle.
\item \textbf{Vrai.} Dans la conformation \textbf{décalée}, les liaisons C--H sont à \ang{60} les unes des autres : les répulsions entre doublets liants sont minimales. Dans la conformation \textbf{éclipsée}, elles se font face (dièdre \ang{0}) : les répulsions sont maximales. La décalée est donc plus stable (écart d'environ \SI{12}{kJ.mol^{-1}} pour l'éthane).
\end{enumerate}
\end{corrige}

\begin{corrige}[Exercice 3 — Texte à trous]
Une molécule qui n'est pas superposable à son image dans un miroir est dite \textbf{chirale}. Elle possède généralement au moins un carbone \textbf{asymétrique}. Deux stéréo-isomères images l'un de l'autre dans un miroir et non superposables sont appelés \textbf{énantiomères} ; les autres stéréo-isomères, non images dans un miroir, sont des \textbf{diastéréoisomères}. L'appareil qui mesure la déviation du plan de la lumière polarisée est le \textbf{polarimètre}. Un mélange équimolaire de deux énantiomères, optiquement inactif, est un mélange \textbf{racémique}. Une molécule qui possède plusieurs carbones asymétriques mais qui reste achirale grâce à un plan de symétrie interne est une forme \textbf{mésomère}.
\end{corrige}

\begin{corrige}[Exercice 4 — Définitions express]
\begin{enumerate}
\item \textbf{Isomérie de chaîne} : isomérie de constitution entre composés de même fonction qui diffèrent par la \textbf{structure du squelette carboné} (chaîne linéaire ou ramifiée). Exemple : le butane \ce{CH3-CH2-CH2-CH3} et le 2-méthylpropane \ce{CH3-CH(CH3)-CH3}, tous deux de formule \ce{C4H10}.
\item \textbf{Isomérie de fonction} : isomérie de constitution entre composés qui possèdent des \textbf{groupes fonctionnels différents}. Exemple : l'éthanol \ce{CH3-CH2-OH} (alcool) et le méthoxyméthane \ce{CH3-O-CH3} (éther-oxyde), tous deux de formule \ce{C2H6O}.
\item \textbf{Pouvoir rotatoire} : propriété d'une substance chirale de faire \textbf{tourner le plan de polarisation de la lumière polarisée} d'un angle $\alpha$, mesuré au polarimètre ; la rotation spécifique $[\alpha]$ est cet angle rapporté à la longueur de la cuve et à la concentration. Il est positif (dextrogyre) ou négatif (lévogyre).
\item \textbf{Conformation chaise} du cyclohexane : conformation la plus stable du cyclohexane, dans laquelle le cycle est plié de sorte que tous les angles de valence valent environ \ang{109,5} et que toutes les liaisons C--H sont décalées deux à deux ; elle rappelle la forme d'un fauteuil (chaise), avec des hydrogènes \emph{axiaux} et \emph{équatoriaux}.
\end{enumerate}
\end{corrige}

\begin{corrige}[Exercice 5 — Isomères de constitution du pentane et du pentanol]
\textbf{1.} Les isomères de formule brute \ce{C5H12} sont au nombre de \textbf{trois} :
\begin{itemize}
\item \ce{CH3-CH2-CH2-CH2-CH3} : \textbf{pentane} (chaîne linéaire) ;
\item \ce{CH3-CH(CH3)-CH2-CH3} : \textbf{2-méthylbutane} ;
\item \ce{CH3-C(CH3)2-CH3} : \textbf{2,2-diméthylpropane} (néopentane).
\end{itemize}
Ces trois composés ont la même chaîne fonctionnelle (alcanes) mais des squelettes différents : ils sont liés par une \textbf{isomérie de chaîne}.

\textbf{2.} Les quatre alcools \textbf{primaires} (groupe \ce{-CH2-OH}) de formule \ce{C5H12O} :
\begin{itemize}
\item \ce{CH3-CH2-CH2-CH2-CH2-OH} : \textbf{pentan-1-ol} ;
\item \ce{CH3-CH(CH3)-CH2-CH2-OH} : \textbf{3-méthylbutan-1-ol} ;
\item \ce{CH3-CH2-CH(CH3)-CH2-OH} : \textbf{2-méthylbutan-1-ol} ;
\item \ce{CH3-C(CH3)2-CH2-OH} : \textbf{2,2-diméthylpropan-1-ol}.
\end{itemize}
On les obtient en fixant \ce{-CH2OH} sur chacun des trois squelettes de la question 1, en éliminant les doublons par symétrie.

\textbf{3.} Le pentan-1-ol \ce{C5H12O} et le 1-méthoxybutane \ce{CH3-O-CH2-CH2-CH2-CH3} ont la \textbf{même formule brute} \ce{C5H12O} (vérifions : \ce{CH3-O-C4H9} donne $1+4=5$ C, $3+9=12$ H, 1 O). Le premier est un \textbf{alcool}, le second un \textbf{éther-oxyde} : ce sont des \textbf{isomères de fonction}.
\end{corrige}

\begin{corrige}[Exercice 6 — Carbone asymétrique et représentation de Cram]
\textbf{1.} Formule semi-développée : \ce{CH3-C*H(OH)-CH2-CH3}. Le carbone n\textsuperscript{o}2 (marqué *) est \textbf{asymétrique} car il est tétraédrique et porte quatre groupes \textbf{tous différents} : \ce{-H}, \ce{-OH}, \ce{-CH3} (méthyle) et \ce{-CH2-CH3} (éthyle).

\textbf{2.} Représentations de Cram des deux énantiomères (images dans un miroir) :
\begin{popfigure}
\centering
\chemfig{C(-[2]H)(-[6]CH_3)(<[4]OH)(<:[0]CH_2CH_3)}
\hspace{2cm}
\chemfig{C(-[2]H)(-[6]CH_3)(<:[4]OH)(<[0]CH_2CH_3)}
\legende{Les deux énantiomères du butan-2-ol : à gauche, \ce{OH} en avant (trait plein) et \ce{CH2CH3} en arrière (trait hachuré) ; à droite, configuration inversée.}
\end{popfigure}
Dans la première représentation, le groupe \ce{-OH} est en avant du plan (trait plein épais) et le groupe éthyle en arrière (trait hachuré) ; dans la seconde, on a permuté les deux : les deux molécules sont images l'une de l'autre dans un miroir et \textbf{non superposables}.

\textbf{3.} Propriétés physiques \textbf{identiques} : température d'ébullition ($\approx \SI{99}{\celsius}$), densité (ou température de fusion, solubilité, indice de réfraction). Propriété qui les \textbf{différencie} : le \textbf{pouvoir rotatoire} — l'un est dextrogyre, l'autre lévogyre, avec des déviations de même valeur absolue et de signes opposés (ils peuvent aussi différer par leur activité biologique).
\end{corrige}

\begin{corrige}[Exercice 7 — Isomérie Z/E]
\textbf{1.} Formules semi-développées :
\begin{itemize}
\item (A) but-2-ène : \ce{CH3-CH=CH-CH3} ;
\item (B) propène : \ce{CH2=CH-CH3} ;
\item (C) pent-2-ène : \ce{CH3-CH=CH-CH2-CH3} ;
\item (D) 2-méthylbut-2-ène : \ce{CH3-C(CH3)=CH-CH3}.
\end{itemize}

\textbf{2.} Condition : chaque carbone de la double liaison doit porter \textbf{deux substituants différents}.
\begin{itemize}
\item \textbf{(A) but-2-ène : oui.} Chaque carbone porte H et \ce{CH3} (différents) : isomérie Z/E.
\item \textbf{(B) propène : non.} Le carbone de \ce{CH2} porte deux H identiques.
\item \textbf{(C) pent-2-ène : oui.} C2 porte H et \ce{CH3} ; C3 porte H et \ce{CH2CH3} : isomérie Z/E.
\item \textbf{(D) 2-méthylbut-2-ène : non.} Le carbone n\textsuperscript{o}2 porte \textbf{deux groupes méthyle identiques}.
\end{itemize}

\textbf{3.} Configurations du pent-2-ène \ce{CH3-CH=CH-CH2-CH3} :
\begin{popfigure}
\centering
\chemfig{CH_3-[:150]C(-[:210]H)=[:30]C(-[:330]H)-[:30]CH_2CH_3}
\hspace{2cm}
\chemfig{CH_3-[:150]C(-[:30]H)=[:210]C(-[:150]H)-[:210,,,,white]}
\legende{(Z)-pent-2-ène (groupes prioritaires \ce{CH3} et \ce{CH2CH3} du même côté) et (E)-pent-2-ène (groupes prioritaires opposés).}
\end{popfigure}
\begin{itemize}
\item \textbf{(Z)-pent-2-ène} : les deux groupes prioritaires (ici \ce{CH3} et \ce{CH2CH3}, prioritaires sur H) sont du \textbf{même côté} de la double liaison (\emph{zusammen} = ensemble) ;
\item \textbf{(E)-pent-2-ène} : les groupes prioritaires sont de part et d'autre de la double liaison (\emph{entgegen} = opposés).
\end{itemize}
Ces deux stéréo-isomères ne sont pas images dans un miroir : ce sont des \textbf{diastéréoisomères}, de propriétés physiques différentes.
\end{corrige}

\begin{corrige}[Exercice 8 — Représentations de Newman]
\textbf{1.} Éthane \ce{CH3-CH3}, regard le long de C--C (le point représente le carbone avant, le cercle le carbone arrière) :
\begin{popfigure}
\centering
\begin{tikzpicture}[scale=0.9]
% Éclipsée
\draw (0,0) circle (0.9);
\fill (0,0) circle (1.5pt);
\foreach \a in {90,210,330}{
  \draw (0,0) -- (\a:0.9) node[pos=1.35]{\scriptsize H};
  \draw (\a:0.9) -- (\a:1.25) node[pos=1.15]{\scriptsize H};
}
\node at (0,-1.7) {\scriptsize éclipsée};
% Décalée
\begin{scope}[xshift=4cm]
\draw (0,0) circle (0.9);
\fill (0,0) circle (1.5pt);
\foreach \a in {90,210,330}{
  \draw (0,0) -- (\a:0.9) node[pos=1.35]{\scriptsize H};
}
\foreach \a in {30,150,270}{
  \draw (\a:0.9) -- (\a:1.25) node[pos=1.15]{\scriptsize H};
}
\node at (0,-1.7) {\scriptsize décalée};
\end{scope}
\end{tikzpicture}
\legende{Projections de Newman de l'éthane : forme éclipsée (gauche) et forme décalée (droite).}
\end{popfigure}
Dans la forme éclipsée, les H de l'arrière sont cachés derrière ceux de l'avant (on les dessine légèrement décalés pour les rendre visibles) ; dans la forme décalée, ils occupent les positions intermédiaires (\ang{60}).

\textbf{2.} Butane, regard le long de C2--C3 (les deux groupes \ce{CH3} remplacent un H sur chaque carbone) :
\begin{popfigure}
\centering
\begin{tikzpicture}[scale=0.9]
% anti
\draw (0,0) circle (0.9);
\fill (0,0) circle (1.5pt);
\draw (0,0) -- (90:0.9) node[pos=1.4]{\scriptsize \ce{CH3}};
\draw (0,0) -- (210:0.9) node[pos=1.35]{\scriptsize H};
\draw (0,0) -- (330:0.9) node[pos=1.35]{\scriptsize H};
\draw (270:0.9) -- (270:1.3) node[pos=1.2]{\scriptsize \ce{CH3}};
\draw (30:0.9) -- (30:1.25) node[pos=1.15]{\scriptsize H};
\draw (150:0.9) -- (150:1.25) node[pos=1.15]{\scriptsize H};
\node at (0,-1.8) {\scriptsize anti (décalée)};
% éclipsée syn
\begin{scope}[xshift=5cm]
\draw (0,0) circle (0.9);
\fill (0,0) circle (1.5pt);
\draw (0,0) -- (90:0.9) node[pos=1.4]{\scriptsize \ce{CH3}};
\draw (0,0) -- (210:0.9) node[pos=1.35]{\scriptsize H};
\draw (0,0) -- (330:0.9) node[pos=1.35]{\scriptsize H};
\draw (90:0.9) -- (90:1.3) node[pos=1.2]{\scriptsize \ce{CH3}};
\draw (210:0.9) -- (210:1.25) node[pos=1.15]{\scriptsize H};
\draw (330:0.9) -- (330:1.25) node[pos=1.15]{\scriptsize H};
\node at (0,-1.8) {\scriptsize éclipsée syn};
\end{scope}
\end{tikzpicture}
\legende{Projections de Newman du butane (regard C2--C3) : conformation anti (gauche) et conformation éclipsée syn (droite).}
\end{popfigure}

\textbf{3.} Classement par énergie potentielle \textbf{croissante} (de la plus stable à la moins stable) :
\[
\text{décalée (éthane)} < \text{anti (butane)} < \text{éclipsée (éthane)} < \text{éclipsée syn (butane)}
\]
Justification par les encombrements stériques :
\begin{itemize}
\item \textbf{Décalée (éthane)} : aucune éclipse, interactions H/H minimales $\Rightarrow$ la plus stable ;
\item \textbf{Anti (butane)} : décalée, mais les groupes \ce{CH3}, plus volumineux que H, créent des interactions H/Me résiduelles (gauche) $\Rightarrow$ un peu moins stable que l'éthane décalé ;
\item \textbf{Éclipsée (éthane)} : trois éclipses H/H $\Rightarrow$ énergie élevée ;
\item \textbf{Éclipsée syn (butane)} : une éclipse Me/Me (très encombrante) et deux éclipses H/H $\Rightarrow$ la moins stable (énergie maximale).
\end{itemize}
\end{corrige}

\begin{corrige}[Exercice 9 — Type Bac : les isomères de \ce{C4H10O}]
\textbf{1.} Masse molaire :
\[
M(\ce{C4H10O}) = 4\times 12 + 10\times 1 + 16 = 48 + 10 + 16 = \SI{74}{g.mol^{-1}}.
\]

\textbf{2.} Les \textbf{sept} isomères de constitution :
\begin{itemize}
\item \ce{CH3-CH2-CH2-CH2-OH} : butan-1-ol ;
\item \ce{CH3-CH2-CH(OH)-CH3} : butan-2-ol ;
\item \ce{CH3-CH(CH3)-CH2-OH} : 2-méthylpropan-1-ol ;
\item \ce{CH3-C(CH3)(OH)-CH3} : 2-méthylpropan-2-ol ;
\item \ce{CH3-CH2-CH2-O-CH3} : 1-méthoxypropane ;
\item \ce{CH3-CH(CH3)-O-CH3} : 2-méthoxypropane ;
\item \ce{CH3-CH2-O-CH2-CH3} : éthoxyéthane.
\end{itemize}

\textbf{3.} Deux familles fonctionnelles : les \textbf{alcools} (groupe hydroxyle \ce{-OH}, quatre premiers) et les \textbf{éthers-oxydes} (groupe \ce{-O-} entre deux chaînes carbonées, trois derniers).

\textbf{4.} Exemples de couples :
\begin{enumerate}
\item Isomérie de \textbf{chaîne} : butan-1-ol / 2-méthylpropan-1-ol ; butan-2-ol / 2-méthylpropan-2-ol ;
\item Isomérie de \textbf{position} : butan-1-ol / butan-2-ol ; 1-méthoxypropane / 2-méthoxypropane ;
\item Isomérie de \textbf{fonction} : butan-1-ol / éthoxyéthane (alcool / éther-oxyde).
\end{enumerate}

\textbf{5.} Le \textbf{butan-2-ol} \ce{CH3-C*H(OH)-CH2-CH3} possède un carbone asymétrique (C2 porte H, \ce{OH}, \ce{CH3}, \ce{C2H5}). Ses deux énantiomères :
\begin{popfigure}
\centering
\chemfig{C(-[2]H)(-[6]CH_3)(<[4]OH)(<:[0]CH_2CH_3)}
\hspace{2cm}
\chemfig{C(-[2]H)(-[6]CH_3)(<:[4]OH)(<[0]CH_2CH_3)}
\legende{Énantiomères du butan-2-ol en représentation de Cram.}
\end{popfigure}

\textbf{6.} Estérification du butan-1-ol par l'acide éthanoïque :
\[ \ce{CH3COOH + CH3CH2CH2CH2OH <=> CH3COOCH2CH2CH2CH3 + H2O} \]
Le produit organique est l'\textbf{éthanoate de butyle} (ester à odeur fruitée). Réaction lente, limitée (équilibre), catalysée par \ce{H2SO4}.

\textbf{Barème indicatif (sur 10)} : Q1 : 1 pt ; Q2 : 3,5 pts (0,5 pt par isomère nommé) ; Q3 : 1 pt ; Q4 : 1,5 pt ; Q5 : 1,5 pt ; Q6 : 1,5 pt.

\textbf{Points de vigilance} : ne pas oublier les éthers-oxydes (erreur classique : n'écrire que les quatre alcools) ; vérifier que le 2-méthylpropan-2-ol n'a \emph{pas} de carbone asymétrique (deux \ce{CH3} identiques) ; l'équation d'estérification s'écrit avec une double flèche (réaction limitée).
\end{corrige}

\begin{corrige}[Exercice 10 — Type Bac : le 2-chlorobutane et l'activité optique]
\textbf{1.} Formule brute : \ce{C4H9Cl}.
\[
M = 4\times 12 + 9\times 1 + 35{,}5 = 48 + 9 + 35{,}5 = \SI{92,5}{g.mol^{-1}}.
\]

\textbf{2.} \ce{CH3-C*HCl-CH2-CH3} : le carbone n\textsuperscript{o}2 est tétraédrique et porte quatre groupes différents : \ce{-H}, \ce{-Cl}, \ce{-CH3} et \ce{-CH2-CH3}. C'est un \textbf{carbone asymétrique} : la molécule est chirale.

\textbf{3.} Énantiomères en perspective de Cram :
\begin{popfigure}
\centering
\chemfig{C(-[2]H)(-[6]CH_3)(<[4]Cl)(<:[0]CH_2CH_3)}
\hspace{2cm}
\chemfig{C(-[2]H)(-[6]CH_3)(<:[4]Cl)(<[0]CH_2CH_3)}
\legende{Énantiomères du 2-chlorobutane en représentation de Cram.}
\end{popfigure}

\textbf{4.} Projections de Fischer (chaîne verticale, \ce{CH3} en haut car le carbone le plus proche du substituant prioritaire ; traits horizontaux vers l'avant) :
\begin{popfigure}
\centering
\begin{tikzpicture}[scale=0.8]
\draw (0,0) -- (0,2.4); \draw (-0.9,1.2) -- (0.9,1.2);
\node at (0,2.7) {\ce{CH3}}; \node at (0,-0.3) {\ce{C2H5}};
\node at (-1.2,1.2) {\ce{Cl}}; \node at (1.2,1.2) {\ce{H}};
\node at (0,-0.9) {\scriptsize énantiomère 1};
\begin{scope}[xshift=4cm]
\draw (0,0) -- (0,2.4); \draw (-0.9,1.2) -- (0.9,1.2);
\node at (0,2.7) {\ce{CH3}}; \node at (0,-0.3) {\ce{C2H5}};
\node at (-1.2,1.2) {\ce{H}}; \node at (1.2,1.2) {\ce{Cl}};
\node at (0,-0.9) {\scriptsize énantiomère 2};
\end{scope}
\end{tikzpicture}
\legende{Projections de Fischer des deux énantiomères du 2-chlorobutane. Les liaisons horizontales sortent du plan.}
\end{popfigure}
Les deux projections ne diffèrent que par l'échange de \ce{Cl} et \ce{H} sur la ligne horizontale : elles sont images dans un miroir et non superposables.

\textbf{5.} Une déviation \textbf{nulle} pour un échantillon pur s'explique par la présence d'un \textbf{mélange racémique} : mélange équimolaire des deux énantiomères, dont les pouvoirs rotatoires, égaux en valeur absolue et de signes opposés, se compensent exactement. Le mélange est optiquement inactif.

\textbf{6.} Propriétés physiques \textbf{communes} : température d'ébullition, densité (ou indice de réfraction, solubilité). Propriété biologique pouvant les \textbf{différencier} : l'activité pharmacologique — un seul énantiomère est souvent actif sur un récepteur biologique (cible elle-même chirale), l'autre pouvant être inactif, voire nocif.

\textbf{Barème indicatif (sur 10)} : Q1 : 1 pt ; Q2 : 1,5 pt ; Q3 : 2 pts ; Q4 : 2 pts ; Q5 : 1,5 pt ; Q6 : 2 pts.

\textbf{Points de vigilance} : en Fischer, les traits horizontaux pointent \emph{vers l'avant} et les verticaux vers l'arrière ; on ne fait \emph{tourner} une projection de Fischer que de \ang{180} dans le plan (jamais \ang{90}) ; une déviation nulle ne signifie pas « molécule achirale » : penser au racémique.
\end{corrige}

\begin{corrige}[Exercice 11 — Type Bac : l'acide tartrique]
\textbf{1.} Formule brute : \ce{C4H6O6}.
\[
M = 4\times 12 + 6\times 1 + 6\times 16 = 48 + 6 + 96 = \SI{150}{g.mol^{-1}}.
\]

\textbf{2.} \ce{HOOC-C*H(OH)-C*H(OH)-COOH}. Les carbones n\textsuperscript{o}2 et n\textsuperscript{o}3 sont chacun tétraédriques et liés à quatre groupes différents : \ce{-H}, \ce{-OH}, \ce{-COOH} et \ce{-CH(OH)-COOH}. Ce sont deux \textbf{carbones asymétriques}. (Les carbones des groupes \ce{-COOH} sont $sp^2$, plans : non asymétriques.)

\textbf{3.} Avec $n = 2$ carbones asymétriques, le nombre théorique de stéréo-isomères serait $2^n = 4$. Mais la molécule possède une \textbf{symétrie interne} (les deux moitiés \ce{HOOC-CH(OH)-} sont identiques) : la configuration dans laquelle les deux carbones ont des configurations opposées possède un \textbf{plan de symétrie} ; elle est superposable à son image (forme \textbf{mésomère}, achirale). Les deux configurations « symétriques » attendues se confondent donc en une seule. Bilan : \textbf{3 stéréo-isomères} : un couple d'énantiomères (+) et (--) et une forme mésomère.

\textbf{4.} Projections de Fischer (chaîne verticale, \ce{COOH} en haut et en bas) :
\begin{popfigure}
\centering
\begin{tikzpicture}[scale=0.75]
% (+)
\draw (0,0.6) -- (0,3.4); \draw (-0.8,1.4) -- (0.8,1.4); \draw (-0.8,2.6) -- (0.8,2.6);
\node at (0,3.7) {\scriptsize \ce{COOH}}; \node at (0,0.3) {\scriptsize \ce{COOH}};
\node at (-1.1,2.6) {\scriptsize \ce{OH}}; \node at (1.1,2.6) {\scriptsize \ce{H}};
\node at (-1.1,1.4) {\scriptsize \ce{OH}}; \node at (1.1,1.4) {\scriptsize \ce{H}};
\node at (0,-0.4) {\scriptsize (1) : (+)};
% (--)
\begin{scope}[xshift=3.6cm]
\draw (0,0.6) -- (0,3.4); \draw (-0.8,1.4) -- (0.8,1.4); \draw (-0.8,2.6) -- (0.8,2.6);
\node at (0,3.7) {\scriptsize \ce{COOH}}; \node at (0,0.3) {\scriptsize \ce{COOH}};
\node at (-1.1,2.6) {\scriptsize \ce{H}}; \node at (1.1,2.6) {\scriptsize \ce{OH}};
\node at (-1.1,1.4) {\scriptsize \ce{H}}; \node at (1.1,1.4) {\scriptsize \ce{OH}};
\node at (0,-0.4) {\scriptsize (2) : (--)};
\end{scope}
% méso
\begin{scope}[xshift=7.2cm]
\draw (0,0.6) -- (0,3.4); \draw (-0.8,1.4) -- (0.8,1.4); \draw (-0.8,2.6) -- (0.8,2.6);
\node at (0,3.7) {\scriptsize \ce{COOH}}; \node at (0,0.3) {\scriptsize \ce{COOH}};
\node at (-1.1,2.6) {\scriptsize \ce{OH}}; \node at (1.1,2.6) {\scriptsize \ce{H}};
\node at (-1.1,1.4) {\scriptsize \ce{H}}; \node at (1.1,1.4) {\scriptsize \ce{OH}};
\draw[dashed,popink] (-1.5,2) -- (1.5,2);
\node at (0,-0.4) {\scriptsize (3) : mésomère};
\end{scope}
\end{tikzpicture}
\legende{Projections de Fischer des trois stéréo-isomères de l'acide tartrique : (1) acide L-(+)-tartrique, (2) acide D-(--)-tartrique, (3) acide mésotartrique (plan de symétrie horizontal en pointillés).}
\end{popfigure}
(La quatrième combinaison, \ce{OH} à gauche en haut et à droite en bas, est identique à (3) par rotation de \ang{180} dans le plan.)

\textbf{5.}
\begin{enumerate}
\item \textbf{Couple d'énantiomères} : (1) et (2), images dans un miroir, non superposables ;
\item \textbf{Couple de diastéréoisomères} : (1) et (3) [ou (2) et (3)] : stéréo-isomères non images dans un miroir ;
\item \textbf{Forme mésomère} : (3). Elle possède un \textbf{plan de symétrie horizontal} (en pointillés) qui échange les deux moitiés de la molécule ; elle est donc superposable à son image dans un miroir : \textbf{achirale} malgré ses deux carbones asymétriques.
\end{enumerate}

\textbf{6.}
\begin{enumerate}
\item L'appareil de mesure est le \textbf{polarimètre}.
\item Le pouvoir rotatoire de la forme mésomère est \textbf{nul} : molécule achirale (compensation \emph{interne} des effets des deux carbones asymétriques de configurations opposées).
\item Un mélange équimolaire de (1) et (2) a un pouvoir rotatoire \textbf{nul} (compensation externe des déviations opposées) : c'est un \textbf{mélange racémique}, noté $(\pm)$.
\end{enumerate}

\textbf{Barème indicatif (sur 12)} : Q1 : 1 pt ; Q2 : 2 pts ; Q3 : 2 pts ; Q4 : 3 pts ; Q5 : 2 pts ; Q6 : 2 pts.

\textbf{Points de vigilance} : ne pas appliquer aveuglément la règle $2^n$ sans vérifier l'existence d'une symétrie interne ; en Fischer, une rotation de \ang{180} dans le plan laisse la configuration inchangée, une rotation de \ang{90} la change ; bien distinguer compensation \emph{interne} (mésomère) et compensation \emph{externe} (racémique) — dans les deux cas la déviation mesurée est nulle, mais pour des raisons différentes.
\end{corrige}

\newpage

\coursec{Fiche bilan}

\courssub{Carte mentale de la leçon}

\begin{popfigure}
\begin{tikzpicture}[
  central/.style={circle, draw=popink, fill=popink!15, text=popdark, font=\bfseries\small, align=center, minimum size=2.8cm},
  branche/.style={rounded corners=4pt, draw=#1, fill=#1!12, text=popdark, font=\small, align=center, inner sep=4pt},
  lien/.style={thick, draw=#1},
  node distance=0.5cm]
\node[central] (C) {Stéréochimie\\ Isomérie};
\node[branche=popblue, above left=0.6cm and 1.8cm of C] (B1) {\textbf{Isomérie de constitution}\\ chaîne -- position -- fonction\\ (même formule brute)};
\node[branche=poporange, above=0.7cm of C] (B2) {\textbf{Chiralité}\\ carbone asymétrique C*\\ 4 substituants différents};
\node[branche=poppurple, above right=0.6cm and 1.8cm of C] (B3) {\textbf{Énantiomères}\\ images miroir non\\ superposables (Fischer)};
\node[branche=popgreen, right=2.0cm of C] (B4) {\textbf{Activité optique}\\ pouvoir rotatoire $[\alpha]$\\ mélange racémique};
\node[branche=poppink, below right=0.6cm and 1.8cm of C] (B5) {\textbf{Diastéréoisomères}\\ stéréoisomères non\\ images miroir};
\node[branche=popgold, below=0.7cm of C] (B6) {\textbf{Isomérie Z/E}\\ alcènes : Z même côté\\ E côtés opposés (CIP)};
\node[branche=popmuted, below left=0.6cm and 1.8cm of C] (B7) {\textbf{Conformations}\\ Newman : éclipsée/décalée\\ cyclohexane : chaise/bateau};
\draw[lien=popblue] (C) -- (B1);
\draw[lien=poporange] (C) -- (B2);
\draw[lien=poppurple] (C) -- (B3);
\draw[lien=popgreen] (C) -- (B4);
\draw[lien=poppink] (C) -- (B5);
\draw[lien=popgold] (C) -- (B6);
\draw[lien=popmuted] (C) -- (B7);
\end{tikzpicture}
\legende{Carte mentale : les sept idées-forces de la stéréochimie.}
\end{popfigure}

\begin{savaistu}[Stéréochimie et monde vivant : le cas du limonène]
Le limonène (\ce{C10H16}), terpène majeur des écorces d'agrumes, possède un carbone asymétrique. 
Son énantiomère \textbf{(+)-limonène} (dextrogyre) sent l'orange (présent dans l'huile essentielle d'orange du Mont Cameroun), 
tandis que son image miroir \textbf{(-)-limonène} (lévogyre) sent le pin et le turpentine. 
Les récepteurs olfactifs de notre nez sont chiraux : ils distinguent les deux formes. 
C'est la raison pour laquelle la stéréochimie est cruciale en parfumerie, en agroalimentaire (bili-bili, vin de palme) et surtout en pharmacie (ex. thalidomide : un énantiomère sédatif, l'autre tératogène).
\end{savaistu}

\courssub{L'essentiel à retenir}

\begin{aretenir}[Fiche bilan — Stéréochimie]
\textbf{Définitions clés}
\begin{itemize}
  \item \cle{Isomères} : composés de \textbf{même formule brute} mais de propriétés physiques et/ou chimiques différentes.
  \item \cle{Isomérie de constitution} : les atomes ne sont pas reliés dans le même ordre (connectivité différente). Trois types :
        \begin{itemize}
          \item \textbf{Chaîne} : squelette carboné différent (ex. \ce{C4H10} : butane / 2-méthylpropane).
          \item \textbf{Position} : même squelette, même fonction, position différente (ex. \ce{C3H7OH} : propan-1-ol / propan-2-ol).
          \item \textbf{Fonction} : groupes caractéristiques différents (ex. \ce{C2H6O} : éthanol \ce{CH3CH2OH} / diméthyléther \ce{CH3OCH3}).
        \end{itemize}
  \item \cle{Carbone asymétrique} (C*) : carbone \textbf{sp$^3$} (tétraédrique, 4 liaisons simples) lié à \textbf{quatre atomes ou groupes différents}. Noté C* ou \ce{C^@}.
  \item \cle{Chiralité} : une molécule est \textbf{chirale} si elle n'est \textbf{pas superposable à son image dans un miroir} (absence de plan, centre ou axe de symétrie impropre). La présence d'un C* est une condition suffisante fréquente (mais non nécessaire).
  \item \cle{Énantiomères} : couple de stéréoisomères images l'un de l'autre dans un miroir, \textbf{non superposables}. Mêmes propriétés physiques (pf, pe, densité, solubilité) et chimiques en milieu achiral. \textbf{Pouvoirs rotatoires opposés} ($[\alpha]_D$ et $-[\alpha]_D$).
  \item \cle{Diastéréoisomères} : stéréoisomères qui ne sont \textbf{pas} images miroir (ex. isomères Z/E ; molécules à plusieurs C* dont \emph{au moins un} mais \emph{pas tous} ont une configuration inversée). Propriétés physiques \textbf{différentes} (pf, pe, solubilité, $R_f$ CCM).
  \item \cle{Activité optique} : une substance \textbf{optiquement active} fait tourner le plan de polarisation de la lumière polarisée d'un angle $\alpha$ (degré). \textbf{Dextrogyre} ($+$, sens horaire) ou \textbf{lévogyre} ($-$, sens anti-horaire).
  \item \cle{Mélange racémique} : mélange équimolaire (50 \% / 50 \%) de deux énantiomères. Noté $(\pm)$ ou \textit{rac}. \textbf{Inactif} optiquement ($\alpha_{\text{obs}} = 0$).
  \item \cle{Isomérie Z/E (cis/trans)} : pour un alcène \ce{R^1R^2C=CR^3R^4} avec $\ce{R^1}\neq\ce{R^2}$ et $\ce{R^3}\neq\ce{R^4}$. Priorité CIP (numéro atomique) : \textbf{Z} (\emph{zusammen}, groupes prioritaires du \emph{même côté}), \textbf{E} (\emph{entgegen}, côtés \emph{opposés}).
  \item \cle{Conformations} : arrangements spatiaux obtenus par rotation autour d'une liaison simple $\sigma$ (libre à T° ambiante). Éthane : forme \textbf{décalée} (staggered, plus stable, énergie min) et \textbf{éclipsée} (eclipsed, instable, énergie max). Cyclohexane : conformation \textbf{chaise} (stable, sans tension) et \textbf{bateau} (instable, interactions de drapeau).
\end{itemize}

\textbf{Formules et grandeurs à connaître par cœur}
\begin{center}
\begin{tabularx}{\linewidth}{|l|X|}
\hline
\rowcolor{popblueL} \textbf{Grandeur / notion} & \textbf{Expression / règle} \\
\hline
Nombre max de stéréoisomères & $2^n$ pour $n$ carbones asymétriques (si pas de plan de symétrie / méso) \\
\hline
Pouvoir rotatoire spécifique & $[\alpha]_{\lambda}^{T} = \dfrac{\alpha}{\ell \times c}$ \quad ($\ell$ en \si{dm}, $c$ en \si{g.mL^{-1}}, $\lambda$ = raie D Na \SI{589}{nm}, $T$ = \SI{20}{\celsius}) \\
\hline
Mélange racémique & 50 \% ($+$) + 50 \% ($-$) $\Rightarrow$ $\alpha_{\text{obs}} = 0$ \quad ($[\alpha]_{\text{rac}} = 0$) \\
\hline
Condition isomérie Z/E & Alcène \ce{C=C} tétra-substitué ou tri-substitué : chaque C porte 2 groupes \textbf{différents} \\
\hline
Condition de chiralité & Présence d'un C* \textbf{ET} absence de tout élément de symétrie (plan, centre, axe $S_n$) \\
\hline
\end{tabularx}
\end{center}

\textbf{Représentations à maîtriser}
\begin{itemize}
  \item \cle{Perspective (Cram / Cram-Nagata)} : Liaison pleine dans le plan du papier, \textbf{trait plein épais (wedge)} vers l'avant (observateur), \textbf{trait hachuré (dash)} vers l'arrière. Exemple acide (S)-lactique : \chemfig{C(-[2]H)(<[5]CH_3)(<:[7]OH)-C(=[1]O)-[7]OH}.
  \item \cle{Projection de Fischer} : Chaîne carbonée la plus longue \textbf{verticale} (liaisons vers l'arrière, \emph{en dessous} du plan), carbone le plus oxydé en haut. Substituants \textbf{horizontaux} (liaisons vers l'avant, \emph{sortant} du plan). Le C* est l'intersection de la croix. \textbf{Permuter deux groupes} $\Leftrightarrow$ inversion de configuration $\Leftrightarrow$ passage à l'énantiomère. Rotation de 180° dans le plan = identité ; rotation de 90° = inversion.
  \item \cle{Projection de Newman} : Regard le long d'une liaison \ce{C-C}. Carbone \textbf{avant} = point central (3 liaisons à 120°). Carbone \textbf{arrière} = cercle (3 liaisons à 120°). Forme \textbf{décalée (staggered)} = plus stable (répulsions minimisées) ; forme \textbf{éclipsée} = maximum d'énergie (répulsions maximales).
\end{itemize}

\textbf{Méthodes express}
\begin{itemize}
  \item \textbf{Repérer un C*} : Chercher un carbone \textbf{tétraédrique} (4 liaisons simples) portant 4 groupes \textbf{tous différents} (atomes de numéros atomiques différents ou chaînes différentes). Le marquer d'un astérisque \textbf{*}.
  \item \textbf{Distinguer énantiomères / diastéréoisomères} : Construire l'image miroir (ou inverser tous les C* en Fischer). Superposable ? $\rightarrow$ Même molécule (achiral ou méso). Non superposable ? $\rightarrow$ \textbf{Énantiomères} (si \emph{tous} les C* inversés) ou \textbf{Diastéréoisomères} (si \emph{au moins un} C* identique et \emph{au moins un} inversé).
  \item \textbf{Attribuer Z ou E (CIP)} : Sur chaque carbone de la double liaison, classer les deux substituants par numéro atomique croissant (priorité = Z le plus grand). Regarder les deux groupes prioritaires : même côté $\rightarrow$ \textbf{Z} ; côtés opposés $\rightarrow$ \textbf{E}.
  \item \textbf{Différencier deux isomères de fonction} : Identifier les \textbf{groupes caractéristiques} (fonctions) : \ce{-OH} (alcool), \ce{-CHO} (aldéhyde), \ce{-COOH} (acide), \ce{-CO-} (cétone), \ce{-O-} (éther), \ce{-NH2} (amine), etc. Ex. \ce{C3H6O} : propanal \ce{CH3CH2CHO} vs propanone \ce{CH3COCH3} vs prop-2-én-1-ol \ce{CH2=CH-CH2OH}.
\end{itemize}
\end{aretenir}

\begin{aretenir}[Checklist avant le Bac]
\begin{itemize}
  \item[$\square$] Je sais définir et illustrer les trois types d'isomérie de constitution (chaîne, position, fonction) avec des exemples \ce{C_nH_{2n+2}} ou \ce{C_nH_{2n}O}.
  \item[$\square$] Je sais repérer \textbf{tous} les carbones asymétriques d'une molécule (y compris dans les cycles) et les marquer d'un astérisque.
  \item[$\square$] Je sais calculer le nombre maximal de stéréoisomères : $2^n$ (et identifier les cas méso où le nombre est réduit).
  \item[$\square$] Je sais dessiner une molécule chirale en perspective (Cram) avec les trois types de liaisons (plein, wedge, dash) sans erreur de géométrie tétraédrique.
  \item[$\square$] Je sais représenter deux énantiomères en projection de Fischer (chaîne verticale, groupes horizontaux) et passer de l'un à l'autre par permutation de deux groupes.
  \item[$\square$] Je sais distinguer sans hésiter énantiomères (images miroir, tous C* inversés) et diastéréoisomères (autres cas : Z/E, ou C* partiellement inversés).
  \item[$\square$] Je connais la définition du pouvoir rotatoire spécifique $[\alpha]$ et ses unités ; je sais qu'un mélange racémique est optiquement inactif.
  \item[$\square$] Je sais vérifier la condition d'existence de l'isomérie Z/E (chaque C de la double liaison porte 2 groupes différents) et attribuer le descripteur Z ou E selon les règles CIP.
  \item[$\square$] Je sais dessiner les projections de Newman éclipsée et décalée de l'éthane (ou butane) et justifier la stabilité relative (interactions de gauche, répulsions éclipser).
  \item[$\square$] Je sais reconnaître et nommer les conformations chaise (stable) et bateau (instable) du cyclohexane ; identifier les positions axiales/équatoriales.
  \item[$\square$] Je sais expliquer pourquoi la stéréochimie est cruciale en pharmacie / biologie (ex. limonène, thalidomide, adrénaline, sucres, acides aminés) : récepteurs chiraux $\Rightarrow$ activités différentes.
  \item[$\square$] Je vérifie systématiquement mes formules développées : \textbf{4 liaisons pour chaque carbone}, respect de la tétravalence.
\end{itemize}
\end{aretenir}

\findecours

\end{document}
